\documentclass{article}
\usepackage{booktabs}% http://ctan.org/pkg/booktabs
\usepackage{amsmath}
\usepackage[margin=0.9in]{geometry}   % slightly tighter 
\usepackage{array}

\usepackage{graphicx}
\usepackage{caption}
\usepackage{subcaption}
\usepackage{needspace} % optional but recommended
% \usepackage{pdflscape} % if you keep landscape pages

\usepackage{float}
\usepackage{threeparttable}
\usepackage{adjustbox}
\usepackage{pdflscape}  % 
\usepackage{comment}
\usepackage[hidelinks]{hyperref}
\title{Senegal Tax Audits: Comments and Revisions}
%\author{Roldan Enamorado}
\date{February, 2026}
\setcounter{tocdepth}{4}
\newcommand{\OutPath}{Output/}
% Helper: one row (label + Desk + Full)
\newcommand{\BinRow}[3]{%
\textbf{#1} &
\includegraphics[width=0.44\textwidth]{\OutPath#2} &
\includegraphics[width=0.44\textwidth]{\OutPath#3} \\
}

% One nested panel: (a) title + two subplots labeled a.1 and a.2
% Args:
%  #1 panel letter (a, b, c...)
%  #2 panel title (e.g., Deciles / Width 1)
%  #3 desk filename (pdf)
%  #4 full filename (pdf)
\newcommand{\BinNestedPanel}[4]{%
\Needspace{0.33\textheight} % if panel doesn't fit, start it on next page
\par\noindent
\begin{center}
\textbf{Panel #1: #2}
\end{center}
\vspace{-0.4em}
\begin{minipage}{0.49\textwidth}
\centering
\textbf{#1.1} Desk audits\par
\vspace{0.10em}
\includegraphics[width=\linewidth]{\OutPath#3}
\end{minipage}\hfill
\begin{minipage}{0.49\textwidth}
\centering
\textbf{#1.2} Full audits\par
\vspace{0.10em}
\includegraphics[width=\linewidth]{\OutPath#4}
\end{minipage}
\par\vspace{0.6em}
}
% Figure block for pooled-sample binning outputs.
% Args:
%  #1 = base filename prefix (before _Desk/_Full and _decile/_1u/...)
%       e.g. avg_ninspectors_binplot_yhatrf
\newcommand{\BinCompareBinsPanelsWholeSample}[1]{%
\BinNestedPanel{a}{Deciles}{#1\_Desk\_decile\_whole_sample.pdf}{#1\_Full\_decile\_whole_sample.pdf}
\BinNestedPanel{b}{Quintiles}{#1\_Desk\_quintile\_whole_sample.pdf}{#1\_Full\_quintile\_whole_sample.pdf}
\BinNestedPanel{c}{Width 1}{#1\_Desk\_1u\_whole_sample.pdf}{#1\_Full\_1u\_whole_sample.pdf}
\BinNestedPanel{d}{Width 2}{#1\_Desk\_2u\_whole_sample.pdf}{#1\_Full\_2u\_whole_sample.pdf}
\BinNestedPanel{e}{Width 4}{#1\_Desk\_4u\_whole_sample.pdf}{#1\_Full\_4u\_whole_sample.pdf}
\BinNestedPanel{f}{Width 5}{#1\_Desk\_5u\_whole_sample.pdf}{#1\_Full\_5u\_whole_sample.pdf}
}

% Figure block for within-method binning outputs.
% Args:
%  #1 = base filename prefix (before _Desk/_Full and _decile/_1u/...)
\newcommand{\BinCompareBinsPanelsWithinMethod}[1]{%
\BinNestedPanel{a}{Deciles}{#1\_Desk\_decile.pdf}{#1\_Full\_decile.pdf}
\BinNestedPanel{b}{Quintiles}{#1\_Desk\_quintile.pdf}{#1\_Full\_quintile.pdf}
\BinNestedPanel{c}{Width 1}{#1\_Desk\_1u.pdf}{#1\_Full\_1u.pdf}
\BinNestedPanel{d}{Width 2}{#1\_Desk\_2u.pdf}{#1\_Full\_2u.pdf}
\BinNestedPanel{e}{Width 4}{#1\_Desk\_4u.pdf}{#1\_Full\_4u.pdf}
\BinNestedPanel{f}{Width 5}{#1\_Desk\_5u.pdf}{#1\_Full\_5u.pdf}
}

% Captioned figure wrappers by binning scope.
% Args:
%  #1 = label suffix
%  #2 = caption
%  #3 = base filename prefix (before _Desk/_Full and _decile/_1u/...)
\newcommand{\BinCompareBinsPanelsWholeSampleFigure}[3]{%
\par\medskip
\noindent\captionof{figure}{#2}\label{fig:#1}
\par\medskip
\BinCompareBinsPanelsWholeSample{#3}
}

\newcommand{\BinCompareBinsPanelsWithinMethodFigure}[3]{%
\par\medskip
\noindent\captionof{figure}{#2}\label{fig:#1}
\par\medskip
\BinCompareBinsPanelsWithinMethod{#3}
}

% Figure block for executed-case outputs with outcome-specific quantile bins.
% Args:
%  #1 = base filename prefix (before _Desk/_Full and _decile/_quintile)
\newcommand{\BinCompareQuantilePanelsOutcomeAvail}[1]{%
\BinNestedPanel{a}{Deciles}{#1\_Desk\_decile\_outcomeavail.pdf}{#1\_Full\_decile\_outcomeavail.pdf}
\BinNestedPanel{b}{Quintiles}{#1\_Desk\_quintile\_outcomeavail.pdf}{#1\_Full\_quintile\_outcomeavail.pdf}
}

\newcommand{\BinCompareQuantilePanelsOutcomeAvailFigure}[3]{%
\par\medskip
\noindent\captionof{figure}{#2}\label{fig:#1}
\par\medskip
\BinCompareQuantilePanelsOutcomeAvail{#3}
}

%%% helper to add several binplot figures:
\newcommand{\BinCompare}[4]{%
\begin{figure}[!htbp]
\caption{#2}
\label{fig:#1}
\centering
\begin{subfigure}[t]{0.48\textwidth}
  \centering
  \includegraphics[width=\linewidth]{Output/#3}
  \subcaption{Desk audits}
\end{subfigure}\hfill
\begin{subfigure}[t]{0.48\textwidth}
  \centering
  \includegraphics[width=\linewidth]{Output/#4}
  \subcaption{Full audits}
\end{subfigure}
\end{figure}
}
% Main: 5 bin definitions stacked, Desk vs Full
% Args:
%  #1 = label suffix (used in \label{fig:...})
%  #2 = caption
%  #3 = base filename prefix (everything before _Desk/_Full, incl. variable)
%       e.g. avg_ninspectors_binplot_yhatrf
\newcommand{\BinCompareBins}[3]{%
\begin{figure}[!htbp]
\centering
\caption{#2}
\label{fig:#1}
\setlength{\tabcolsep}{3pt}
\renewcommand{\arraystretch}{1.12}
\begin{tabular}{@{}lcc@{}}
\toprule
 & \textbf{Desk audits} & \textbf{Full audits} \\
\midrule
\BinRow{Deciles}{#3\_Desk\_decile.pdf}{#3\_Full\_decile.pdf}
\BinRow{Width 1}{#3\_Desk\_1u.pdf}{#3\_Full\_1u.pdf}
\BinRow{Width 2}{#3\_Desk\_2u.pdf}{#3\_Full\_2u.pdf}
\BinRow{Width 4}{#3\_Desk\_4u.pdf}{#3\_Full\_4u.pdf}
\BinRow{Width 5}{#3\_Desk\_5u.pdf}{#3\_Full\_5u.pdf}
\bottomrule
\end{tabular}
\end{figure}
}

\newcommand{\sym}[1]{\ifmmode^{#1}\else\(^{#1}\)\fi}
\begin{document}

\maketitle
\tableofcontents
\listoffigures
\listoftables

\newpage
\section{New exhibits for paper}

\subsection{Regression Analysis - Description of Approach}

\textbf{Objective.} Using predicted evasion estimates generated by the Random Forest model, we test whether execution of algorithm-assigned cases depends on whether the "first-in-line" for execution algorithm case(s) on a given list are high- vs. low-potential. The "list" unit differs by audit type: for desk audits it is defined at the inspector $\times$ selection-year level, while for full audits it is defined at the bureau (detailed) $\times$ selection-year level. We compare execution rates of algorithm cases across these lists.

\textbf{Steps.}
\begin{enumerate}
    \item \textbf{Load predicted evasion estimates data which includes estimates for selected audits, and restrict by audit type (full or desk audit).} 
    We use the dataset containing predicted evasion (\texttt{yhatrf}) and execution outcomes, and split the analysis by audit type (\texttt{x2==0} desk; \texttt{x2==1} full).

    \item \textbf{Construct predicted-evasion quantile groups (audit-type specific).}
    We compute quintiles of predicted evasion \texttt{yhatrf} in the following way:
    \begin{itemize}
        \item \textbf{Desk audits:} quintiles within \texttt{verificateur\_selection $\times$ selectionyear}.
        \item \textbf{Full audits:} quintiles within \texttt{bureau\_detailed $\times$ selectionyear}.
    \end{itemize}

    \item \textbf{Restrict to algorithm-assigned cases.}
    We keep only algorithm-selected audits (\texttt{algorithm==1}) so that all outcomes and tags refer to the execution of algorithm-assigned cases.

    \item \textbf{Define "top of the list" using execution sequence among algorithm cases.}
    Within each list, we order algorithm cases by the sequencing variable (\texttt{sequencing}) to define early algorithm cases (e.g., Top 1/Top 2, first 10\% or first 20\% of the algorithm sequence).

    \item \textbf{Tag lists based on the potential of first-in-line algorithm cases.}
    For each list, we create indicators for whether an first-in-line algorithm case is:
    (i) \textbf{high-potential} (top quintile of \texttt{yhatrf}, \texttt{yhatrf\_q5==5}), or 
    (ii) \textbf{low-potential} (bottom quintile, \texttt{yhatrf\_q5==1}).
    This generates list-level groups such as "high-potential-first" vs "low-potential-first".

    \item \textbf{Collapse to the list level and compute execution rates.}
    We collapse the data to the list level and compute the mean execution rate of assigned algorithm cases within each list (mean of \texttt{y2} among \texttt{algorithm==1}).

    \item \textbf{Compare execution outcomes across list groups.}
    We present (i) boxplots of list-level algorithm execution rates (unweighted and weighted by the number of assigned algorithm cases), and (ii) regression estimates comparing execution rates across "high-potential-first" and "low-potential-first" lists (with fixed effects varying by audit type).
\end{enumerate}


\begin{landscape}


%%%  Main table 
\begin{table}[H]\centering
\begin{adjustbox}{max width=1.4\textwidth}
\begin{threeparttable}
\caption{Execution of algorithm-assigned audits and high predicted evasion cases at the top of the list [Full audits]}
\label{tab:exec_rate_full_panels}

\footnotesize
\begin{tabular}{l*{9}{c}} \toprule
            &\multicolumn{3}{c}{No FE}                                        &\multicolumn{3}{c}{Selection year FE}                            &\multicolumn{3}{c}{Bureau FE}                                    \\\cmidrule(lr){2-4}\cmidrule(lr){5-7}\cmidrule(lr){8-10}
\multicolumn{1}{l}{\textbf{Outcome: Execution Rate of Algorithm-Selected Cases}} & Top 2 & First 10\% seq & First 20\% seq & Top 2 & First 10\% seq & First 20\% seq & Top 2 & First 10\% seq & First 20\% seq \\ \midrule \multicolumn{10}{l}{\textbf{Panel A: Discrete measures}}\\ \addlinespace
High potential evasion case in list =1&       0.218\sym{*}  &       0.289\sym{**} &       0.207\sym{*}  &       0.176         &       0.224\sym{**} &       0.199\sym{**} &       0.144         &       0.245         &       0.165         \\
            &     (0.124)         &     (0.123)         &     (0.108)         &     (0.105)         &     (0.102)         &     (0.087)         &     (0.185)         &     (0.173)         &     (0.153)         \\
\addlinespace
Low potential evasion case in list =1&      -0.079         &       0.044         &      -0.108         &      -0.079         &      -0.007         &      -0.072         &      -0.067         &       0.085         &      -0.049         \\
            &     (0.108)         &     (0.112)         &     (0.116)         &     (0.105)         &     (0.110)         &     (0.112)         &     (0.122)         &     (0.127)         &     (0.171)         \\
\midrule
Selection Year FE&          No         &          No         &          No         &         Yes         &         Yes         &         Yes         &          No         &          No         &          No         \\
Bureau FE   &          No         &          No         &          No         &          No         &          No         &          No         &         Yes         &         Yes         &         Yes         \\
Observations&          29         &          29         &          29         &          29         &          29         &          29         &          29         &          29         &          29         \\
R-squared   &       0.129         &       0.118         &       0.208         &       0.366         &       0.350         &       0.420         &       0.424         &       0.431         &       0.430         \\

\addlinespace \midrule \multicolumn{10}{l}{\textbf{Panel B: Continuous measure}}\\ \addlinespace

Share of high predicted evasion case at the top&       0.532\sym{**} &       0.530\sym{**} &       0.556\sym{***}&       0.457\sym{**} &       0.403         &       0.481\sym{**} &       0.398         &       0.492         &       0.444         \\
            &     (0.225)         &     (0.240)         &     (0.187)         &     (0.190)         &     (0.238)         &     (0.216)         &     (0.325)         &     (0.333)         &     (0.310)         \\
\midrule
Selection Year FE&          No         &          No         &          No         &         Yes         &         Yes         &         Yes         &          No         &          No         &          No         \\
Bureau FE   &          No         &          No         &          No         &          No         &          No         &          No         &         Yes         &         Yes         &         Yes         \\
Observations&          29         &          29         &          29         &          29         &          29         &          29         &          29         &          29         &          29         \\
R-squared   &       0.112         &       0.076         &       0.097         &       0.350         &       0.313         &       0.343         &       0.416         &       0.422         &       0.417         \\
\bottomrule \end{tabular}


\begin{tablenotes}[flushleft]
\footnotesize
\item Notes: This table reports regression estimates of how the presence of high- or low-predicted-evasion (high-/low-potential) cases at the top of an algorithm list is associated with the execution rate of algorithm-assigned audits. The unit of observation is a bureau-by-selection-year list $(b,t)$. The dependent variable is $y_{bt}$, the execution rate of algorithm-selected cases in bureau $b$ and selection year $t$ (executed algorithm-selected cases divided by assigned algorithm-selected cases). To classify the audit case potential we use the  estimated  predicted evasion from the Random Forest model and partitioned it into quintiles within tax office $\times$ selection year; a  "high" predicted evasion corresponds to the top quintile and "low"' predicted evasion to the bottom quintile. Panel A presents the results of the  discrete specification of an OLS regression through the following equation:  $y_{bt}=\alpha+\beta_H\,\text{HighTop}_{bt}+\beta_L\,\text{LowTop}_{bt}+\delta_t+\mu_b+\varepsilon_{bt}$, where $\text{HighTop}_{bt}=1$ if the top of the ranking contains a high-predicted-evasion case under the relevant top-of-list definition (Top 2, first 10\% of the sequence, first 20\% of the sequence), and 0 otherwise; $\text{LowTop}_{bt}$ is defined analogously for low-predicted-evasion cases. Panel B presents the results of an alternative specification using a continuous measure of high potential cases, by estimating: $y_{bt}=\alpha+\theta\,\text{ShareHighTop}_{bt}+\delta_t+\mu_b+\varepsilon_{bt}$, where $\text{ShareHighTop}_{bt}$ is the share of high-predicted-evasion cases at the top of the list (same top-of-list definitions as in Panel A). Columns labeled No FE omit fixed effects, Selection Year FE include $\delta_t$ only, and Bureau FE include $\mu_b$ only. Standard errors are robust. Stars denote significance at $^{*}p<0.10$, $^{**}p<0.05$, $^{***}p<0.01$.
\end{tablenotes}
\end{threeparttable}
\end{adjustbox}
\end{table}

%%% Adding robustness check with regression similar as table 1 but for inspector selected audits. We should expect no response of inspectors on poor audit case if they we inspector selected. 
\begin{comment}
\begin{table}[H]\centering
\begin{adjustbox}{max width=1.4\textwidth}
\begin{threeparttable}
\caption{Execution of inspector-assigned audits and high predicted evasion cases at the top of the list [Full audits]}
\label{tab:exec_rate_full_panels_inspector}

\footnotesize
\begin{tabular}{l*{9}{c}} \toprule
            &\multicolumn{3}{c}{No FE}                                        &\multicolumn{3}{c}{Selection year FE}                            &\multicolumn{3}{c}{Bureau FE}                                    \\\cmidrule(lr){2-4}\cmidrule(lr){5-7}\cmidrule(lr){8-10}
\multicolumn{1}{l}{\textbf{Outcome: Execution Rate of Inspector-Selected Cases}} & Top 2 & First 10\% seq & First 20\% seq & Top 2 & First 10\% seq & First 20\% seq & Top 2 & First 10\% seq & First 20\% seq \\ \midrule \multicolumn{10}{l}{\textbf{Panel A: Discrete measures}}\\ \addlinespace
High potential evasion case in list =1&      -0.126         &      -0.009         &      -0.153\sym{*}  &      -0.078         &      -0.018         &      -0.163         &      -0.102         &      -0.070         &      -0.249         \\
            &     (0.118)         &     (0.122)         &     (0.086)         &     (0.100)         &     (0.100)         &     (0.144)         &     (0.134)         &     (0.143)         &     (0.173)         \\
\addlinespace
Low potential evasion case in list =1&      -0.058         &       0.088         &      -0.034         &      -0.090         &      -0.090         &      -0.210         &       0.038         &       0.018         &      -0.077         \\
            &     (0.122)         &     (0.145)         &     (0.116)         &     (0.104)         &     (0.121)         &     (0.168)         &     (0.166)         &     (0.165)         &     (0.188)         \\
\midrule
Selection Year FE&          No         &          No         &          No         &         Yes         &         Yes         &         Yes         &          No         &          No         &          No         \\
Bureau FE   &          No         &          No         &          No         &          No         &          No         &          No         &         Yes         &         Yes         &         Yes         \\
Observations&          29         &          29         &          29         &          29         &          29         &          29         &          29         &          29         &          29         \\
R-squared   &       0.044         &       0.014         &       0.053         &       0.413         &       0.401         &       0.441         &       0.406         &       0.382         &       0.454         \\

\addlinespace \midrule \multicolumn{10}{l}{\textbf{Panel B: Continuous measure}}\\ \addlinespace

Share of high predicted evasion case at the top&      -0.217         &      -0.233         &      -0.504\sym{*}  &      -0.107         &      -0.114         &      -0.242         &      -0.225         &      -0.271         &      -0.449\sym{*}  \\
            &     (0.227)         &     (0.212)         &     (0.261)         &     (0.189)         &     (0.159)         &     (0.259)         &     (0.221)         &     (0.196)         &     (0.213)         \\
\midrule
Selection Year FE&          No         &          No         &          No         &         Yes         &         Yes         &         Yes         &          No         &          No         &          No         \\
Bureau FE   &          No         &          No         &          No         &          No         &          No         &          No         &         Yes         &         Yes         &         Yes         \\
Observations&          29         &          29         &          29         &          29         &          29         &          29         &          29         &          29         &          29         \\
R-squared   &       0.038         &       0.047         &       0.109         &       0.400         &       0.402         &       0.414         &       0.405         &       0.413         &       0.439         \\
\bottomrule \end{tabular}


\begin{tablenotes}[flushleft]
\footnotesize
\item Notes: This table reports regression estimates of how the presence of high- or low-predicted-evasion (high-/low-potential) cases at the top of an inspector list is associated with the execution rate of inspector-assigned audits. The unit of observation is a bureau-by-selection-year list $(b,t)$. The dependent variable is $y_{bt}$, the execution rate of inspector-selected cases in bureau $b$ and selection year $t$ (executed inspector-selected cases divided by assigned inspector-selected cases). To classify the audit case potential we use the  estimated  predicted evasion from the Random Forest model and partitioned it into quintiles within tax office $\times$ selection year; a  "high" predicted evasion corresponds to the top quintile and "low"' predicted evasion to the bottom quintile. Panel A presents the results of the discrete specification of an OLS regression through the following equation: $y_{bt}=\alpha+\beta_H\,\text{HighTop}_{bt}+\beta_L\,\text{LowTop}_{bt}+\delta_t+\mu_b+\varepsilon_{bt}$, where $\text{HighTop}_{bt}=1$ if the top of the ranking contains a high-predicted-evasion case under the relevant top-of-list definition (Top 2, first 10\% of the sequence, first 20\% of the sequence), and 0 otherwise; $\text{LowTop}_{bt}$ is defined analogously for low-predicted-evasion cases. Panel B presents the results of an alternative specification using a continuous measure of high potential cases, by estimating: $y_{bt}=\alpha+\theta\,\text{ShareHighTop}_{bt}+\delta_t+\mu_b+\varepsilon_{bt}$, where $\text{ShareHighTop}_{bt}$ is the share of high-predicted-evasion cases at the top of the list (same top-of-list definitions as in Panel A). Columns labeled No FE omit fixed effects, Selection Year FE include $\delta_t$ only, and Bureau FE include $\mu_b$ only. Standard errors are robust. Stars denote significance at $^{*}p<0.10$, $^{**}p<0.05$, $^{***}p<0.01$.
\end{tablenotes}
\end{threeparttable}
\end{adjustbox}
\end{table}
\end{comment}
\end{landscape}

\begin{comment}

\begin{figure}[H]
    \centering
    \caption{Execution Rates and Audit Success, Algorithm-Selected Full Audits}
    \begin{threeparttable}
        

    \begin{tabular}{cc}
    \multicolumn{2}{c}{A: Share of Void Cases} \\
    A1: 2018 vs 2019     &  A2: 2018 vs 2019-2020 average\\
    \includegraphics[width=0.5\linewidth]{Output/scatter_share_v_Algorithm_2018_2019_full.pdf}  &  \includegraphics[width=0.5\linewidth]{Output/scatter_share_v_Algorithm_2018_avg_full.pdf} \\
    \multicolumn{2}{c}{B: Below-Median Evasion Rate} \\
    B1: 2018 vs 2019     &  B2: 2018 vs 2019-2020 average\\
    \includegraphics[width=0.5\linewidth]{Output/scatter_share_med_er_Algorithm_2018_2019_full.pdf}  &  \includegraphics[width=0.5\linewidth]{Output/scatter_share_med_er_Algorithm_2018_avg_full.pdf} \\
    \multicolumn{2}{c}{C: Bottom-Quartile Evasion Rate} \\
    C1: 2018 vs 2019     &  C2: 2018 vs 2019-2020 average\\
    \includegraphics[width=0.5\linewidth]{Output/scatter_share_botq_er_Algorithm_2018_2019_full.pdf} & 
    \includegraphics[width=0.5\linewidth]{Output/scatter_share_botq_er_Algorithm_2018_avg_full.pdf}   \\
    \end{tabular}
    \begin{tablenotes}
    \footnotesize \item Notes: This figure relates tax office-level exposure in year $t$ to low-yield algorithm-selected audits to the execution rate of algorithm-selected audits in year $t+1$. The sample includes selected full audits only (excluding safety cases) and is aggregated at the tax office $\times$ year level. The estimation sample keeps offices observed in at least two periods in the relevant comparison window (2018 vs 2019, or 2018 vs the 2019--2020 average specification). An audit is classified as void when it is executed and uncovered evasion value equals zero. Evasion rate is defined for executed audits as uncovered evasion value divided by uncovered evasion value plus average turnover; average turnover is constructed as the average reported sales (2014--2020). Panel A uses the share of void audits; Panel B the share of audits below the median evasion-rate cutoff; Panel C the share in the bottom quartile of the evasion-rate distribution. Median and quartile evasion-rate cutoffs are computed within selection-year $\times$ selection method. For Panels A--C, low-yield shares on the horizontal axis use executed algorithm-selected audits as denominators. The vertical-axis execution measure is the algorithm execution rate, defined as executed algorithm-selected audits divided by assigned algorithm-selected audits. Column 1 compares 2018 and 2019. Column 2 compares 2018 with the average of 2019 and 2020, computed as the row mean. Points in the scatterplot represent tax offices observed during the whole analysis period (2018--2020), and lines show unweighted linear fits.

    \end{tablenotes}
        \end{threeparttable}
    \label{fig:effortresonse-scatter}
%    \begin{figuresnotes}
        
\end{figure} 
\end{comment}

%%% Figure two: Shares relative to inspectors selected 
\begin{comment}
\begin{figure}[H]
    \centering
    \caption{Effort Response to Audit Success, Algorithm-Selected Full Audits (Relative to Inspector Selection)}
    \label{fig:effortresponse-scatter-relins}
    \begin{threeparttable}
        

    \begin{tabular}{cc}
    \multicolumn{2}{c}{A: Share of Void Cases (Algorithm / Inspectors)} \\
    A1: 2018 vs 2019     &  A2: 2018 vs 2019-2020 average\\
    \includegraphics[width=0.5\linewidth]{Output/scatter_share_v_relIns_Algorithm_2018_2019_full.pdf}  &
    \includegraphics[width=0.5\linewidth]{Output/scatter_share_v_relIns_Algorithm_2018_avg_full.pdf} \\
    \multicolumn{2}{c}{B: Below-Median Evasion Rate (Algorithm / Inspectors)} \\
    B1: 2018 vs 2019     &  B2: 2018 vs 2019-2020 average\\
    \includegraphics[width=0.5\linewidth]{Output/scatter_share_med_er_relIns_Algorithm_2018_2019_full.pdf}  &
    \includegraphics[width=0.5\linewidth]{Output/scatter_share_med_er_relIns_Algorithm_2018_avg_full.pdf} \\
    \multicolumn{2}{c}{C: Bottom-Quartile Evasion Rate (Algorithm / Inspectors)} \\
    C1: 2018 vs 2019     &  C2: 2018 vs 2019-2020 average\\
    \includegraphics[width=0.5\linewidth]{Output/scatter_share_botq_er_relIns_Algorithm_2018_2019_full.pdf} &
    \includegraphics[width=0.5\linewidth]{Output/scatter_share_botq_er_relIns_Algorithm_2018_avg_full.pdf} \\
    \end{tabular}
    \begin{tablenotes}
\footnotesize \item Notes: This figure relates tax office level exposure in year $t$ to low-yield algorithm-selected audits to the relative execution rate of algorithm-selected audits in year $t+1$, with all measures expressed as Algorithm relative to Inspectors within the same tax office-year. The sample includes selected full audits only, and is aggregated at the tax office $\times$ year level. The estimation sample further keeps only offices observed in at least two periods in the relevant comparison window (2018 vs 2019, or the 2018 vs 2019--2020 average specification). An audit is classified as void when it is executed and uncovered evasion value equals zero. Evasion rate is defined for executed audits as uncovered evasion value divided by uncovered evasion value plus average turnover; average turnover is constructed as the average reported sales (2014-2020). Panel A uses the share of void audits; Panel B the share of audits below the median evasion-rate cutoff; Panel C the share in the bottom quartile of the evasion-rate distribution. Median and quartile evasion rate cutoffs are computed within selection-year and  selection method. For Panels A--C, low yield shares use method-specific executed audits as denominators and are presented in the horizontal axis. The vertical-axis execution measure uses method-specific assigned audits as denominator. In the relative specification, each Algorithm measure is divided by the corresponding Inspectors measure. Ratios are set to missing when inspector denominators are zero, and observations with zero method-specific denominators are excluded for the corresponding measure. Column 1 compares 2018 and 2019. Column 2 compares 2018 with the average of 2019 and 2020, computed as the row mean. Points in the scatterplot represent tax offices observed during the whole analysis period (2018--2020), and lines show unweighted linear fits. 
    \end{tablenotes}
    \end{threeparttable}
\end{figure}
\end{comment}


\begin{figure}[H]
    \centering
    \caption{Predicted evasion vs. number of inspectors and duration, Full audits [Executed cases sample -- preferred]}
    \label{fig:binplots_pred_evasion_effort_full_whole_sample}
    \begin{threeparttable}
    \setlength{\tabcolsep}{2pt}
    \renewcommand{\arraystretch}{0.95}
    \begin{tabular}{@{}cc@{}}
    \multicolumn{2}{c}{A: Average number of inspectors} \\
    A1: Deciles of predicted evasion     &  A2: Quintiles of predicted evasion\\
    \includegraphics[width=0.46\linewidth]{Output/avg_ninspectors_binplot_yhatrf_Full_decile_whole_sample.pdf}  &
    \includegraphics[width=0.46\linewidth]{Output/avg_ninspectors_binplot_yhatrf_Full_quintile_whole_sample.pdf} \\
    \multicolumn{2}{c}{B: Average duration [Admin. data]} \\
    B1: Deciles of predicted evasion     &  B2: Quintiles of predicted evasion\\
    \includegraphics[width=0.46\linewidth]{Output/avg_duration_admin_binplot_yhatrf_Full_decile_whole_sample.pdf}  &
    \includegraphics[width=0.46\linewidth]{Output/avg_duration_admin_binplot_yhatrf_Full_quintile_whole_sample.pdf} \\
    \multicolumn{2}{c}{C: Average duration [Taxpayer data]} \\
    C1: Deciles of predicted evasion     &  C2: Quintiles of predicted evasion\\
    \includegraphics[width=0.46\linewidth]{Output/avg_duration_tp_binplot_yhatrf_Full_decile_whole_sample.pdf} &
    \includegraphics[width=0.46\linewidth]{Output/avg_duration_tp_binplot_yhatrf_Full_quintile_whole_sample.pdf} \\
    \multicolumn{2}{c}{D: Average duration [Inspector survey]} \\
    D1: Deciles of predicted evasion     &  D2: Quintiles of predicted evasion\\
    \includegraphics[width=0.46\linewidth]{Output/avg_duration_self_binplot_yhatrf_Full_decile_whole_sample.pdf} &
    \includegraphics[width=0.46\linewidth]{Output/avg_duration_self_binplot_yhatrf_Full_quintile_whole_sample.pdf} \\
    \end{tabular}
    \begin{tablenotes}
    
    \tiny \item Notes: Predicted-evasion bins are defined using the pooled sample of selected full audits and are then kept fixed when computing outcomes for executed cases. Each panel reports mean outcomes by selection method and predicted-evasion bin, comparing algorithm-selected and inspector-selected cases. Panel A shows the number of inspectors; Panel B administrative duration; Panel C taxpayer-reported duration; and Panel D inspector/self-reported duration. Shaded bands indicate 95\% confidence intervals for each method-by-bin cell. Column 1 uses predicted-evasion deciles, while Column 2 uses predicted-evasion quintiles.
    \end{tablenotes}
    \end{threeparttable}
\end{figure}

%% same figure but with bins constructed on executed audits  + outcome availability subsample:
\begin{figure}[H]
    \centering
    \caption{Predicted evasion vs. number of inspectors and duration, Full audits [Outcome-specific executed sample]}
    \label{fig:binplots_pred_evasion_effort_full_outcomeavail}
    \begin{threeparttable}
    \setlength{\tabcolsep}{2pt}
    \renewcommand{\arraystretch}{0.95}
    \begin{tabular}{@{}cc@{}}
    \multicolumn{2}{c}{A: Average number of inspectors} \\
    A1: Deciles of predicted evasion     &  A2: Quintiles of predicted evasion\\
    \includegraphics[width=0.46\linewidth]{Output/avg_ninspectors_binplot_yhatrf_Full_decile_outcomeavail.pdf}  &
    \includegraphics[width=0.46\linewidth]{Output/avg_ninspectors_binplot_yhatrf_Full_quintile_outcomeavail.pdf} \\
    \multicolumn{2}{c}{B: Average duration [Admin. data]} \\
    B1: Deciles of predicted evasion     &  B2: Quintiles of predicted evasion\\
    \includegraphics[width=0.46\linewidth]{Output/avg_duration_admin_binplot_yhatrf_Full_decile_outcomeavail.pdf}  &
    \includegraphics[width=0.46\linewidth]{Output/avg_duration_admin_binplot_yhatrf_Full_quintile_outcomeavail.pdf} \\
    \multicolumn{2}{c}{C: Average duration [Taxpayer data]} \\
    C1: Deciles of predicted evasion     &  C2: Quintiles of predicted evasion\\
    \includegraphics[width=0.46\linewidth]{Output/avg_duration_tp_binplot_yhatrf_Full_decile_outcomeavail.pdf} &
    \includegraphics[width=0.46\linewidth]{Output/avg_duration_tp_binplot_yhatrf_Full_quintile_outcomeavail.pdf} \\
    \multicolumn{2}{c}{D: Average duration [Inspector survey]} \\
    D1: Deciles of predicted evasion     &  D2: Quintiles of predicted evasion\\
    \includegraphics[width=0.46\linewidth]{Output/avg_duration_self_binplot_yhatrf_Full_decile_outcomeavail.pdf} &
    \includegraphics[width=0.46\linewidth]{Output/avg_duration_self_binplot_yhatrf_Full_quintile_outcomeavail.pdf} \\
    \end{tabular}
    \begin{tablenotes}
    
    \tiny \item Notes: Predicted-evasion bins are defined separately for each outcome using the pooled eligible sample of executed full audits with non-missing predicted evasion and non-missing outcome data, and these outcome-specific cutoffs are then applied to both algorithm-selected and inspector-selected cases. Panel A shows the number of inspectors; Panel B administrative duration; Panel C taxpayer-reported duration; and Panel D inspector/self-reported duration. Shaded bands indicate 95\% confidence intervals for each method-by-bin cell. Column 1 uses predicted-evasion deciles, while Column 2 uses predicted-evasion quintiles.
    \end{tablenotes}
    \end{threeparttable}
\end{figure}

\begin{figure}[H]
    \centering
    \caption{Execution Rates and Audit Success, Algorithm-Selected Full Audits}
    \label{fig:effortresponse-scatter-success}
    \begin{threeparttable}
    \begin{tabular}{cc}
    \multicolumn{2}{c}{A: Share of Detected Audits} \\
    A1: 2018 vs 2019     &  A2: 2018 vs 2019-2020 average\\
    \includegraphics[width=0.5\linewidth]{Output/scatter_share_det_Algorithm_2018_2019_full_success.pdf}  &
    \includegraphics[width=0.5\linewidth]{Output/scatter_share_det_Algorithm_2018_avg_full_success.pdf} \\
    \multicolumn{2}{c}{B: Above-Median Evasion Rate} \\
    B1: 2018 vs 2019     &  B2: 2018 vs 2019-2020 average\\
    \includegraphics[width=0.5\linewidth]{Output/scatter_share_abvmed_er_Algorithm_2018_2019_full_success.pdf}  &
    \includegraphics[width=0.5\linewidth]{Output/scatter_share_abvmed_er_Algorithm_2018_avg_full_success.pdf} \\
    \multicolumn{2}{c}{C: Top-Quartile Evasion Rate} \\
    C1: 2018 vs 2019     &  C2: 2018 vs 2019-2020 average\\
    \includegraphics[width=0.5\linewidth]{Output/scatter_share_topq_er_Algorithm_2018_2019_full_success.pdf} &
    \includegraphics[width=0.5\linewidth]{Output/scatter_share_topq_er_Algorithm_2018_avg_full_success.pdf} \\
    \end{tabular}
    \begin{tablenotes}
    \footnotesize \item Notes: This figure relates tax office-level exposure in year $t$ to successful algorithm-selected audits to the execution rate of algorithm-selected audits in year $t+1$. The sample includes selected full audits only (excluding safety cases) and is aggregated at the tax office $\times$ year level. The estimation sample keeps offices observed in at least two periods in the relevant comparison window (2018 vs 2019, or 2018 vs the 2019--2020 average specification). An audit is classified as detected when it is executed and uncovered evasion value is strictly positive (complement of a void audit). Evasion rate is defined for executed audits as uncovered evasion value divided by uncovered evasion value plus average turnover; average turnover is constructed as the average reported sales (2014--2020). Panel A uses the share of detected audits; Panel B the share of audits above the median evasion-rate cutoff; Panel C the share in the top quartile of the evasion-rate distribution. Median and quartile evasion-rate cutoffs are computed within selection-year $\times$ selection method. For Panels A--C, shares on the horizontal axis use executed algorithm-selected audits as denominators. The vertical-axis execution measure is the algorithm execution rate, defined as executed algorithm-selected audits divided by assigned algorithm-selected audits. Column 1 compares 2018 and 2019. Column 2 compares 2018 with the average of 2019 and 2020, computed as the row mean. Points represent tax offices observed during the whole analysis period (2018--2020), and lines show unweighted linear fits.
    \end{tablenotes}
    \end{threeparttable}
\end{figure}

%%%% New figure with just 3 panels comparing 2018 vs 2019 for full audits.
\begin{figure}[H]
    \centering
    \caption{Execution Rates and Audit Success, Algorithm-Selected Full Audits: 2018 vs 2019}
    \label{fig:effortresponse-scatter-success}

    \begin{tabular}{cc}
    A: Share of Detected Audits &
    B: Above-Median Evasion Rate \\[0.3em]
    \includegraphics[width=0.42\textwidth]{Output/scatter_share_det_Algorithm_2018_2019_full_success.pdf} &
    \includegraphics[width=0.42\textwidth]{Output/scatter_share_abvmed_er_Algorithm_2018_2019_full_success.pdf} \\
    \end{tabular}

    \vspace{0.5em}

    \begin{tabular}{c}
    C: Top-Quartile Evasion Rate \\[0.3em]
    \includegraphics[width=0.42\textwidth]{Output/scatter_share_topq_er_Algorithm_2018_2019_full_success.pdf} \\
    \end{tabular}

    \vspace{0.5em}

    \begin{minipage}{\textwidth}
    \footnotesize
    \textit{Notes:} This figure relates tax office-level exposure in year $t$ to successful algorithm-selected audits to the execution rate of algorithm-selected audits in year $t+1$. The sample includes selected full audits only (excluding safety cases) and is aggregated at the tax office $\times$ year level. The estimation sample keeps offices observed in both 2018 and 2019. An audit is classified as detecting evasion when it is executed and the uncovered evasion value is strictly positive. Evasion rate is defined for executed audits as uncovered evasion value divided by uncovered evasion value plus average turnover; average turnover is constructed as the average reported sales (2014--2020). Panel A uses the share of audits that detect evasion, conditional on audit implementation; Panel B the share of audits above the median evasion-rate cutoff; and Panel C the share in the top quartile of the evasion-rate distribution. Median and quartile evasion-rate cutoffs are computed within tax office $\times$ selection year. In all panels, shares on the horizontal axis use executed algorithm-selected audits as denominators. The vertical-axis execution measure is the algorithm execution rate, defined as executed algorithm-selected audits divided by assigned algorithm-selected audits. Each plot compares 2018 and 2019. Points represent tax offices observed in both years, and lines show unweighted linear fits.
    \end{minipage}
\end{figure}

%%% Relative to inspectors within selection year X tax office 
\begin{figure}[H]
    \centering
    \caption{Execution Rates and Audit Success relative to Inspector-selected Cases, Full Audits: 2018 vs 2019}
    \label{fig:effortresponse-scatter-success}

    \begin{tabular}{cc}
    A: Share of Detected Audits &
    B: Above-Median Evasion Rate \\[0.3em]
    \includegraphics[width=0.42\textwidth]{Output/scatter_share_det_relIns_Algorithm_2018_2019_full_success.pdf} &
    \includegraphics[width=0.42\textwidth]{Output/scatter_share_abvmed_er_relIns_Algorithm_2018_2019_full_success.pdf} \\
    \end{tabular}

    \vspace{0.5em}

    \begin{tabular}{c}
    C: Top-Quartile Evasion Rate \\[0.3em]
    \includegraphics[width=0.42\textwidth]{Output/scatter_share_topq_er_relIns_Algorithm_2018_2019_full_success.pdf} \\
    \end{tabular}

    \vspace{0.5em}

    \begin{minipage}{\textwidth}
    \footnotesize
    \textit{Notes:} This figure relates tax office-level exposure in year $t$ to successful algorithm-selected audits to the execution rate of algorithm-selected audits in year $t+1$. The sample includes selected full audits only (excluding safety cases) and is aggregated at the tax office $\times$ year level. The estimation sample keeps offices observed in both 2018 and 2019. An audit is classified as detecting evasion when it is executed and the uncovered evasion value is strictly positive. Evasion rate is defined for executed audits as uncovered evasion value divided by uncovered evasion value plus average turnover; average turnover is constructed as the average reported sales (2014--2020). Panel A uses the share of audits that detect evasion, conditional on audit implementation, relative to the corresponding share among inspector-selected audits; Panel B uses the share of audits above the median evasion-rate cutoff, relative to the corresponding share among inspector-selected audits; and Panel C uses the share in the top quartile of the evasion-rate distribution, relative to the corresponding share among inspector-selected audits. Median and quartile evasion-rate cutoffs are computed within tax office $\times$ selection year. In all panels, the horizontal axis measures the algorithm-selected share relative to the corresponding inspector-selected share within the same tax office and year. The vertical-axis execution measure is the algorithm execution rate, defined as executed algorithm-selected audits divided by assigned algorithm-selected audits. Each plot compares 2018 and 2019. Points represent tax offices observed in both years, and lines show unweighted linear fits.
    \end{minipage}
    \end{figure}



\begin{comment}
\begin{figure}[H]
    \centering
    \caption{Effort Response to Audit Success, Algorithm-Selected Full Audits (Relative to Inspector Selection)}
    \label{fig:effortresponse-scatter-relins-success}
    \begin{threeparttable}
    \begin{tabular}{cc}
    \multicolumn{2}{c}{A: Share of Detected Audits (Algorithm / Inspectors)} \\
    A1: 2018 vs 2019     &  A2: 2018 vs 2019-2020 average\\
    \includegraphics[width=0.5\linewidth]{Output/scatter_share_det_relIns_Algorithm_2018_2019_full_success.pdf}  &
    \includegraphics[width=0.5\linewidth]{Output/scatter_share_det_relIns_Algorithm_2018_avg_full_success.pdf} \\
    \multicolumn{2}{c}{B: Above-Median Evasion Rate (Algorithm / Inspectors)} \\
    B1: 2018 vs 2019     &  B2: 2018 vs 2019-2020 average\\
    \includegraphics[width=0.5\linewidth]{Output/scatter_share_abvmed_er_relIns_Algorithm_2018_2019_full_success.pdf}  &
    \includegraphics[width=0.5\linewidth]{Output/scatter_share_abvmed_er_relIns_Algorithm_2018_avg_full_success.pdf} \\
    \multicolumn{2}{c}{C: Top-Quartile Evasion Rate (Algorithm / Inspectors)} \\
    C1: 2018 vs 2019     &  C2: 2018 vs 2019-2020 average\\
    \includegraphics[width=0.5\linewidth]{Output/scatter_share_topq_er_relIns_Algorithm_2018_2019_full_success.pdf} &
    \includegraphics[width=0.5\linewidth]{Output/scatter_share_topq_er_relIns_Algorithm_2018_avg_full_success.pdf} \\
    \end{tabular}
    \begin{tablenotes}
    \footnotesize \item Notes: This figure relates tax office-level exposure in year $t$ to successful algorithm-selected audits to the relative execution rate of algorithm-selected audits in year $t+1$, with all measures expressed as Algorithm relative to Inspectors within the same tax office-year. The sample includes selected full audits only and is aggregated at the tax office $\times$ year level. The estimation sample keeps only offices observed in at least two periods in the relevant comparison window (2018 vs 2019, or 2018 vs 2019--2020 average). An audit is classified as detected when it is executed and uncovered evasion value is strictly positive. Panel A uses the detection share; Panel B the share above the median evasion-rate cutoff; Panel C the share in the top quartile of the evasion-rate distribution. Median and quartile evasion-rate cutoffs are computed within selection-year $\times$ selection method. For Panels A--C, shares use method-specific executed audits as denominators. The vertical-axis execution measure uses method-specific assigned audits as denominator. In the relative specification, each Algorithm measure is divided by the corresponding Inspectors measure. Ratios are set to missing when inspector denominators are zero, and observations with zero method-specific denominators are excluded for the corresponding measure. Column 1 compares 2018 and 2019. Column 2 compares 2018 with the average of 2019 and 2020. Points represent tax offices observed during 2018--2020, and lines show unweighted linear fits.
    \end{tablenotes}
    \end{threeparttable}
\end{figure}
\end{comment}


\section{Audit productivity}

\subsection{Table 4 with predicted evasion controls}

This subsection presents the Table 4 baseline replication together with variants that keep the original nine-column layout and add predicted-evasion controls.

\newcommand{\TableFourRobustnessNotes}[1]{%
\footnotesize
\item Notes: * 0.10 ** 0.05 *** 0.01 significance levels. OLS results of a regression of the main audit outcomes on the selection method of the case, for full and desk audits separately. Specifications for full audits control for tax office x year fixed effects, which corresponds to the selection list for full audits. For desk audits, we show specifications with tax office x year fixed effects, and inspector x year fixed effects, the latter corresponding to the list for desk audits. The first outcome (Columns 1 to 3) is the probability that the case on the list was actually executed by the inspectors. The second outcome (Columns 4 to 6) is the probability that the audit resulted in a penalty (initial or final notice), conditional on the case being conducted. The third outcome (Columns 7 to 9) is the log amount of the evasion plus penalties uncovered, conditional on the case ending in positive penalty. The evasion amount is obtained from the final notice, which we complement with the initial notice when the final notice is not available. The sample includes all cases selected in the 2018, 2019, and 2020 audit programs. The data includes the selection and audits of the Large Taxpayer Unit, Medium Taxpayer Units 1 and 2, Liberal Professionals, and the regional SME units of Dakar Plateau, Grand Dakar, Pikine Guediawaye, and Ngor Almadies. Robust standard errors (Huber-White) are shown in parentheses, and clustered at the inspector x year level when inspector fixed effects are included. Lee bounds are shown in brackets for the coefficient on the algorithm-selection dummy, computed based on the attrition between algorithm and inspector cases for variables after audit execution. #1 Table C.1 shows intent-to-treat effects, Table C.3 shows robustness to controlling for audit slot fixed effects, and Table C.4 shows the results when using the evasion rate (instead of the evasion amount) as the outcome. Figure D.1 shows the distribution of audit list (= inspector x year) fixed effects for columns 3, 6 and 9.
}

\begin{landscape}

\begin{table}[H]\centering
\begin{adjustbox}{max width=1.4\textwidth}
\begin{threeparttable}
\caption{Algorithm Selection and Audit Outcomes}
\label{tab:table4_main_outcomes_replicated}

\footnotesize
\input{Output/table4_main_outcomes_replicated_with_lee.tex}

\begin{tablenotes}[flushleft]
\TableFourRobustnessNotes{This replicated specification does not add predicted-evasion controls.}
\end{tablenotes}
\end{threeparttable}
\end{adjustbox}
\end{table}

\begin{table}[H]\centering
\begin{adjustbox}{max width=1.4\textwidth}
\begin{threeparttable}
\caption{Algorithm Selection and Audit Outcomes, with Predicted Evasion Control}
\label{tab:table4_main_outcomes_yhatrf_control}

\footnotesize
\begin{tabular}{lccc|ccc|ccc} \toprule
\multicolumn{1}{l}{} & \multicolumn{3}{c}{P(Execution)} & \multicolumn{3}{c}{P(Detection | Execution)} & \multicolumn{3}{c}{log(Evasion) | Detection} \\\cmidrule(lr){2-4}\cmidrule(lr){5-7}\cmidrule(lr){8-10} \multicolumn{1}{l}{} & Full audits & Desk audits & Desk audits & Full audits & Desk audits & Desk audits & Full audits & Desk audits & Desk audits \\ \multicolumn{1}{l}{} & (1) & (2) & (3) & (4) & (5) & (6) & (7) & (8) & (9) \\ \midrule
Algorithm   &       -0.16\sym{***}&       -0.04\sym{**} &       -0.04\sym{**} &        0.06\sym{***}&        0.03\sym{***}&        0.03\sym{***}&        0.05         &       -0.03         &        0.01         \\
            &      (0.03)         &      (0.02)         &      (0.02)         &      (0.01)         &      (0.01)         &      (0.01)         &      (0.10)         &      (0.09)         &      (0.10)         \\
\addlinespace
 & & & & [0.03 ***,0.07***] & & [0.02**,0.03***] & [0.05,0.09] & & [-0.08,0.08] \\  Inspectors x Overlap&        0.13\sym{*}  &        0.04         &        0.04         &        0.01         &        0.03         &        0.01         &        0.21         &       -0.13         &       -0.22         \\
            &      (0.07)         &      (0.04)         &      (0.04)         &      (0.02)         &      (0.02)         &      (0.02)         &      (0.25)         &      (0.22)         &      (0.26)         \\
\addlinespace
Algorithm x Random&                     &        0.01         &        0.01         &                     &        0.02         &        0.01         &                     &        0.02         &       -0.02         \\
            &                     &      (0.03)         &      (0.03)         &                     &      (0.02)         &      (0.02)         &                     &      (0.13)         &      (0.14)         \\
\addlinespace
Predicted evasion&        0.03\sym{***}&        0.02\sym{***}&        0.03\sym{***}&        0.08\sym{***}&        0.09\sym{***}&        0.09\sym{***}&        1.06\sym{***}&        0.67\sym{***}&        0.66\sym{***}\\
            &      (0.01)         &      (0.00)         &      (0.00)         &      (0.00)         &      (0.00)         &      (0.00)         &      (0.04)         &      (0.03)         &      (0.03)         \\
\midrule
Tax center x Year&         Yes         &         Yes         &         Yes         &         Yes         &         Yes         &         Yes         &         Yes         &         Yes         &         Yes         \\
Inspector x Year&          No         &          No         &         Yes         &          No         &          No         &         Yes         &          No         &          No         &         Yes         \\
N           &         944         &        2731         &        2731         &         508         &        1011         &         992         &         457         &         747         &         728         \\
R2          &        0.29         &        0.25         &        0.35         &        0.86         &        0.89         &        0.91         &        0.79         &        0.63         &        0.68         \\
Mean outcome&        0.53         &        0.37         &        0.37         &        0.89         &        0.73         &        0.73         &       19.26         &       17.76         &       17.73         \\
\bottomrule \end{tabular}


\begin{tablenotes}[flushleft]
\TableFourRobustnessNotes{This robustness specification additionally controls linearly for the Random-Forest predicted evasion measure (\texttt{yhatrf}).}
\end{tablenotes}
\end{threeparttable}
\end{adjustbox}
\end{table}

\begin{table}[H]\centering
\begin{adjustbox}{max width=1.4\textwidth}
\begin{threeparttable}
\caption{Algorithm Selection and Audit Outcomes, with Predicted Evasion Quadratic Control}
\label{tab:table4_main_outcomes_yhatrf_quadratic_control}

\footnotesize
\input{Output/table4_main_outcomes_yhatrf_quadratic_control_with_lee.tex}

\begin{tablenotes}[flushleft]
\TableFourRobustnessNotes{This robustness specification additionally controls for the Random-Forest predicted evasion measure (\texttt{yhatrf}) and its square.}
\end{tablenotes}
\end{threeparttable}
\end{adjustbox}
\end{table}

\begin{table}[H]\centering
\begin{adjustbox}{max width=1.4\textwidth}
\begin{threeparttable}
\caption{Algorithm Selection and Audit Outcomes, with Predicted Evasion Decile Controls}
\label{tab:table4_main_outcomes_yhatrf_deciles_control}

\footnotesize
\begin{tabular}{lccc|ccc|ccc} \toprule
\multicolumn{1}{l}{} & \multicolumn{3}{c}{P(Execution)} & \multicolumn{3}{c}{P(Detection | Execution)} & \multicolumn{3}{c}{log(Evasion) | Detection} \\\cmidrule(lr){2-4}\cmidrule(lr){5-7}\cmidrule(lr){8-10} \multicolumn{1}{l}{} & Full audits & Desk audits & Desk audits & Full audits & Desk audits & Desk audits & Full audits & Desk audits & Desk audits \\ \multicolumn{1}{l}{} & (1) & (2) & (3) & (4) & (5) & (6) & (7) & (8) & (9) \\ \midrule
Algorithm   &       -0.09\sym{***}&       -0.03\sym{*}  &       -0.03\sym{*}  &        0.04\sym{**} &        0.01         &        0.01         &        0.11         &       -0.15         &       -0.02         \\
            &      (0.03)         &      (0.02)         &      (0.02)         &      (0.02)         &      (0.01)         &      (0.02)         &      (0.11)         &      (0.10)         &      (0.11)         \\
\addlinespace
 & & & & [0.01 ,0.04***] & & [-0.01,0.02] & [0.02,0.14] & & [-0.14*,0.05] \\  Inspectors x Overlap&        0.13\sym{**} &        0.07\sym{*}  &        0.06         &        0.06\sym{*}  &        0.08\sym{**} &        0.04         &        0.50\sym{**} &       -0.05         &       -0.14         \\
            &      (0.06)         &      (0.04)         &      (0.04)         &      (0.03)         &      (0.03)         &      (0.04)         &      (0.25)         &      (0.20)         &      (0.26)         \\
\addlinespace
Algorithm x Random&                     &        0.02         &        0.03         &                     &        0.04\sym{*}  &        0.01         &                     &        0.04         &       -0.08         \\
            &                     &      (0.03)         &      (0.03)         &                     &      (0.02)         &      (0.03)         &                     &      (0.13)         &      (0.14)         \\
\midrule
Tax center x Year&         Yes         &         Yes         &         Yes         &         Yes         &         Yes         &         Yes         &         Yes         &         Yes         &         Yes         \\
Inspector x Year&          No         &          No         &         Yes         &          No         &          No         &         Yes         &          No         &          No         &         Yes         \\
N           &         944         &        2731         &        2731         &         508         &        1011         &         992         &         457         &         747         &         728         \\
R2          &        0.37         &        0.32         &        0.39         &        0.75         &        0.80         &        0.76         &        0.76         &        0.61         &        0.67         \\
Mean outcome&        0.53         &        0.37         &        0.37         &        0.89         &        0.73         &        0.73         &       19.26         &       17.76         &       17.73         \\
\bottomrule \end{tabular}


\begin{tablenotes}[flushleft]
\TableFourRobustnessNotes{This robustness specification additionally controls for pooled-sample deciles of the Random-Forest predicted evasion measure (\texttt{yhatrf}).}
\end{tablenotes}
\end{threeparttable}
\end{adjustbox}
\end{table}

\begin{table}[H]\centering
\begin{adjustbox}{max width=1.4\textwidth}
\begin{threeparttable}
\caption{Algorithm Selection and Audit Outcomes, with Predicted Evasion Quintile Controls}
\label{tab:table4_main_outcomes_yhatrf_quintiles_control}

\footnotesize
\input{Output/table4_main_outcomes_yhatrf_quintiles_control_with_lee.tex}

\begin{tablenotes}[flushleft]
\TableFourRobustnessNotes{This robustness specification additionally controls for pooled-sample quintiles of the Random-Forest predicted evasion measure (\texttt{yhatrf}).}
\end{tablenotes}
\end{threeparttable}
\end{adjustbox}
\end{table}

\end{landscape}

\subsection{Table 5 with predicted evasion controls}

This section presents the full-audit Table 5 top-panel variants for the audit-productivity discussion. The tables below keep the original sample and layout and add predicted-evasion controls in the robustness specifications.

\begin{landscape}

\begin{table}[H]\centering
\begin{adjustbox}{max width=1.4\textwidth}
\begin{threeparttable}
\caption{Resources Allocated to Audits and Audit Productivity, Full Audits (Panels A1 and B1 only)}
\label{tab:table5_fullaudits_replicated}

\footnotesize
\begin{tabular}{lcccc|cccc} \toprule
            &\multicolumn{4}{c}{\textbf{A1: Resource Outcomes}}                                     &\multicolumn{4}{c}{\textbf{B1: Productivity Outcomes}}                                 \\\cmidrule(lr){2-5}\cmidrule(lr){6-9}
\multicolumn{1}{l}{} & \shortstack{Number of Agents} & \shortstack{Duration in Days\\(Taxpayer Survey)} & \shortstack{Days from Start to Conf.\\(Admin Data)} & \shortstack{Days Working on Case\\(Self-Reported)} & \shortstack{Evasion/ Number of\\Agents} & \shortstack{Evasion/Duration\\(Taxpayer Survey)} & \shortstack{Evasion/Duration\\(Admin. Data)} & \shortstack{Evasion/Days Working\\(Self-Reported)} \\ \multicolumn{1}{l}{} & (1) & (2) & (3) & (4) & (5) & (6) & (7) & (8) \\ \midrule
Algorithm   &       -0.26\sym{***}&       -8.54\sym{**} &      -30.18\sym{**} &      -27.26         &       -0.55\sym{***}&       -0.12         &       -0.50\sym{*}  &       -0.74         \\
            &      (0.08)         &      (3.72)         &     (13.75)         &     (30.60)         &      (0.18)         &      (0.39)         &      (0.26)         &      (0.49)         \\
\addlinespace
Inspectors x Overlap&        0.04         &      -25.57         &       -9.83         &                     &        0.30         &        0.96         &       -0.01         &                     \\
            &      (0.17)         &     (21.89)         &     (27.66)         &                     &      (0.44)         &      (2.20)         &      (0.57)         &                     \\
\midrule
N           &         508         &         214         &         282         &          51         &         457         &          98         &         253         &          40         \\
R2          &        0.21         &        0.25         &        0.30         &        0.25         &        0.31         &        0.24         &        0.25         &        0.40         \\
Mean outcome&        2.87         &       28.06         &      160.07         &      177.07         &       18.19         &       15.90         &       15.00         &       13.62         \\
\bottomrule \end{tabular}


\begin{tablenotes}[flushleft]
\footnotesize
\item Notes: This table replicates the top panels (A1 and B1) of the full-audit version of Table 5 using the original specification. Panel A1 reports resource outcomes and Panel B1 reports productivity outcomes. The sample includes selected full audits and excludes safety cases. Column 1 reports the number of agents recorded on the case. Column 2 uses the taxpayer-survey duration measure. Column 3 uses the administrative measure of days from the start of the audit to confirmation and includes the original date availability controls. Column 4 uses inspectors' self-reported days worked on the case. Columns 5 to 8 use the log of evasion divided by the corresponding resource measure. All regressions include Algorithm, Inspectors x Overlap, Random, and Replacement indicators and absorb inspectorclusteryear fixed effects. Robust standard errors are shown in parentheses. Stars denote significance at $^{*}p<0.10$, $^{**}p<0.05$, $^{***}p<0.01$.
\end{tablenotes}
\end{threeparttable}
\end{adjustbox}
\end{table}

\begin{table}[H]\centering
\begin{adjustbox}{max width=1.4\textwidth}
\begin{threeparttable}
\caption{Resources Allocated to Audits and Audit Productivity, Full Audits, with Predicted Evasion Control (Panels A1 and B1 only)}
\label{tab:table5_fullaudits_yhatrf_control}

\footnotesize
\begin{tabular}{lcccc|cccc} \toprule
            &\multicolumn{4}{c}{\textbf{A1: Resource Outcomes}}                                     &\multicolumn{4}{c}{\textbf{B1: Productivity Outcomes}}                                 \\\cmidrule(lr){2-5}\cmidrule(lr){6-9}
\multicolumn{1}{l}{} & \shortstack{Number of Agents} & \shortstack{Duration in Days\\(Taxpayer Survey)} & \shortstack{Days from Start to Conf.\\(Admin Data)} & \shortstack{Days Working on Case\\(Self-Reported)} & \shortstack{Evasion/ Number of\\Agents} & \shortstack{Evasion/Duration\\(Taxpayer Survey)} & \shortstack{Evasion/Duration\\(Admin. Data)} & \shortstack{Evasion/Days Working\\(Self-Reported)} \\ \multicolumn{1}{l}{} & (1) & (2) & (3) & (4) & (5) & (6) & (7) & (8) \\ \midrule
Algorithm   &       -0.19\sym{***}&       -8.18\sym{**} &      -28.65\sym{**} &      -11.43         &        0.12         &        0.86\sym{***}&        0.16         &       -0.22         \\
            &      (0.06)         &      (3.78)         &     (14.01)         &     (25.35)         &      (0.10)         &      (0.24)         &      (0.17)         &      (0.27)         \\
\addlinespace
Inspectors x Overlap&       -0.09         &      -25.98         &      -10.64         &                     &        0.23         &        1.36         &        0.21         &                     \\
            &      (0.13)         &     (22.09)         &     (27.78)         &                     &      (0.24)         &      (1.46)         &      (0.34)         &                     \\
\addlinespace
Predicted evasion&        0.15\sym{***}&        0.40         &        1.72         &       13.35\sym{***}&        1.01\sym{***}&        1.31\sym{***}&        1.03\sym{***}&        1.04\sym{***}\\
            &      (0.01)         &      (0.59)         &      (1.96)         &      (1.84)         &      (0.04)         &      (0.10)         &      (0.06)         &      (0.11)         \\
\midrule
N           &         508         &         214         &         282         &          51         &         457         &          98         &         253         &          40         \\
R2          &        0.52         &        0.25         &        0.30         &        0.51         &        0.78         &        0.74         &        0.71         &        0.70         \\
Mean outcome&        2.87         &       28.06         &      160.07         &      177.07         &       18.19         &       15.90         &       15.00         &       13.62         \\
\bottomrule \end{tabular}


\begin{tablenotes}[flushleft]
\footnotesize
\item Notes: * 0.10 ** 0.05 *** 0.01 significance levels. OLS results of regressions of outcomes on the selection method of the case, with fixed effects at the inspector-cluster $\times$ year level. The specifications additionally control for predicted evasion. The sample includes all full audit cases selected from the 2018, 2019, and 2020 audit programs, excluding safety cases, for which we had data on the outcome variables. Panel A1 shows the results for the resource outcomes, and Panel B1 shows the results for the productivity outcomes. Column 1 shows the results for the number of agents involved in the audit, which is obtained from the audit records. Column 2 shows the results for audit duration as reported by taxpayers in the taxpayer survey. The survey asked them about their last full audit experience. Column 3 uses the difference between the date of initial notice and the date of the start of the audit (either verification announcement or information request or start date of audit) as a proxy for the duration of the audit. This measure is available only for cases that had an initial notice. We drop negative values, trim the data at the 99th percentile, and control for dummies indicating availability of date of initial notice and final notice. These controls are helpful because the definition of the start date may vary depending on the availability of information. Column 4 measures the duration of the audit using the self-reported number of days inspectors worked on a case, reported on a spreadsheet that inspectors needed to fill out during the 2019 audit program. Columns 5 to 8 use as outcomes the log of the ratio of evasion over the different cost measures used in Columns 1 to 4. We multiplied the answer by the number of inspectors that worked in the case. Robust standard errors (Huber-White formula) are shown in parentheses.
\end{tablenotes}
\end{threeparttable}
\end{adjustbox}
\end{table}

\begin{table}[H]\centering
\begin{adjustbox}{max width=1.4\textwidth}
\begin{threeparttable}
\caption{Resources Allocated to Audits and Audit Productivity, Full Audits, with Predicted Evasion Quadratic Control (Panels A1 and B1 only)}
\label{tab:table5_fullaudits_yhatrf_quadratic_control}

\footnotesize
\begin{tabular}{lcccc|cccc} \toprule
            &\multicolumn{4}{c}{\textbf{A1: Resource Outcomes}}                                     &\multicolumn{4}{c}{\textbf{B1: Productivity Outcomes}}                                 \\\cmidrule(lr){2-5}\cmidrule(lr){6-9}
\multicolumn{1}{l}{} & \shortstack{Number of Agents} & \shortstack{Duration in Days\\(Taxpayer Survey)} & \shortstack{Days from Start to Conf.\\(Admin Data)} & \shortstack{Days Working on Case\\(Self-Reported)} & \shortstack{Evasion/ Number of\\Agents} & \shortstack{Evasion/Duration\\(Taxpayer Survey)} & \shortstack{Evasion/Duration\\(Admin. Data)} & \shortstack{Evasion/Days Working\\(Self-Reported)} \\ \multicolumn{1}{l}{} & (1) & (2) & (3) & (4) & (5) & (6) & (7) & (8) \\ \midrule
Algorithm   &       -0.18\sym{***}&      -10.33\sym{**} &      -26.46\sym{*}  &      -16.02         &        0.12         &        0.85\sym{***}&        0.18         &       -0.25         \\
            &      (0.06)         &      (4.18)         &     (14.24)         &     (24.60)         &      (0.10)         &      (0.24)         &      (0.17)         &      (0.26)         \\
\addlinespace
Inspectors x Overlap&       -0.09         &      -26.27         &       -9.54         &                     &        0.25         &        1.72         &        0.23         &                     \\
            &      (0.13)         &     (21.84)         &     (27.64)         &                     &      (0.24)         &      (1.45)         &      (0.35)         &                     \\
\addlinespace
Predicted evasion&        0.08         &       10.55\sym{**} &       -8.71         &       47.60         &        0.36         &       -1.75         &       -0.28         &       -2.68         \\
            &      (0.06)         &      (5.07)         &     (12.33)         &     (31.52)         &      (0.43)         &      (1.26)         &      (0.52)         &      (2.11)         \\
\addlinespace
Predicted evasion squared&        0.00         &       -0.34\sym{*}  &        0.37         &       -1.34         &        0.02         &        0.08\sym{**} &        0.04\sym{**} &        0.10\sym{*}  \\
            &      (0.00)         &      (0.17)         &      (0.45)         &      (1.25)         &      (0.01)         &      (0.04)         &      (0.01)         &      (0.06)         \\
\midrule
N           &         508         &         214         &         282         &          51         &         457         &          98         &         253         &          40         \\
R2          &        0.52         &        0.27         &        0.30         &        0.52         &        0.78         &        0.75         &        0.72         &        0.71         \\
Mean outcome&        2.87         &       28.06         &      160.07         &      177.07         &       18.19         &       15.90         &       15.00         &       13.62         \\
\bottomrule \end{tabular}


\begin{tablenotes}[flushleft]
\footnotesize
\item Notes: * 0.10 ** 0.05 *** 0.01 significance levels. OLS results of regressions of outcomes on the selection method of the case, with fixed effects at the inspector-cluster $\times$ year level. The specifications additionally control for predicted evasion and its square. The sample includes all full audit cases selected from the 2018, 2019, and 2020 audit programs, excluding safety cases, for which we had data on the outcome variables. Panel A1 shows the results for the resource outcomes, and Panel B1 shows the results for the productivity outcomes. Column 1 shows the results for the number of agents involved in the audit, which is obtained from the audit records. Column 2 shows the results for audit duration as reported by taxpayers in the taxpayer survey. The survey asked them about their last full audit experience. Column 3 uses the difference between the date of initial notice and the date of the start of the audit (either verification announcement or information request or start date of audit) as a proxy for the duration of the audit. This measure is available only for cases that had an initial notice. We drop negative values, trim the data at the 99th percentile, and control for dummies indicating availability of date of initial notice and final notice. These controls are helpful because the definition of the start date may vary depending on the availability of information. Column 4 measures the duration of the audit using the self-reported number of days inspectors worked on a case, reported on a spreadsheet that inspectors needed to fill out during the 2019 audit program. Columns 5 to 8 use as outcomes the log of the ratio of evasion over the different cost measures used in Columns 1 to 4. We multiplied the answer by the number of inspectors that worked in the case. Robust standard errors (Huber-White formula) are shown in parentheses.
\end{tablenotes}
\end{threeparttable}
\end{adjustbox}
\end{table}

\begin{table}[H]\centering
\begin{adjustbox}{max width=1.4\textwidth}
\begin{threeparttable}
\caption{Resources Allocated to Audits and Audit Productivity, Full Audits, with Predicted Evasion Decile Controls (Panels A1 and B1 only)}
\label{tab:table5_fullaudits_yhatrf_deciles_control}

\footnotesize
\begin{tabular}{lcccc|cccc} \toprule
            &\multicolumn{4}{c}{\textbf{A1: Resource Outcomes}}                                     &\multicolumn{4}{c}{\textbf{B1: Productivity Outcomes}}                                 \\\cmidrule(lr){2-5}\cmidrule(lr){6-9}
\multicolumn{1}{l}{} & \shortstack{Number of Agents} & \shortstack{Duration in Days\\(Taxpayer Survey)} & \shortstack{Days from Start to Conf.\\(Admin Data)} & \shortstack{Days Working on Case\\(Self-Reported)} & \shortstack{Evasion/ Number of\\Agents} & \shortstack{Evasion/Duration\\(Taxpayer Survey)} & \shortstack{Evasion/Duration\\(Admin. Data)} & \shortstack{Evasion/Days Working\\(Self-Reported)} \\ \multicolumn{1}{l}{} & (1) & (2) & (3) & (4) & (5) & (6) & (7) & (8) \\ \midrule
Algorithm   &       -0.20\sym{***}&      -10.35\sym{**} &      -23.59         &      -21.35         &        0.05         &        0.84\sym{***}&        0.11         &       -0.46         \\
            &      (0.06)         &      (4.68)         &     (14.49)         &     (21.61)         &      (0.11)         &      (0.26)         &      (0.19)         &      (0.33)         \\
\addlinespace
Inspectors x Overlap&       -0.10         &      -28.01         &       -5.10         &                     &        0.21         &        0.61         &        0.23         &                     \\
            &      (0.13)         &     (20.74)         &     (27.54)         &                     &      (0.27)         &      (1.90)         &      (0.37)         &                     \\
\midrule
N           &         508         &         214         &         282         &          51         &         457         &          98         &         253         &          40         \\
R2          &        0.51         &        0.29         &        0.32         &        0.57         &        0.76         &        0.71         &        0.66         &        0.76         \\
Mean outcome&        2.87         &       28.06         &      160.07         &      177.07         &       18.19         &       15.90         &       15.00         &       13.62         \\
\bottomrule \end{tabular}


\begin{tablenotes}[flushleft]
\footnotesize
\item Notes: * 0.10 ** 0.05 *** 0.01 significance levels. OLS results of regressions of outcomes on the selection method of the case, with fixed effects at the inspector-cluster $\times$ year level. The specifications additionally control for predicted-evasion deciles formed on the pooled regression-eligible full-audit sample. The sample includes all full audit cases selected from the 2018, 2019, and 2020 audit programs, excluding safety cases, for which we had data on the outcome variables. Panel A1 shows the results for the resource outcomes, and Panel B1 shows the results for the productivity outcomes. Column 1 shows the results for the number of agents involved in the audit, which is obtained from the audit records. Column 2 shows the results for audit duration as reported by taxpayers in the taxpayer survey. The survey asked them about their last full audit experience. Column 3 uses the difference between the date of initial notice and the date of the start of the audit (either verification announcement or information request or start date of audit) as a proxy for the duration of the audit. This measure is available only for cases that had an initial notice. We drop negative values, trim the data at the 99th percentile, and control for dummies indicating availability of date of initial notice and final notice. These controls are helpful because the definition of the start date may vary depending on the availability of information. Column 4 measures the duration of the audit using the self-reported number of days inspectors worked on a case, reported on a spreadsheet that inspectors needed to fill out during the 2019 audit program. Columns 5 to 8 use as outcomes the log of the ratio of evasion over the different cost measures used in Columns 1 to 4. We multiplied the answer by the number of inspectors that worked in the case. Robust standard errors (Huber-White formula) are shown in parentheses.
\end{tablenotes}
\end{threeparttable}
\end{adjustbox}
\end{table}

\begin{table}[H]\centering
\begin{adjustbox}{max width=1.4\textwidth}
\begin{threeparttable}
\caption{Resources Allocated to Audits and Audit Productivity, Full Audits, with Predicted Evasion Quintile Controls (Panels A1 and B1 only)}
\label{tab:table5_fullaudits_yhatrf_quintiles_control}

\footnotesize
\begin{tabular}{lcccc|cccc} \toprule
            &\multicolumn{4}{c}{\textbf{A1: Resource Outcomes}}                                     &\multicolumn{4}{c}{\textbf{B1: Productivity Outcomes}}                                 \\\cmidrule(lr){2-5}\cmidrule(lr){6-9}
\multicolumn{1}{l}{} & \shortstack{Number of Agents} & \shortstack{Duration in Days\\(Taxpayer Survey)} & \shortstack{Days from Start to Conf.\\(Admin Data)} & \shortstack{Days Working on Case\\(Self-Reported)} & \shortstack{Evasion/ Number of\\Agents} & \shortstack{Evasion/Duration\\(Taxpayer Survey)} & \shortstack{Evasion/Duration\\(Admin. Data)} & \shortstack{Evasion/Days Working\\(Self-Reported)} \\ \multicolumn{1}{l}{} & (1) & (2) & (3) & (4) & (5) & (6) & (7) & (8) \\ \midrule
Algorithm   &       -0.20\sym{***}&       -8.56\sym{**} &      -27.06\sym{*}  &       -7.38         &       -0.02         &        0.76\sym{***}&       -0.02         &       -0.37         \\
            &      (0.06)         &      (4.08)         &     (14.34)         &     (21.59)         &      (0.12)         &      (0.28)         &      (0.19)         &      (0.30)         \\
\addlinespace
Inspectors x Overlap&       -0.10         &      -27.41         &       -6.37         &                     &        0.21         &        0.44         &        0.14         &                     \\
            &      (0.13)         &     (21.58)         &     (27.54)         &                     &      (0.28)         &      (1.72)         &      (0.39)         &                     \\
\midrule
N           &         508         &         214         &         282         &          51         &         457         &          98         &         253         &          40         \\
R2          &        0.50         &        0.27         &        0.32         &        0.54         &        0.72         &        0.66         &        0.62         &        0.75         \\
Mean outcome&        2.87         &       28.06         &      160.07         &      177.07         &       18.19         &       15.90         &       15.00         &       13.62         \\
\bottomrule \end{tabular}


\begin{tablenotes}[flushleft]
\footnotesize
\item Notes: * 0.10 ** 0.05 *** 0.01 significance levels. OLS results of regressions of outcomes on the selection method of the case, with fixed effects at the inspector-cluster $\times$ year level. The specifications additionally control for predicted-evasion quintiles formed on the pooled regression-eligible full-audit sample. The sample includes all full audit cases selected from the 2018, 2019, and 2020 audit programs, excluding safety cases, for which we had data on the outcome variables. Panel A1 shows the results for the resource outcomes, and Panel B1 shows the results for the productivity outcomes. Column 1 shows the results for the number of agents involved in the audit, which is obtained from the audit records. Column 2 shows the results for audit duration as reported by taxpayers in the taxpayer survey. The survey asked them about their last full audit experience. Column 3 uses the difference between the date of initial notice and the date of the start of the audit (either verification announcement or information request or start date of audit) as a proxy for the duration of the audit. This measure is available only for cases that had an initial notice. We drop negative values, trim the data at the 99th percentile, and control for dummies indicating availability of date of initial notice and final notice. These controls are helpful because the definition of the start date may vary depending on the availability of information. Column 4 measures the duration of the audit using the self-reported number of days inspectors worked on a case, reported on a spreadsheet that inspectors needed to fill out during the 2019 audit program. Columns 5 to 8 use as outcomes the log of the ratio of evasion over the different cost measures used in Columns 1 to 4. We multiplied the answer by the number of inspectors that worked in the case. Robust standard errors (Huber-White formula) are shown in parentheses.
\end{tablenotes}
\end{threeparttable}
\end{adjustbox}
\end{table}

\end{landscape}

\subsection{Table 6 with predicted evasion controls}

This subsection presents the baseline Table 6 replication together with variants that keep the original sample and layout and add predicted-evasion controls.

\begin{table}[H]\centering
\begin{adjustbox}{max width=1.4\textwidth}
\begin{threeparttable}
\caption{Algorithm Selection, Perceived Audit Efficiency and Corruption}
\label{tab:table6_indices_replicated}

\footnotesize
\begin{tabular}{lccc|ccc} \toprule
\multicolumn{7}{l}{\textbf{A: Surveyed Firms (Self-Reporting a Recent Audit)}} \\ \multicolumn{1}{l}{Outcome:} & \multicolumn{3}{c}{Efficiency Index} & \multicolumn{3}{c}{Corruption Index} \\ \multicolumn{1}{l}{} & Full Audits & Desk Audits & All Audits & Full Audits & Desk Audits & All Audits \\ \multicolumn{1}{l}{} & (1) & (2) & (3) & (4) & (5) & (6) \\ \midrule
Algorithm   &       -0.41\sym{***}&        0.01         &       -0.17         &        0.05         &       -0.05         &       -0.02         \\
            &      (0.15)         &      (0.14)         &      (0.10)         &      (0.15)         &      (0.14)         &      (0.10)         \\
\addlinespace
Inspectors x Overlap&        0.45\sym{**} &        0.04         &        0.16         &       -0.56\sym{*}  &       -0.12         &       -0.18         \\
            &      (0.18)         &      (0.26)         &      (0.21)         &      (0.29)         &      (0.27)         &      (0.20)         \\
\addlinespace
Algorithm x Random&                     &       -0.09         &       -0.05         &                     &       -0.08         &       -0.04         \\
            &                     &      (0.19)         &      (0.18)         &                     &      (0.18)         &      (0.16)         \\
\midrule
N           &         199         &         273         &         472         &         199         &         272         &         471         \\
R2          &        0.10         &        0.02         &        0.03         &        0.11         &        0.10         &        0.05         \\
Mean outcome&        0.00         &        0.00         &        0.00         &       -0.07         &        0.05         &        0.00         \\

\midrule \multicolumn{7}{l}{\textbf{B: Only Audited Firms (as per Administrative Audit Data)}} \\ \multicolumn{1}{l}{} & (1) & (2) & (3) & (4) & (5) & (6) \\ \midrule

Algorithm   &       -0.21         &        0.36         &        0.02         &        0.19         &       -0.28         &       -0.03         \\
            &      (0.26)         &      (0.23)         &      (0.16)         &      (0.25)         &      (0.23)         &      (0.16)         \\
\addlinespace
Inspectors x Overlap&        0.36         &        0.33         &       -0.01         &       -0.64         &       -0.33         &       -0.03         \\
            &      (0.46)         &      (0.39)         &      (0.30)         &      (0.70)         &      (0.37)         &      (0.32)         \\
\addlinespace
Algorithm x Random&                     &       -0.04         &        0.08         &                     &       -0.06         &       -0.06         \\
            &                     &      (0.26)         &      (0.25)         &                     &      (0.28)         &      (0.24)         \\
\midrule
N           &          89         &         120         &         211         &          88         &         122         &         212         \\
R2          &        0.25         &        0.15         &        0.11         &        0.25         &        0.17         &        0.09         \\
Mean outcome&        0.00         &       -0.06         &       -0.02         &       -0.14         &       -0.03         &       -0.08         \\
\bottomrule \end{tabular}


\begin{tablenotes}[flushleft]
\footnotesize
\item Notes: * 0.10 ** 0.05 *** 0.01 significance levels. OLS results with fixed effects at the year and tax office level for outcomes from the taxpayer survey. Given the lower number of observations in the survey, we do not run the desk audit regression with year $\times$ inspector fixed effects, as is done with outcomes from the administrative data. Robust standard errors are in parentheses. The table shows the association between audit selection method and taxpayers' perception of the efficiency of audits (Columns 1 to 3), and of the incidence of corruption during audits (Columns 4 to 6). The sample of surveyed taxpayers is a subsample of the audit cases selected by the yearly audit programs between 2018 and 2020. Panel A shows the results for all firms interviewed in the survey, conditional on them self-reporting an audit in recent years. Panel B is restricted to firms for which we can confirm that an audit as part of the yearly audit program took place, according to the administrative audit reports. The two outcomes are (1) an index of perceived efficiency of audits, and (2) an index of perceived corruption of audits. Each measure is constructed by combining three questions into a normalized index following Anderson (2008). To measure efficiency, we combine the questions on taxpayers' assessment of the auditor's (i) technical knowledge, (ii) efficiency during the audit, and (iii) capacity to uncover all evaded taxes. For corruption, we combine the questions on the perception of (i) the auditor's dishonesty during the audit, (ii) the frequency of bribes paid to inspectors among firms similar to the respondent's, and (iii) the existence of a preferential tax audit treatment for firms connected with the tax administration. The coefficients are measured in terms of standard deviation of the respective indexes.
\end{tablenotes}
\end{threeparttable}
\end{adjustbox}
\end{table}

\begin{table}[H]\centering
\begin{adjustbox}{max width=1.4\textwidth}
\begin{threeparttable}
\caption{Algorithm Selection, Perceived Audit Efficiency and Corruption, with Predicted Evasion Control}
\label{tab:table6_indices_yhatrf_control}

\footnotesize
\begin{tabular}{lccc|ccc} \toprule
\multicolumn{7}{l}{\textbf{A: Surveyed Firms (Self-Reporting a Recent Audit)}} \\ \multicolumn{1}{l}{Outcome:} & \multicolumn{3}{c}{Efficiency Index} & \multicolumn{3}{c}{Corruption Index} \\ \multicolumn{1}{l}{} & Full Audits & Desk Audits & All Audits & Full Audits & Desk Audits & All Audits \\ \multicolumn{1}{l}{} & (1) & (2) & (3) & (4) & (5) & (6) \\ \midrule
Algorithm   &       -0.40\sym{**} &       -0.01         &       -0.18\sym{*}  &       -0.00         &       -0.06         &       -0.04         \\
            &      (0.16)         &      (0.14)         &      (0.10)         &      (0.15)         &      (0.13)         &      (0.10)         \\
\addlinespace
Inspectors x Overlap&        0.43\sym{**} &        0.05         &        0.17         &       -0.46         &       -0.12         &       -0.17         \\
            &      (0.18)         &      (0.26)         &      (0.21)         &      (0.30)         &      (0.27)         &      (0.20)         \\
\addlinespace
Algorithm x Random&                     &       -0.09         &       -0.06         &                     &       -0.08         &       -0.05         \\
            &                     &      (0.19)         &      (0.18)         &                     &      (0.18)         &      (0.16)         \\
\addlinespace
Predicted evasion&        0.02         &       -0.02         &       -0.01         &       -0.06         &       -0.01         &       -0.02         \\
            &      (0.04)         &      (0.02)         &      (0.01)         &      (0.03)         &      (0.02)         &      (0.01)         \\
\midrule
N           &         199         &         273         &         472         &         199         &         272         &         471         \\
R2          &        0.10         &        0.03         &        0.03         &        0.12         &        0.10         &        0.06         \\
Mean outcome&        0.00         &        0.00         &        0.00         &       -0.07         &        0.05         &        0.00         \\

\midrule \multicolumn{7}{l}{\textbf{B: Only Audited Firms (as per Administrative Audit Data)}} \\ \multicolumn{1}{l}{} & (1) & (2) & (3) & (4) & (5) & (6) \\ \midrule

Algorithm   &       -0.36         &        0.31         &       -0.05         &        0.27         &       -0.24         &        0.00         \\
            &      (0.27)         &      (0.23)         &      (0.17)         &      (0.25)         &      (0.22)         &      (0.16)         \\
\addlinespace
Inspectors x Overlap&        0.18         &        0.40         &       -0.06         &       -0.27         &       -0.13         &        0.10         \\
            &      (0.24)         &      (0.36)         &      (0.30)         &      (0.49)         &      (0.31)         &      (0.28)         \\
\addlinespace
Algorithm x Random&                     &       -0.00         &        0.07         &                     &       -0.03         &       -0.04         \\
            &                     &      (0.27)         &      (0.24)         &                     &      (0.25)         &      (0.23)         \\
\addlinespace
Predicted evasion&        0.01         &       -0.03         &       -0.02         &       -0.04         &        0.00         &       -0.01         \\
            &      (0.06)         &      (0.02)         &      (0.02)         &      (0.06)         &      (0.02)         &      (0.02)         \\
\midrule
N           &          89         &         120         &         211         &          88         &         122         &         212         \\
R2          &        0.10         &        0.10         &        0.04         &        0.16         &        0.13         &        0.05         \\
Mean outcome&        0.00         &       -0.06         &       -0.02         &       -0.14         &       -0.03         &       -0.08         \\
\bottomrule \end{tabular}


\begin{tablenotes}[flushleft]
\footnotesize
\item Notes: * 0.10 ** 0.05 *** 0.01 significance levels. OLS results with fixed effects at the year and tax office level for outcomes from the taxpayer survey. The specifications additionally control for predicted evasion. Given the lower number of observations in the survey, we do not run the desk audit regression with year $\times$ inspector fixed effects, as is done with outcomes from the administrative data. Robust standard errors are in parentheses. The table shows the association between audit selection method and taxpayers' perception of the efficiency of audits (Columns 1 to 3), and of the incidence of corruption during audits (Columns 4 to 6). The sample of surveyed taxpayers is a subsample of the audit cases selected by the yearly audit programs between 2018 and 2020. Panel A shows the results for all firms interviewed in the survey, conditional on them self-reporting an audit in recent years. Panel B is restricted to firms for which we can confirm that an audit as part of the yearly audit program took place, according to the administrative audit reports. The two outcomes are (1) an index of perceived efficiency of audits, and (2) an index of perceived corruption of audits. Each measure is constructed by combining three questions into a normalized index following Anderson (2008). To measure efficiency, we combine the questions on taxpayers' assessment of the auditor's (i) technical knowledge, (ii) efficiency during the audit, and (iii) capacity to uncover all evaded taxes. For corruption, we combine the questions on the perception of (i) the auditor's dishonesty during the audit, (ii) the frequency of bribes paid to inspectors among firms similar to the respondent's, and (iii) the existence of a preferential tax audit treatment for firms connected with the tax administration. The coefficients are measured in terms of standard deviation of the respective indexes.
\end{tablenotes}
\end{threeparttable}
\end{adjustbox}
\end{table}

\begin{table}[H]\centering
\begin{adjustbox}{max width=\textwidth}
\begin{threeparttable}
\caption{Algorithm Selection, Perceived Audit Efficiency and Corruption, with Predicted Evasion Quadratic Control}
\label{tab:table6_indices_yhatrf_quadratic_control}

\footnotesize
\begin{tabular}{lccc|ccc} \toprule
\multicolumn{7}{l}{\textbf{A: Surveyed Firms (Self-Reporting a Recent Audit)}} \\ \multicolumn{1}{l}{Outcome:} & \multicolumn{3}{c}{Efficiency Index} & \multicolumn{3}{c}{Corruption Index} \\ \multicolumn{1}{l}{} & Full Audits & Desk Audits & All Audits & Full Audits & Desk Audits & All Audits \\ \multicolumn{1}{l}{} & (1) & (2) & (3) & (4) & (5) & (6) \\ \midrule
Algorithm   &       -0.39\sym{**} &       -0.01         &       -0.17\sym{*}  &       -0.02         &       -0.06         &       -0.06         \\
            &      (0.16)         &      (0.14)         &      (0.11)         &      (0.15)         &      (0.13)         &      (0.10)         \\
\addlinespace
Inspectors x Overlap&        0.42\sym{**} &        0.05         &        0.18         &       -0.45         &       -0.13         &       -0.19         \\
            &      (0.18)         &      (0.27)         &      (0.21)         &      (0.31)         &      (0.27)         &      (0.20)         \\
\addlinespace
Algorithm x Random&                     &       -0.09         &       -0.06         &                     &       -0.09         &       -0.06         \\
            &                     &      (0.19)         &      (0.18)         &                     &      (0.18)         &      (0.16)         \\
\addlinespace
Predicted evasion&       -0.05         &       -0.00         &       -0.04         &        0.05         &        0.04         &        0.07         \\
            &      (0.30)         &      (0.07)         &      (0.05)         &      (0.23)         &      (0.07)         &      (0.06)         \\
\addlinespace
Predicted evasion squared&        0.00         &       -0.00         &        0.00         &       -0.00         &       -0.00         &       -0.00\sym{*}  \\
            &      (0.01)         &      (0.00)         &      (0.00)         &      (0.01)         &      (0.00)         &      (0.00)         \\
\midrule
N           &         199         &         273         &         472         &         199         &         272         &         471         \\
R2          &        0.10         &        0.03         &        0.03         &        0.12         &        0.10         &        0.06         \\
Mean outcome&        0.00         &        0.00         &        0.00         &       -0.07         &        0.05         &        0.00         \\

\midrule \multicolumn{7}{l}{\textbf{B: Only Audited Firms (as per Administrative Audit Data)}} \\ \multicolumn{1}{l}{} & (1) & (2) & (3) & (4) & (5) & (6) \\ \midrule

Algorithm   &       -0.39         &        0.31         &       -0.03         &        0.33         &       -0.23         &       -0.00         \\
            &      (0.27)         &      (0.23)         &      (0.17)         &      (0.25)         &      (0.22)         &      (0.17)         \\
\addlinespace
Inspectors x Overlap&        0.17         &        0.40         &       -0.03         &       -0.25         &       -0.12         &        0.09         \\
            &      (0.24)         &      (0.35)         &      (0.29)         &      (0.49)         &      (0.31)         &      (0.28)         \\
\addlinespace
Algorithm x Random&                     &        0.02         &        0.11         &                     &        0.00         &       -0.06         \\
            &                     &      (0.26)         &      (0.24)         &                     &      (0.26)         &      (0.24)         \\
\addlinespace
Predicted evasion&        0.14         &       -0.08         &       -0.11         &       -0.34         &       -0.08         &        0.02         \\
            &      (0.37)         &      (0.13)         &      (0.08)         &      (0.32)         &      (0.12)         &      (0.08)         \\
\addlinespace
Predicted evasion squared&       -0.00         &        0.00         &        0.00         &        0.01         &        0.00         &       -0.00         \\
            &      (0.01)         &      (0.01)         &      (0.00)         &      (0.01)         &      (0.01)         &      (0.00)         \\
\midrule
N           &          89         &         120         &         211         &          88         &         122         &         212         \\
R2          &        0.11         &        0.10         &        0.04         &        0.18         &        0.13         &        0.06         \\
Mean outcome&        0.00         &       -0.06         &       -0.02         &       -0.14         &       -0.03         &       -0.08         \\
\bottomrule \end{tabular}


\begin{tablenotes}[flushleft]
\footnotesize
\item Notes: * 0.10 ** 0.05 *** 0.01 significance levels. OLS results with fixed effects at the year and tax office level for outcomes from the taxpayer survey. The specifications additionally control for predicted evasion and its square. Given the lower number of observations in the survey, we do not run the desk audit regression with year $\times$ inspector fixed effects, as is done with outcomes from the administrative data. Robust standard errors are in parentheses. The table shows the association between audit selection method and taxpayers' perception of the efficiency of audits (Columns 1 to 3), and of the incidence of corruption during audits (Columns 4 to 6). The sample of surveyed taxpayers is a subsample of the audit cases selected by the yearly audit programs between 2018 and 2020. Panel A shows the results for all firms interviewed in the survey, conditional on them self-reporting an audit in recent years. Panel B is restricted to firms for which we can confirm that an audit as part of the yearly audit program took place, according to the administrative audit reports. The two outcomes are (1) an index of perceived efficiency of audits, and (2) an index of perceived corruption of audits. Each measure is constructed by combining three questions into a normalized index following Anderson (2008). To measure efficiency, we combine the questions on taxpayers' assessment of the auditor's (i) technical knowledge, (ii) efficiency during the audit, and (iii) capacity to uncover all evaded taxes. For corruption, we combine the questions on the perception of (i) the auditor's dishonesty during the audit, (ii) the frequency of bribes paid to inspectors among firms similar to the respondent's, and (iii) the existence of a preferential tax audit treatment for firms connected with the tax administration. The coefficients are measured in terms of standard deviation of the respective indexes.
\end{tablenotes}
\end{threeparttable}
\end{adjustbox}
\end{table}

\begin{table}[H]\centering
\begin{adjustbox}{max width=\textwidth}
\begin{threeparttable}
\caption{Algorithm Selection, Perceived Audit Efficiency and Corruption, with Predicted Evasion Decile Controls}
\label{tab:table6_indices_yhatrf_deciles_control}

\footnotesize
\begin{tabular}{lccc|ccc} \toprule
\multicolumn{7}{l}{\textbf{A: Surveyed Firms (Self-Reporting a Recent Audit)}} \\ \multicolumn{1}{l}{Outcome:} & \multicolumn{3}{c}{Efficiency Index} & \multicolumn{3}{c}{Corruption Index} \\ \multicolumn{1}{l}{} & Full Audits & Desk Audits & All Audits & Full Audits & Desk Audits & All Audits \\ \multicolumn{1}{l}{} & (1) & (2) & (3) & (4) & (5) & (6) \\ \midrule
Algorithm   &       -0.33\sym{*}  &       -0.05         &       -0.17         &       -0.03         &       -0.05         &       -0.02         \\
            &      (0.17)         &      (0.14)         &      (0.11)         &      (0.15)         &      (0.13)         &      (0.10)         \\
\addlinespace
Inspectors x Overlap&        0.47\sym{**} &        0.03         &        0.16         &       -0.44         &       -0.10         &       -0.14         \\
            &      (0.21)         &      (0.28)         &      (0.21)         &      (0.33)         &      (0.28)         &      (0.21)         \\
\addlinespace
Algorithm x Random&                     &       -0.08         &       -0.04         &                     &       -0.14         &       -0.07         \\
            &                     &      (0.19)         &      (0.18)         &                     &      (0.18)         &      (0.16)         \\
\midrule
N           &         199         &         273         &         472         &         199         &         272         &         471         \\
R2          &        0.13         &        0.07         &        0.04         &        0.14         &        0.14         &        0.07         \\
Mean outcome&        0.00         &        0.00         &        0.00         &       -0.07         &        0.05         &        0.00         \\

\midrule \multicolumn{7}{l}{\textbf{B: Only Audited Firms (as per Administrative Audit Data)}} \\ \multicolumn{1}{l}{} & (1) & (2) & (3) & (4) & (5) & (6) \\ \midrule

Algorithm   &       -0.30         &        0.31         &       -0.04         &        0.21         &       -0.24         &        0.06         \\
            &      (0.30)         &      (0.24)         &      (0.17)         &      (0.26)         &      (0.22)         &      (0.16)         \\
\addlinespace
Inspectors x Overlap&        0.21         &        0.41         &        0.02         &       -0.29         &       -0.14         &        0.08         \\
            &      (0.27)         &      (0.37)         &      (0.30)         &      (0.55)         &      (0.37)         &      (0.31)         \\
\addlinespace
Algorithm x Random&                     &       -0.13         &        0.03         &                     &        0.01         &        0.03         \\
            &                     &      (0.28)         &      (0.25)         &                     &      (0.28)         &      (0.24)         \\
\midrule
N           &          89         &         120         &         211         &          88         &         122         &         212         \\
R2          &        0.15         &        0.18         &        0.08         &        0.24         &        0.21         &        0.10         \\
Mean outcome&        0.00         &       -0.06         &       -0.02         &       -0.14         &       -0.03         &       -0.08         \\
\bottomrule \end{tabular}


\begin{tablenotes}[flushleft]
\footnotesize
\item Notes: * 0.10 ** 0.05 *** 0.01 significance levels. OLS results with fixed effects at the year and tax office level for outcomes from the taxpayer survey. The specifications additionally control for predicted-evasion deciles formed on the pooled regression-eligible survey-analysis sample. Given the lower number of observations in the survey, we do not run the desk audit regression with year $\times$ inspector fixed effects, as is done with outcomes from the administrative data. Robust standard errors are in parentheses. The table shows the association between audit selection method and taxpayers' perception of the efficiency of audits (Columns 1 to 3), and of the incidence of corruption during audits (Columns 4 to 6). The sample of surveyed taxpayers is a subsample of the audit cases selected by the yearly audit programs between 2018 and 2020. Panel A shows the results for all firms interviewed in the survey, conditional on them self-reporting an audit in recent years. Panel B is restricted to firms for which we can confirm that an audit as part of the yearly audit program took place, according to the administrative audit reports. The two outcomes are (1) an index of perceived efficiency of audits, and (2) an index of perceived corruption of audits. Each measure is constructed by combining three questions into a normalized index following Anderson (2008). To measure efficiency, we combine the questions on taxpayers' assessment of the auditor's (i) technical knowledge, (ii) efficiency during the audit, and (iii) capacity to uncover all evaded taxes. For corruption, we combine the questions on the perception of (i) the auditor's dishonesty during the audit, (ii) the frequency of bribes paid to inspectors among firms similar to the respondent's, and (iii) the existence of a preferential tax audit treatment for firms connected with the tax administration. The coefficients are measured in terms of standard deviation of the respective indexes.
\end{tablenotes}
\end{threeparttable}
\end{adjustbox}
\end{table}

\begin{table}[H]\centering
\begin{adjustbox}{max width=\textwidth}
\begin{threeparttable}
\caption{Algorithm Selection, Perceived Audit Efficiency and Corruption, with Predicted Evasion Quintile Controls}
\label{tab:table6_indices_yhatrf_quintiles_control}

\footnotesize
\begin{tabular}{lccc|ccc} \toprule
\multicolumn{7}{l}{\textbf{A: Surveyed Firms (Self-Reporting a Recent Audit)}} \\ \multicolumn{1}{l}{Outcome:} & \multicolumn{3}{c}{Efficiency Index} & \multicolumn{3}{c}{Corruption Index} \\ \multicolumn{1}{l}{} & Full Audits & Desk Audits & All Audits & Full Audits & Desk Audits & All Audits \\ \multicolumn{1}{l}{} & (1) & (2) & (3) & (4) & (5) & (6) \\ \midrule
Algorithm   &       -0.41\sym{***}&        0.01         &       -0.16         &        0.02         &       -0.10         &       -0.05         \\
            &      (0.16)         &      (0.14)         &      (0.10)         &      (0.15)         &      (0.14)         &      (0.10)         \\
\addlinespace
Inspectors x Overlap&        0.43\sym{**} &        0.04         &        0.17         &       -0.42         &       -0.11         &       -0.19         \\
            &      (0.20)         &      (0.26)         &      (0.20)         &      (0.31)         &      (0.26)         &      (0.18)         \\
\addlinespace
Algorithm x Random&                     &       -0.08         &       -0.05         &                     &       -0.07         &       -0.07         \\
            &                     &      (0.20)         &      (0.18)         &                     &      (0.18)         &      (0.17)         \\
\midrule
N           &         199         &         273         &         472         &         199         &         272         &         471         \\
R2          &        0.10         &        0.03         &        0.03         &        0.12         &        0.11         &        0.07         \\
Mean outcome&        0.00         &        0.00         &        0.00         &       -0.07         &        0.05         &        0.00         \\

\midrule \multicolumn{7}{l}{\textbf{B: Only Audited Firms (as per Administrative Audit Data)}} \\ \multicolumn{1}{l}{} & (1) & (2) & (3) & (4) & (5) & (6) \\ \midrule

Algorithm   &       -0.41         &        0.32         &       -0.04         &        0.28         &       -0.32         &       -0.00         \\
            &      (0.26)         &      (0.24)         &      (0.17)         &      (0.25)         &      (0.23)         &      (0.16)         \\
\addlinespace
Inspectors x Overlap&        0.21         &        0.45         &        0.03         &       -0.31         &       -0.07         &        0.05         \\
            &      (0.25)         &      (0.33)         &      (0.29)         &      (0.50)         &      (0.40)         &      (0.29)         \\
\addlinespace
Algorithm x Random&                     &        0.04         &        0.08         &                     &       -0.02         &       -0.06         \\
            &                     &      (0.28)         &      (0.25)         &                     &      (0.27)         &      (0.24)         \\
\midrule
N           &          89         &         120         &         211         &          88         &         122         &         212         \\
R2          &        0.12         &        0.11         &        0.04         &        0.21         &        0.14         &        0.07         \\
Mean outcome&        0.00         &       -0.06         &       -0.02         &       -0.14         &       -0.03         &       -0.08         \\
\bottomrule \end{tabular}


\begin{tablenotes}[flushleft]
\footnotesize
\item Notes: * 0.10 ** 0.05 *** 0.01 significance levels. OLS results with fixed effects at the year and tax office level for outcomes from the taxpayer survey. The specifications additionally control for predicted-evasion quintiles formed on the pooled regression-eligible survey-analysis sample. Given the lower number of observations in the survey, we do not run the desk audit regression with year $\times$ inspector fixed effects, as is done with outcomes from the administrative data. Robust standard errors are in parentheses. The table shows the association between audit selection method and taxpayers' perception of the efficiency of audits (Columns 1 to 3), and of the incidence of corruption during audits (Columns 4 to 6). The sample of surveyed taxpayers is a subsample of the audit cases selected by the yearly audit programs between 2018 and 2020. Panel A shows the results for all firms interviewed in the survey, conditional on them self-reporting an audit in recent years. Panel B is restricted to firms for which we can confirm that an audit as part of the yearly audit program took place, according to the administrative audit reports. The two outcomes are (1) an index of perceived efficiency of audits, and (2) an index of perceived corruption of audits. Each measure is constructed by combining three questions into a normalized index following Anderson (2008). To measure efficiency, we combine the questions on taxpayers' assessment of the auditor's (i) technical knowledge, (ii) efficiency during the audit, and (iii) capacity to uncover all evaded taxes. For corruption, we combine the questions on the perception of (i) the auditor's dishonesty during the audit, (ii) the frequency of bribes paid to inspectors among firms similar to the respondent's, and (iii) the existence of a preferential tax audit treatment for firms connected with the tax administration. The coefficients are measured in terms of standard deviation of the respective indexes.
\end{tablenotes}
\end{threeparttable}
\end{adjustbox}
\end{table}

\subsection{Table 8 with predicted evasion controls}

This subsection presents the full-audit Table 8 baseline replication together with variants that keep the original observable-characteristics ladder and add predicted-evasion controls in a sixth column.

\newcommand{\TableEightRobustnessNotes}[1]{%
\footnotesize
\item Notes: * 0.10 ** 0.05 *** 0.01 significance levels. OLS results of a regression of the main audit outcomes on the selection method of the case for full audits. We show specifications for full audits controlling for tax office x year fixed effects, which corresponds to the selection list for full audits. Each column progressively adds control variables at the firm level: Column 1 has no controls, Column 2 controls for turnover (log), Column 3 controls in addition for the profit rate (profit over turnover), Column 4 controls in addition for the firm productivity (turnover divided by the wage bill), and Column 5 controls in addition for further firm characteristics. #1 The first outcome (A) is the probability that the case on the list was conducted by the inspectors. The second outcome (B) is the probability that an audit results in a penalty (initial or final notice), conditional on the case being conducted. The third outcome (C) is the log amount of the evasion plus penalties, conditional on the case ending in positive penalty. The evasion amount is obtained from the final notice and complemented with the initial notice when the final notice is not available. The sample includes all full audt cases selected in the 2018, 2019, and 2020 audit programs. The data includes the selection and audits of the Large Taxpayer Unit, Medium Taxpayer Units 1 and 2, Liberal Professionals, and the regional SME units of Dakar Plateau, Grand Dakar, Pikine Guediawaye, and Ngor Almadies. Robust standard errors (Huber-White) are shown in parentheses, and inspector x year levels when inspector fixed effects are included.
}

\begin{table}[H]\centering
\begin{adjustbox}{max width=\textwidth}
\begin{threeparttable}
\caption{Explaining the Execution Gap for Full Audits: The Role of Observable Characteristics}
\label{tab:table8_fullaudits_replicated}

\footnotesize
\input{Output/table8_fullaudits_replicated.tex}

\begin{tablenotes}[flushleft]
\TableEightRobustnessNotes{}
\end{tablenotes}
\end{threeparttable}
\end{adjustbox}
\end{table}

\begin{table}[H]\centering
\begin{adjustbox}{max width=\textwidth}
\begin{threeparttable}
\caption{Explaining the Execution Gap for Full Audits: The Role of Observable Characteristics, with Predicted Evasion Control}
\label{tab:table8_fullaudits_yhatrf_control}

\footnotesize
\input{Output/table8_fullaudits_yhatrf_control.tex}

\begin{tablenotes}[flushleft]
\TableEightRobustnessNotes{Column 6 additionally controls linearly for predicted evasion.}
\end{tablenotes}
\end{threeparttable}
\end{adjustbox}
\end{table}

\begin{table}[H]\centering
\begin{adjustbox}{max width=\textwidth}
\begin{threeparttable}
\caption{Explaining the Execution Gap for Full Audits: The Role of Observable Characteristics, with Predicted Evasion Quadratic Control}
\label{tab:table8_fullaudits_yhatrf_quadratic_control}

\footnotesize
\input{Output/table8_fullaudits_yhatrf_quadratic_control.tex}

\begin{tablenotes}[flushleft]
\TableEightRobustnessNotes{Column 6 additionally controls for predicted evasion and its square.}
\end{tablenotes}
\end{threeparttable}
\end{adjustbox}
\end{table}

\begin{table}[H]\centering
\begin{adjustbox}{max width=\textwidth}
\begin{threeparttable}
\caption{Explaining the Execution Gap for Full Audits: The Role of Observable Characteristics, with Predicted Evasion Decile Controls}
\label{tab:table8_fullaudits_yhatrf_deciles_control}

\footnotesize
\input{Output/table8_fullaudits_yhatrf_deciles_control.tex}

\begin{tablenotes}[flushleft]
\TableEightRobustnessNotes{Column 6 additionally controls for pooled-sample predicted-evasion deciles.}
\end{tablenotes}
\end{threeparttable}
\end{adjustbox}
\end{table}

\begin{table}[H]\centering
\begin{adjustbox}{max width=\textwidth}
\begin{threeparttable}
\caption{Explaining the Execution Gap for Full Audits: The Role of Observable Characteristics, with Predicted Evasion Quintile Controls}
\label{tab:table8_fullaudits_yhatrf_quintiles_control}

\footnotesize
\input{Output/table8_fullaudits_yhatrf_quintiles_control.tex}

\begin{tablenotes}[flushleft]
\TableEightRobustnessNotes{Column 6 additionally controls for pooled-sample predicted-evasion quintiles.}
\end{tablenotes}
\end{threeparttable}
\end{adjustbox}
\end{table}
\section{Inspector effort analysis}


   \begin{comment}
\subsection{Date definition [Pre-Post]}

\begin{table}[h]
    \caption{Number of executed desk audit cases by date statistics avaialability}
    {
\def\sym#1{\ifmmode^{#1}\else\(^{#1}\)\fi}
\begin{tabular}{l*{4}{c}}
\toprule
                    &   Algorithm&  Inspectors&      Random&       Total\\
\midrule
Total Executed Cases&         364&         476&         171&        1011\\
Initial and final date Avail.&          97&         159&          53&         309\\
Valid Duration stat available &          92&         137&          53&         282\\
Has either Conf. or Notif. date&         184&         224&          71&         479\\
\bottomrule
\end{tabular}
}

    \label{tab:date_stats}
\end{table}

\begin{itemize}
    \item{Date availability stats are better for full audits (57\% has both opening and closing dates vrs 30\% in desk audits.}
\end{itemize}

\begin{figure}
    \caption{Selection year and opening date per of audit cases}
    \centering
    \includegraphics[scale=0.65]{Output/selection_year_opening_dates.pdf}
    \label{fig:placeholder}
\end{figure}
\begin{itemize}
    \item{There are some cases that started before the selection year.} 
\end{itemize}
\begin{table}[H]
\centering
    \caption{Number of executed  cases by audit type and date availability}
    \begin{tabular}{l*{3}{c}}
\toprule
            & Desk audits& Full audits&       Total\\
\midrule
Has only opening date&         222&          10&         232\\
            &            &            &            \\
\addlinespace
Has only closing date&         479&         206&         685\\
            &            &            &            \\
\addlinespace
Has both    &         309&         292&         601\\
            &            &            &            \\
\addlinespace
Has neither &           1&           0&           1\\
            &            &            &            \\
\addlinespace
Total       &        1011&         508&        1519\\
            &            &            &            \\
\bottomrule
\end{tabular}

    \label{tab:date_stats}
\end{table}

\end{comment}

\subsection{Definition of void audit cases}
\begin{table}[h]
    \centering
    \caption{Number of void desk audit cases}
    {
\def\sym#1{\ifmmode^{#1}\else\(^{#1}\)\fi}
\begin{tabular}{l*{4}{c}}
\toprule
Void Audit  &   Algorithm&  Inspectors&      Random&       Total\\
\midrule
No          &         264&         367&         116&         747\\
Yes         &         100&         109&          55&         264\\
Total       &         364&         476&         171&        1011\\
\bottomrule
\end{tabular}
}

    \label{tab:void_aud_stats}    
\end{table}



\begin{table}[H]
    \centering
    \caption{Number of null cases and their notification/confirmation status}
    \begin{adjustbox}{max width=\textwidth}
    {
\def\sym#1{\ifmmode^{#1}\else\(^{#1}\)\fi}
\begin{tabular}{l*{6}{c}}
\toprule
            &\multicolumn{3}{c}{Desk audits}                                  &\multicolumn{3}{c}{Full audits}                                  \\
            &Evasion Detected         &No evasion detected         &       Total         &Evasion Detected         &No evasion detected         &       Total         \\
\midrule
Only Notified&         286         &           4         &         290         &         151         &          18         &         169         \\
            &                     &                     &                     &                     &                     &                     \\
\addlinespace
Only Confirmed&          53         &           4         &          57         &          39         &           0         &          39         \\
            &                     &                     &                     &                     &                     &                     \\
\addlinespace
Notified and Confirmed&         408         &           2         &         410         &         267         &           2         &         269         \\
            &                     &                     &                     &                     &                     &                     \\
\addlinespace
Neither     &           0         &         254         &         254         &           0         &          31         &          31         \\
            &                     &                     &                     &                     &                     &                     \\
\addlinespace
Total       &         747         &         264         &        1011         &         457         &          51         &         508         \\
            &                     &                     &                     &                     &                     &                     \\
\midrule
\(N\)       &        1011         &                     &                     &         508         &                     &                     \\
\bottomrule
\end{tabular}
}

    \label{tab:nullsandconfirmation}
    \end{adjustbox}
\end{table}

\begin{table}[H]
    \centering
    \caption{Number of notified desk audits according to  penalty/droit simple registered.}
    \begin{adjustbox}{max width=\textwidth}
    \begin{tabular}{l*{7}{c}}
\toprule
            &\multicolumn{4}{c}{Desk audits}                    &\multicolumn{3}{c}{Full audits}       \\
            &   Algorithm&  Inspectors&      Random&       Total&   Algorithm&  Inspectors&       Total\\
\midrule
has penalités only&           4&           0&           0&           4&           1&           0&           1\\
            &            &            &            &            &            &            &            \\
\addlinespace
has both    &         239&         339&         112&         690&         149&         268&         417\\
            &            &            &            &            &            &            &            \\
\addlinespace
has neither &           3&           2&           1&           6&           6&          14&          20\\
            &            &            &            &            &            &            &            \\
\addlinespace
Total       &         246&         341&         113&         700&         156&         282&         438\\
            &            &            &            &            &            &            &            \\
\bottomrule
\end{tabular}

    \label{tab:nullsandconfirmation}
    \end{adjustbox}
\end{table}


From what I understand of the paper, the workflow of inspectors is as follows:
\begin{itemize}
\item{Irregularities detected  $\rightarrow$ evasion value  $\rightarrow$ notification to the taxpayer} 
\item{Taxpayer may dispute or not $\rightarrow$ confirmation}
\end{itemize}

However, after doing some  tabs I noticed that: 


\begin{itemize}
\item{Some executed cases with positive evasion value skipped the notification step and went directly to confirmation, which shouldn't be an issue if using y3==0 for void audit.}
\item{Some cases with zero evasion were either notified or confirmed, which is weird right? If there was no evasion detected (evasionvalue==0 or y4==0) no notification or confirmation should've taken place.}
\end{itemize}
\begin{itemize}
    \item{Note: not all inspectors that were assigned desk audits executed an algorithm case.}
\end{itemize}

\subsection{Definitions of low evasion rates cases}
\textbf{steps}
\begin{enumerate}
    \item{Estimate evasion rate of each case as $\frac{\text{Evasion Value}}{\text{Evasion value}+ \text{Mean Turnover (2018-2020)}}$ } 
    \item{
    Estimate within inspector cluster year (list) - method [Alg + Rand / Inspectors] quartiles according to evasion rate.}
    \item{Tag cases in bottom 25\% (bad) or below median  evasion rate}
    \item{Estimate the shares of cases on the bottom or below median evasion rate per inspector,  selection method and year.}
    \item{Run scatterplots per inspector and given method, comparing share of low evasion rate cases in year t and execution rate of that method in year t +1.}
\end{enumerate}
\begin{figure}
    \caption{Distribution of evasion rates across years}
    \centering
    \includegraphics[scale=0.65]{Output/kdensity_evasionrate_syear.pdf}
    \label{fig:kdensity_evasionrate}
\end{figure}
\begin{comment}
\begin{figure}[H]
    \centering
    \caption{Distribution of void audit cases by inspector and selection method}
    \includegraphics[scale=0.65]{Output/kdensity_void_audit_rate.pdf}
    \label{fig:void_audit_plot}
\end{figure}


\begin{table}[h]
    \centering
    \caption{Inspector's share of void desk audit cases, by method}
    {
\def\sym#1{\ifmmode^{#1}\else\(^{#1}\)\fi}
\begin{tabular}{l*{1}{ccccc}}
\toprule
            &         min&        mean&         max&          sd&       count\\
\midrule
Algorithm   &           0&       9.690&         100&       16.20&          81\\
Inspectors  &           0&       10.36&         100&       17.28&          87\\
Random      &           0&       9.654&         100&       20.35&          65\\
\bottomrule
\multicolumn{6}{l}{\footnotesize Total number of Inspectors: 123}\\
\end{tabular}
}

    \label{tab:stats_insp_vaudits}    
\end{table}



\begin{table}[h]
    \centering
    \caption{Inspectors audit stats, by method}
    \begin{tabular}{l*{4}{c}} \hline
            &   Algorithm&  Inspectors&      Random&       Total\\
\hline \multicolumn{5}{l}{\textbf{Panel A: No. assigned cases}} \\
mean        &    9.145299&    11.18103&    3.956522&    22.20325\\
sd          &    6.962046&    7.656312&    3.182759&    16.94783\\
count       &         117&         116&          92&         123\\
\hline
            &   Algorithm&  Inspectors&      Random&       Total\\
\hline \multicolumn{5}{l}{\textbf{Panel B: No. executed cases}} \\
mean        &    3.111111&    4.103448&    1.858696&    8.219512\\
sd          &    4.280626&    5.388199&    2.406865&    10.78671\\
count       &         117&         116&          92&         123\\
\hline \end{tabular}

    \label{tab:stats_insp_vaudits}    
\end{table}


    

\begin{table}[h]
    \centering
    \caption{Inspectors executed and void audit shares, by method}
    \begin{tabular}{l*{4}{c}} \hline
            &   Algorithm&  Inspectors&      Random\\
\hline \multicolumn{4}{l}{\textbf{Panel A: Share Execution  over implementation}} \\
mean        &    35.17204&    47.50068&    27.08887\\
sd          &    27.04335&    27.38804&    30.29245\\
count       &         101&         102&          85\\
\hline
            &   Algorithm&  Inspectors&      Random\\
\hline \multicolumn{4}{l}{\textbf{Panel B: Share void audit over implementation}} \\
mean        &    9.690092&     10.3646&    9.654304\\
sd          &     16.2033&    17.27603&    20.34805\\
count       &          81&          87&          65\\
\hline \end{tabular}

    \label{tab:stats_insp_vaudits}    
\end{table}
\end{comment}



\section{Approaches to assess inspectors effort}
\subsection{Approach 1: pre-post void audit case  [discarded for now]}
\textbf{Issues encountered:}
\begin{itemize}
    \item{Few desk audit inspectors have a valid open - closure audit date for all of their executed cases: 41 out of 123.}
    \item{Only 34 out 123  have experienced a void audit case}
    \item{Only 2 inspectors have both: i) valid dates for all audits + ii) encountered a void audit}
\end{itemize}
\textbf{Possible solutions to missing dates?}
\begin{itemize}
    \item{Impute starting date values $+$ 1 month}
    \item{use average duration? subsampling by void versus non-void?}
\end{itemize}

\subsection{Approach 2: Comparing previous year performance with next year's execution rate }
In this subsection we compare the inspectors' performance between two periods based on selection year. We proceed to plot per selection method, i) the share of executed cases from  year $t$ vrs $t+1$, and ii) the share of void audits in $t$ versus the share of execution of assigned audits in $t+1$.  

\begin{itemize}
    \item{Sample: Only desk - audit inspectors that are observable for at least two periods. 
    Furthermore, we divide the sample  between inspectors that changed tax office in the next period and those that remained in the same tax office.}
\end{itemize}

\textbf{Definition of shares}

Let Alg $=$ algorithm  + random cases.
Let $N$ denote the number of cases per inspector $i$ at (selection) year $t$. Then, 

\textbf{Share of void cases:}

\begin{itemize}
    \item{$\text{Share of all void cases}_{it} = \frac{N^\text{Total Void}_{it}}{N^{\text{Total Executed}}_{it}}$}

    \item{$\text{Share of void alg cases}_{it} = \frac{N^\text{Void Alg}_{it}}{N^{\text{Total Alg Executed}}_{it}}$}
    
     \item{$\text{Share of void Inspector cases}_{it} = \frac{N^\text{Void Inspector}_{it}}{N^{\text{Total Inspector Executed}}_{it}}$}
\end{itemize}

\textbf{Share of executed cases}
\begin{itemize}
    \item{$\text{Share of all executed cases}_{it} = \frac{N^\text{Total Executed}_{it}}{N^{\text{Total Assigned}}_{it}}$}

     \item{$\text{Share of executed Alg cases}_{it} = \frac{N^\text{Executed Alg}_{it}}{N^{\text{Total Alg Assigned}}_{it}}$}

     \item{$\text{Share of executed Inspector cases}_{it} = \frac{N^\text{Executed Inspector}_{it}}{N^{\text{Total Inspector Assigned}}_{it}}$}
\end{itemize}

\textbf{}

\begin{table}[H]
\caption{Number of inspectors by period and subsample (Desk audits}
    \centering
    \resizebox{\textwidth}{!}{%
{
\def\sym#1{\ifmmode^{#1}\else\(^{#1}\)\fi}
\begin{tabular}{l*{9}{c}}
\toprule
            &\multicolumn{3}{c}{2018–2019}                                    &\multicolumn{3}{c}{2019–2020}                                    &\multicolumn{3}{c}{2018 – avg(2019–2020)}                        \\\cmidrule(lr){2-4}\cmidrule(lr){5-7}\cmidrule(lr){8-10}
            &\multicolumn{1}{c}{(1)}&\multicolumn{1}{c}{(2)}&\multicolumn{1}{c}{(3)}&\multicolumn{1}{c}{(4)}&\multicolumn{1}{c}{(5)}&\multicolumn{1}{c}{(6)}&\multicolumn{1}{c}{(7)}&\multicolumn{1}{c}{(8)}&\multicolumn{1}{c}{(9)}\\
            &\multicolumn{1}{c}{Same bureau}&\multicolumn{1}{c}{Changed bureau}&\multicolumn{1}{c}{Total}&\multicolumn{1}{c}{Same bureau}&\multicolumn{1}{c}{Changed bureau}&\multicolumn{1}{c}{Total}&\multicolumn{1}{c}{Same bureau}&\multicolumn{1}{c}{Changed bureau}&\multicolumn{1}{c}{Total}\\
\midrule
Total Inspectors&          81         &           1         &          82         &          86         &           2         &          88         &          96         &           6         &         102         \\
            &                     &                     &                     &                     &                     &                     &                     &                     &                     \\
Inspectors observable two periods&          13         &           1         &          14         &          23         &           2         &          25         &          51         &           6         &          57         \\
            &                     &                     &                     &                     &                     &                     &                     &                     &                     \\
\bottomrule
\end{tabular}
}

    }
\end{table}


\begin{table}[H]
\caption{Number of inspectors that registered execution and void audits by period and sub-sample (Desk audits)}
    \centering
    \resizebox{\textwidth}{!}{%
        {
\def\sym#1{\ifmmode^{#1}\else\(^{#1}\)\fi}
\begin{tabular}{l*{9}{c}}
\toprule
            &\multicolumn{3}{c}{2018–2019}                                    &\multicolumn{3}{c}{2019–2020}                                    &\multicolumn{3}{c}{2018 – avg(2019–2020)}                        \\\cmidrule(lr){2-4}\cmidrule(lr){5-7}\cmidrule(lr){8-10}
            &\multicolumn{1}{c}{(1)}&\multicolumn{1}{c}{(2)}&\multicolumn{1}{c}{(3)}&\multicolumn{1}{c}{(4)}&\multicolumn{1}{c}{(5)}&\multicolumn{1}{c}{(6)}&\multicolumn{1}{c}{(7)}&\multicolumn{1}{c}{(8)}&\multicolumn{1}{c}{(9)}\\
            &\multicolumn{1}{c}{Same bureau}&\multicolumn{1}{c}{Changed bureau}&\multicolumn{1}{c}{Total}&\multicolumn{1}{c}{Same bureau}&\multicolumn{1}{c}{Changed bureau}&\multicolumn{1}{c}{Total}&\multicolumn{1}{c}{Same bureau}&\multicolumn{1}{c}{Changed bureau}&\multicolumn{1}{c}{Total}\\
\midrule
Year t      &                     &                     &                     &                     &                     &                     &                     &                     &                     \\
Executed any audit&          13         &           1         &          14         &          23         &           2         &          25         &          30         &           5         &          35         \\
            &                     &                     &                     &                     &                     &                     &                     &                     &                     \\
Executed algo-selected audit&          13         &           1         &          14         &          22         &           2         &          24         &          28         &           5         &          33         \\
            &                     &                     &                     &                     &                     &                     &                     &                     &                     \\
Executed inspector-selected audit&          11         &           1         &          12         &          23         &           2         &          25         &          24         &           5         &          29         \\
            &                     &                     &                     &                     &                     &                     &                     &                     &                     \\
Reported any void audit&           3         &           0         &           3         &          21         &           2         &          23         &           6         &           2         &           8         \\
            &                     &                     &                     &                     &                     &                     &                     &                     &                     \\
Reported void algo audit&           0         &           0         &           0         &          18         &           2         &          20         &           3         &           1         &           4         \\
            &                     &                     &                     &                     &                     &                     &                     &                     &                     \\
Reported void inspector audit&           3         &           0         &           3         &          17         &           1         &          18         &           3         &           1         &           4         \\
            &                     &                     &                     &                     &                     &                     &                     &                     &                     \\
Reported any below-median evasion rate&          13         &           1         &          14         &          21         &           2         &          23         &          27         &           5         &          32         \\
            &                     &                     &                     &                     &                     &                     &                     &                     &                     \\
Reported below-median evasion rate algo audit&          11         &           1         &          12         &          20         &           2         &          22         &          21         &           5         &          26         \\
            &                     &                     &                     &                     &                     &                     &                     &                     &                     \\
Reported below-median evasion rate  inspector audit&          10         &           1         &          11         &          18         &           2         &          20         &          18         &           5         &          23         \\
            &                     &                     &                     &                     &                     &                     &                     &                     &                     \\
Reported any bot-quartile evasion rate&          10         &           1         &          11         &          19         &           2         &          21         &          22         &           5         &          27         \\
            &                     &                     &                     &                     &                     &                     &                     &                     &                     \\
Reported bot-quartile  evasion rate algo audit&           9         &           1         &          10         &          17         &           2         &          19         &          17         &           4         &          21         \\
            &                     &                     &                     &                     &                     &                     &                     &                     &                     \\
Reported bot-quartile  evasion rate  inspector audit&           7         &           1         &           8         &          16         &           1         &          17         &          12         &           5         &          17         \\
            &                     &                     &                     &                     &                     &                     &                     &                     &                     \\
Year t+1    &                     &                     &                     &                     &                     &                     &                     &                     &                     \\
Executed any audit&          13         &           1         &          14         &          21         &           2         &          23         &          45         &           5         &          50         \\
            &                     &                     &                     &                     &                     &                     &                     &                     &                     \\
Executed algo-selected audit&          12         &           1         &          13         &          17         &           0         &          17         &          32         &           3         &          35         \\
            &                     &                     &                     &                     &                     &                     &                     &                     &                     \\
Executed inspector-selected audit&          13         &           1         &          14         &          17         &           2         &          19         &          45         &           5         &          50         \\
            &                     &                     &                     &                     &                     &                     &                     &                     &                     \\
Reported any void audit&          10         &           0         &          10         &           1         &           0         &           1         &          26         &           2         &          28         \\
            &                     &                     &                     &                     &                     &                     &                     &                     &                     \\
Reported void algo audit&           9         &           0         &           9         &           1         &           0         &           1         &          22         &           2         &          24         \\
            &                     &                     &                     &                     &                     &                     &                     &                     &                     \\
Reported void inspector audit&           8         &           0         &           8         &           0         &           0         &           0         &          21         &           1         &          22         \\
            &                     &                     &                     &                     &                     &                     &                     &                     &                     \\
Reported any below-median evasion rate&          12         &           1         &          13         &          17         &           1         &          18         &          41         &           4         &          45         \\
            &                     &                     &                     &                     &                     &                     &                     &                     &                     \\
Reported below-median evasion rate algo audit&          12         &           0         &          12         &          13         &           0         &          13         &          30         &           2         &          32         \\
            &                     &                     &                     &                     &                     &                     &                     &                     &                     \\
Reported below-median evasion rate  inspector audit&          11         &           1         &          12         &           9         &           1         &          10         &          36         &           4         &          40         \\
            &                     &                     &                     &                     &                     &                     &                     &                     &                     \\
Reported any bot-quartile evasion rate&          12         &           0         &          12         &          12         &           0         &          12         &          38         &           3         &          41         \\
            &                     &                     &                     &                     &                     &                     &                     &                     &                     \\
Reported bot-quartile  evasion rate algo audit&          12         &           0         &          12         &           7         &           0         &           7         &          28         &           2         &          30         \\
            &                     &                     &                     &                     &                     &                     &                     &                     &                     \\
Reported bot-quartile  evasion rate  inspector audit&          10         &           0         &          10         &           7         &           0         &           7         &          31         &           2         &          33         \\
            &                     &                     &                     &                     &                     &                     &                     &                     &                     \\
\bottomrule
\end{tabular}
}

    }
\end{table}

\begin{table}[H]
\caption{Inspector average execution and void case rates by period and sub-sample (Desk audits)}
    \centering
    \resizebox{\textwidth}{!}{%
        {
\def\sym#1{\ifmmode^{#1}\else\(^{#1}\)\fi}
\begin{tabular}{l*{9}{c}}
\toprule
            &\multicolumn{3}{c}{2018–2019}                                    &\multicolumn{3}{c}{2019–2020}                                    &\multicolumn{3}{c}{2018 – avg(2019–2020)}                        \\\cmidrule(lr){2-4}\cmidrule(lr){5-7}\cmidrule(lr){8-10}
            &\multicolumn{1}{c}{(1)}&\multicolumn{1}{c}{(2)}&\multicolumn{1}{c}{(3)}&\multicolumn{1}{c}{(4)}&\multicolumn{1}{c}{(5)}&\multicolumn{1}{c}{(6)}&\multicolumn{1}{c}{(7)}&\multicolumn{1}{c}{(8)}&\multicolumn{1}{c}{(9)}\\
            &\multicolumn{1}{c}{Same bureau}&\multicolumn{1}{c}{Changed bureau}&\multicolumn{1}{c}{Total}&\multicolumn{1}{c}{Same bureau}&\multicolumn{1}{c}{Changed bureau}&\multicolumn{1}{c}{Total}&\multicolumn{1}{c}{Same bureau}&\multicolumn{1}{c}{Changed bureau}&\multicolumn{1}{c}{Total}\\
\midrule
Year t      &                     &                     &                     &                     &                     &                     &                     &                     &                     \\
Exec rate, all audits&       43.08         &       53.33         &       43.81         &       61.59         &       43.52         &       60.14         &       36.50         &          58         &       39.33         \\
            &     (21.01)         &         (.)         &     (20.37)         &     (26.51)         &     (9.166)         &     (25.94)         &     (21.58)         &     (11.45)         &     (21.71)         \\
Exec rate, algo audits&       44.12         &       61.11         &       45.33         &       60.82         &       50.00         &       59.96         &       37.83         &       54.44         &       40.02         \\
            &     (25.20)         &         (.)         &     (24.64)         &     (28.25)         &     (23.57)         &     (27.63)         &     (27.28)         &     (14.91)         &     (26.46)         \\
Exec rate, inspector audits&       41.82         &       41.67         &       41.81         &       62.53         &       37.50         &       60.53         &       34.15         &       63.33         &       37.99         \\
            &     (28.58)         &         (.)         &     (27.45)         &     (28.64)         &     (5.893)         &     (28.31)         &     (28.19)         &     (17.28)         &     (28.63)         \\
Void rate, all audits&       2.585         &           0         &       2.401         &       41.14         &       28.33         &       40.11         &       4.731         &       3.152         &       4.506         \\
            &     (4.940)         &         (.)         &     (4.796)         &     (25.24)         &     (16.50)         &     (24.65)         &     (11.66)         &     (4.400)         &     (10.89)         \\
Void rate, algo audits&           0         &           0         &           0         &       41.57         &       22.50         &       39.98         &       6.548         &       3.333         &       6.061         \\
            &         (0)         &         (.)         &         (0)         &     (28.61)         &     (3.536)         &     (27.87)         &     (21.44)         &     (7.454)         &     (19.90)         \\
Void rate, inspector audits&       8.636         &           0         &       7.917         &       38.55         &          30         &       37.87         &       3.958         &           4         &       3.966         \\
            &     (16.45)         &         (.)         &     (15.88)         &     (29.55)         &     (42.43)         &     (29.68)         &     (11.70)         &     (8.944)         &     (11.13)         \\
Share of below median ev. rate, all audits&       48.17         &       43.75         &       47.86         &       44.06         &       53.33         &       44.80         &       49.85         &       47.37         &       49.49         \\
            &     (22.64)         &         (.)         &     (21.78)         &     (23.03)         &     (18.86)         &     (22.53)         &     (28.53)         &     (12.09)         &     (26.69)         \\
Share of below median ev. rate, algo audits&       43.99         &       45.45         &       44.10         &       48.44         &       47.50         &       48.36         &       44.41         &       50.76         &       45.37         \\
            &     (27.95)         &         (.)         &     (26.86)         &     (22.87)         &     (38.89)         &     (23.31)         &     (33.37)         &     (24.03)         &     (31.89)         \\
Share of below median ev. rate, inspector audits&       47.47         &          40         &       46.85         &       41.34         &          55         &       42.43         &       51.27         &       44.56         &       50.11         \\
            &     (26.75)         &         (.)         &     (25.60)         &     (27.68)         &     (7.071)         &     (26.81)         &     (38.64)         &     (21.13)         &     (36.01)         \\
Share of bottom quartile ev. rate, all audits&       23.80         &          25         &       23.88         &       35.51         &          45         &       36.27         &       29.92         &       23.98         &       29.07         \\
            &     (16.26)         &         (.)         &     (15.63)         &     (23.59)         &     (7.071)         &     (22.78)         &     (27.36)         &     (10.91)         &     (25.63)         \\
Share of bottom quartile ev. rate, algo audits&       22.74         &       27.27         &       23.06         &       40.34         &       47.50         &       40.94         &       23.59         &       22.12         &       23.37         \\
            &     (17.58)         &         (.)         &     (16.94)         &     (27.91)         &     (38.89)         &     (27.95)         &     (24.96)         &     (18.39)         &     (23.84)         \\
Share of bottom quartile ev. rate, inspector audits&       24.48         &          20         &       24.11         &       31.09         &          30         &       31.01         &       32.40         &       25.89         &       31.28         \\
            &     (30.23)         &         (.)         &     (28.85)         &     (23.45)         &     (42.43)         &     (24.07)         &     (42.23)         &     (11.27)         &     (38.59)         \\
Year t+1    &                     &                     &                     &                     &                     &                     &                     &                     &                     \\
Exec rate, all audits&       54.84         &       41.67         &       53.90         &       19.34         &       8.333         &       18.46         &       32.07         &       22.66         &       31.08         \\
            &     (19.38)         &         (.)         &     (18.95)         &     (11.93)         &         (0)         &     (11.82)         &     (23.09)         &     (15.59)         &     (22.50)         \\
Exec rate, algo audits&       53.21         &          50         &       52.98         &       19.39         &           0         &       17.84         &       29.04         &       16.67         &       27.73         \\
            &     (25.71)         &         (.)         &     (24.72)         &     (18.45)         &         (0)         &     (18.46)         &     (27.29)         &     (21.08)         &     (26.82)         \\
Exec rate, inspector audits&       57.68         &       33.33         &       55.94         &       19.29         &       11.11         &       18.63         &       33.36         &       23.15         &       32.29         \\
            &     (21.02)         &         (.)         &     (21.22)         &     (18.81)         &         (0)         &     (18.15)         &     (24.46)         &     (15.61)         &     (23.79)         \\
Void rate, all audits&       37.14         &           0         &       34.48         &       1.587         &           0         &       1.449         &       19.11         &       5.667         &       17.77         \\
            &     (29.19)         &         (.)         &     (29.75)         &     (7.274)         &         (0)         &     (6.950)         &     (25.20)         &     (8.788)         &     (24.35)         \\
Void rate, algo audits&       39.35         &           0         &       36.32         &       1.961         &           .         &       1.961         &       26.52         &          15         &       25.54         \\
            &     (29.74)         &         (.)         &     (30.49)         &     (8.085)         &         (.)         &     (8.085)         &     (27.27)         &     (13.23)         &     (26.44)         \\
Void rate, inspector audits&       35.37         &           0         &       32.84         &           0         &           0         &           0         &       17.89         &           6         &       16.70         \\
            &     (37.39)         &         (.)         &     (37.14)         &         (0)         &         (0)         &         (0)         &     (25.12)         &     (13.42)         &     (24.38)         \\
Share of below median ev. rate, all audits&       49.55         &          20         &       47.44         &       42.62         &          50         &       43.26         &       52.12         &       39.67         &       50.88         \\
            &     (21.22)         &         (.)         &     (21.86)         &     (26.69)         &     (70.71)         &     (29.65)         &     (28.71)         &     (37.09)         &     (29.44)         \\
Share of below median ev. rate, algo audits&       57.41         &           0         &       52.99         &       56.08         &           .         &       56.08         &       55.37         &       31.67         &       53.34         \\
            &     (15.53)         &         (.)         &     (21.78)         &     (39.41)         &         (.)         &     (39.41)         &     (25.68)         &     (38.84)         &     (27.11)         \\
Share of below median ev. rate, inspector audits&       44.27         &          50         &       44.68         &       33.82         &          50         &       35.53         &       48.87         &          46         &       48.58         \\
            &     (25.84)         &         (.)         &     (24.87)         &     (36.38)         &     (70.71)         &     (38.47)         &     (34.15)         &     (31.90)         &     (33.63)         \\
Share of bottom quartile ev. rate, all audits&       38.64         &           0         &       35.88         &       21.27         &           0         &       19.42         &       37.75         &          24         &       36.37         \\
            &     (19.27)         &         (.)         &     (21.20)         &     (21.17)         &         (0)         &     (21.10)         &     (28.12)         &     (30.70)         &     (28.36)         \\
Share of bottom quartile ev. rate, algo audits&       49.77         &           0         &       45.94         &       27.84         &           .         &       27.84         &       41.73         &       31.67         &       40.87         \\
            &     (18.91)         &         (.)         &     (22.77)         &     (39.16)         &         (.)         &     (39.16)         &     (28.76)         &     (38.84)         &     (29.17)         \\
Share of bottom quartile ev. rate, inspector audits&       29.35         &           0         &       27.26         &       16.18         &           0         &       14.47         &       32.66         &          21         &       31.49         \\
            &     (23.21)         &         (.)         &     (23.64)         &     (21.54)         &         (0)         &     (20.94)         &     (30.14)         &     (32.86)         &     (30.27)         \\
\bottomrule
\end{tabular}
}

    }
\end{table}


\begin{table}[H]
\caption{Number of tax offices  by period and subsample (Full - Audits)} 
    \centering
    
    \resizebox{\textwidth}{!}{%
\input{Output/bureau_sample_stats.tex}
    }
\end{table}

\begin{table}[H]
\caption{Number of Tax offices  that registered execution and void audits by period and sub-sample (Full - Audits)}
    \centering
    \resizebox{\textwidth}{!}{%
        \small {
\def\sym#1{\ifmmode^{#1}\else\(^{#1}\)\fi}
\begin{tabular}{l*{3}{c}}
\toprule
            &\multicolumn{1}{c}{(1)}&\multicolumn{1}{c}{(2)}&\multicolumn{1}{c}{(3)}\\
            &\multicolumn{1}{c}{2018–2019}&\multicolumn{1}{c}{2019–2020}&\multicolumn{1}{c}{2018 – avg(2019–2020)}\\
\midrule
Year t      &                     &                     &                     \\
Executed any audit&           7         &          11         &           7         \\
            &                     &                     &                     \\
Executed algo-selected audit&           7         &          11         &           7         \\
            &                     &                     &                     \\
Executed inspector-selected audit&           7         &          11         &           7         \\
            &                     &                     &                     \\
Reported any void audit&           6         &           9         &           6         \\
            &                     &                     &                     \\
Reported void algo audit&           6         &           4         &           6         \\
            &                     &                     &                     \\
Reported void inspector audit&           4         &           9         &           4         \\
            &                     &                     &                     \\
Reported any below-median evasion rate&           7         &          11         &           7         \\
            &                     &                     &                     \\
Reported below-median evasion rate algo audit&           7         &           7         &           7         \\
            &                     &                     &                     \\
Reported below-median evasion rate  inspector audit&           7         &          11         &           7         \\
            &                     &                     &                     \\
Reported any bot-quartile evasion rate&           7         &          11         &           7         \\
            &                     &                     &                     \\
Reported bot-quartile  evasion rate algo audit&           7         &           7         &           7         \\
            &                     &                     &                     \\
Reported bot-quartile  evasion rate  inspector audit&           7         &          10         &           7         \\
            &                     &                     &                     \\
Year t+1    &                     &                     &                     \\
Executed any audit&           7         &          10         &          11         \\
            &                     &                     &                     \\
Executed algo-selected audit&           7         &           6         &          11         \\
            &                     &                     &                     \\
Executed inspector-selected audit&           7         &          10         &          11         \\
            &                     &                     &                     \\
Reported any void audit&           6         &           2         &           9         \\
            &                     &                     &                     \\
Reported void algo audit&           4         &           0         &           4         \\
            &                     &                     &                     \\
Reported void inspector audit&           6         &           2         &           9         \\
            &                     &                     &                     \\
Reported any below-median evasion rate&           7         &           8         &          11         \\
            &                     &                     &                     \\
Reported below-median evasion rate algo audit&           6         &           4         &           8         \\
            &                     &                     &                     \\
Reported below-median evasion rate  inspector audit&           7         &           7         &          11         \\
            &                     &                     &                     \\
Reported any bot-quartile evasion rate&           7         &           7         &          11         \\
            &                     &                     &                     \\
Reported bot-quartile  evasion rate algo audit&           6         &           3         &           7         \\
            &                     &                     &                     \\
Reported bot-quartile  evasion rate  inspector audit&           6         &           7         &          10         \\
            &                     &                     &                     \\
\bottomrule
\end{tabular}
}

    }
\end{table}



\begin{table}[H]
\caption{Tax office average execution and void case rates by period and sub-sample (Full - Audits)}
    \centering
    \resizebox{\textwidth}{!}{%
        {
\def\sym#1{\ifmmode^{#1}\else\(^{#1}\)\fi}
\begin{tabular}{l*{3}{c}}
\toprule
            &\multicolumn{1}{c}{(1)}&\multicolumn{1}{c}{(2)}&\multicolumn{1}{c}{(3)}\\
            &\multicolumn{1}{c}{2018–2019}&\multicolumn{1}{c}{2019–2020}&\multicolumn{1}{c}{2018 – avg(2019–2020)}\\
\midrule
Year t      &                     &                     &                     \\
Exec rate, all audits&       64.72         &       56.90         &       64.72         \\
            &     (15.74)         &     (18.92)         &     (15.74)         \\
Exec rate, algo audits&       52.84         &       50.10         &       52.84         \\
            &     (22.50)         &     (26.63)         &     (22.50)         \\
Exec rate, inspector audits&       75.67         &       63.46         &       75.67         \\
            &     (15.50)         &     (18.65)         &     (15.50)         \\
Void rate, all audits&       11.42         &       16.86         &       11.42         \\
            &     (6.523)         &     (14.18)         &     (6.523)         \\
Void rate, algo audits&       15.84         &       11.30         &       15.84         \\
            &     (12.40)         &     (20.82)         &     (12.40)         \\
Void rate, inspector audits&       9.325         &       18.93         &       9.325         \\
            &     (9.347)         &     (16.64)         &     (9.347)         \\
Share of below median ev. rate, all audits&       48.20         &       43.65         &       48.20         \\
            &     (8.663)         &     (21.50)         &     (8.663)         \\
Share of below median ev. rate, algo audits&       50.58         &       46.04         &       50.58         \\
            &     (10.96)         &     (42.06)         &     (10.96)         \\
Share of below median ev. rate, inspector audits&       47.39         &       43.21         &       47.39         \\
            &     (14.02)         &     (17.15)         &     (14.02)         \\
Share of bottom quartile ev. rate, all audits&       24.64         &       25.74         &       24.64         \\
            &     (10.84)         &     (15.07)         &     (10.84)         \\
Share of bottom quartile ev. rate, algo audits&       26.47         &       29.11         &       26.47         \\
            &     (12.44)         &     (31.33)         &     (12.44)         \\
Share of bottom quartile ev. rate, inspector audits&       24.16         &       23.89         &       24.16         \\
            &     (14.78)         &     (14.41)         &     (14.78)         \\
Year t+1    &                     &                     &                     \\
Exec rate, all audits&       67.30         &       28.43         &       42.66         \\
            &     (9.983)         &     (23.68)         &     (15.00)         \\
Exec rate, algo audits&       62.06         &       21.86         &       35.98         \\
            &     (19.85)         &     (23.85)         &     (16.30)         \\
Exec rate, inspector audits&       71.67         &       34.43         &       48.94         \\
            &     (15.44)         &     (28.66)         &     (18.79)         \\
Void rate, all audits&       15.54         &       3.958         &       11.74         \\
            &     (13.95)         &     (10.51)         &     (10.19)         \\
Void rate, algo audits&       17.76         &           0         &       10.32         \\
            &     (24.26)         &         (0)         &     (20.83)         \\
Void rate, inspector audits&       13.31         &       4.386         &       12.11         \\
            &     (11.13)         &     (10.70)         &     (8.835)         \\
Share of below median ev. rate, all audits&       52.40         &       39.24         &       41.93         \\
            &     (16.66)         &     (22.74)         &     (18.23)         \\
Share of below median ev. rate, algo audits&       58.06         &       43.91         &       48.76         \\
            &     (35.14)         &     (40.09)         &     (37.01)         \\
Share of below median ev. rate, inspector audits&       48.61         &       39.49         &       41.50         \\
            &     (17.37)         &     (32.64)         &     (16.81)         \\
Share of bottom quartile ev. rate, all audits&       26.05         &       20.19         &       23.19         \\
            &     (9.876)         &     (17.15)         &     (11.46)         \\
Share of bottom quartile ev. rate, algo audits&       31.46         &       15.17         &       26.70         \\
            &     (19.18)         &     (23.26)         &     (24.44)         \\
Share of bottom quartile ev. rate, inspector audits&       21.12         &       24.71         &       23.18         \\
            &     (14.77)         &     (29.64)         &     (17.33)         \\
\bottomrule
\end{tabular}
}

    }
\end{table}





% =========================================================
% HELPER MACRO (put in your preamble, or right before figures)
% =========================================================
% Tune this: try 0.28-0.33 depending on how tall your plots are
\newcommand{\panelH}{0.60\textheight}

\newcommand{\twopanels}[4]{%
\begin{subfigure}{0.49\textwidth}
  \centering
  \caption{#1 (Desk Audits)}
  \includegraphics[width=\linewidth,height=\panelH,keepaspectratio]{#2}
\end{subfigure}\hfill
\begin{subfigure}{0.49\textwidth}
  \centering
  \caption{#1 (Full Audits)}
  \includegraphics[width=\linewidth,height=\panelH,keepaspectratio]{#3}
\end{subfigure}
\vspace{#4}
}
% =========================================================
% FIGURES (baseline vs _full for each row)
% =========================================================

%%% Share of execution of all audits t vrs t+1
\begin{figure}[H]
    \centering
    \caption{Share of execution of all $\text{audits}_t$ vrs Share of execution of all $\text{audits}_{t+1}$}

    \twopanels{2018 vrs 2019}
      {Output/scatter_share_exec_all_2018_2019.pdf}
      {Output/scatter_share_exec_all_2018_2019_full.pdf}
      {0.6em}

    \twopanels{2019 vrs 2020}
      {Output/scatter_share_exec_all_2019_2020.pdf}
      {Output/scatter_share_exec_all_2019_2020_full.pdf}
      {0.6em}

    \twopanels{2018 vrs (2019--2020)}
      {Output/scatter_share_exec_all_2018_avg.pdf}
      {Output/scatter_share_exec_all_2018_avg_full.pdf}
      {0em}
\end{figure}

% share of execution algo audits t vrs  t+1
\begin{figure}[H]
    \centering
    \caption{Share of execution of $\text{Algorithm}_t$ vrs Share of execution of $\text{Algorithm}_{t+1}$}

    \twopanels{2018 vrs 2019}
      {Output/scatter_share_exec_ALG_2018_2019.pdf}
      {Output/scatter_share_exec_Algorithm_2018_2019_full.pdf}
      {0.35em}

    \twopanels{2019 vrs 2020}
      {Output/scatter_share_exec_ALG_2019_2020.pdf}
      {Output/scatter_share_exec_Algorithm_2019_2020_full.pdf}
      {0.35em}

    \twopanels{2018 vrs (2019--2020)}
      {Output/scatter_share_exec_ALG_2018_avg.pdf}
      {Output/scatter_share_exec_Algorithm_2018_avg_full.pdf}
      {0em}
\end{figure}

% share of execution Inspector audits t vrs  t+1
\begin{figure}[H]
    \centering
    \caption{Share of execution of $\text{Inspector}_t$ vrs Share of execution of $\text{Inspector}_{t+1}$}

    \twopanels{2018 vrs 2019}
      {Output/scatter_share_exec_Inspectors_2018_2019.pdf}
      {Output/scatter_share_exec_Inspectors_2018_2019_full.pdf}
      {0.6em}

    \twopanels{2019 vrs 2020}
      {Output/scatter_share_exec_Inspectors_2019_2020.pdf}
      {Output/scatter_share_exec_Inspectors_2019_2020_full.pdf}
      {0.6em}

    \twopanels{2018 vrs (2019--2020)}
      {Output/scatter_share_exec_Inspectors_2018_avg.pdf}
      {Output/scatter_share_exec_Inspectors_2018_avg_full.pdf}
      {0em}
\end{figure}


% share of all void execution vrs total executed/assigned t+1
\begin{figure}[H]
    \centering
    \caption{Share of total $\text{void}_t$ vrs Share of total execution of $\text{assigned}_{t+1}$}

    \twopanels{2018 vrs 2019}
      {Output/scatter_share_v_all_2018_2019.pdf}
      {Output/scatter_share_v_all_2018_2019_full.pdf}
      {0.6em}

    \twopanels{2019 vrs 2020}
      {Output/scatter_share_v_all_2019_2020.pdf}
      {Output/scatter_share_v_all_2019_2020_full.pdf}
      {0.6em}

    \twopanels{2018 vrs (2019--2020)}
      {Output/scatter_share_v_all_2018_avg.pdf}
      {Output/scatter_share_v_all_2018_avg_full.pdf}
      {0em}
\end{figure}

% share of algo void execution vrs share of execution t+1
\begin{figure}[H]
    \centering
    \caption{Share of void $\text{algo}_t$ vrs share of execution of $\text{algo}_{t+1}$}

    \twopanels{2018 vrs 2019}
      {Output/scatter_share_v_ALG_2018_2019.pdf}
      {Output/scatter_share_v_Algorithm_2018_2019_full.pdf}
      {0.35em}

    \twopanels{2019 vrs 2020}
      {Output/scatter_share_v_ALG_2019_2020.pdf}
      {Output/scatter_share_v_Algorithm_2019_2020_full.pdf}
      {0.35em}

    \twopanels{2018 vrs (2019--2020)}
      {Output/scatter_share_v_ALG_2018_avg.pdf}
      {Output/scatter_share_v_Algorithm_2018_avg_full.pdf}
      {0em}
\end{figure}

% share of Inspectors void execution vrs share of inspector selected executions  t+1
\begin{figure}[H]
    \centering
    \caption{Share of void $\text{Inspector-selected}_t$ vrs share of execution of $\text{Inspector-selected}_{t+1}$}

    \twopanels{2018 vrs 2019}
      {Output/scatter_share_v_Inspectors_2018_2019.pdf}
      {Output/scatter_share_v_Inspectors_2018_2019_full.pdf}
      {0.6em}

    \twopanels{2019 vrs 2020}
      {Output/scatter_share_v_Inspectors_2019_2020.pdf}
      {Output/scatter_share_v_Inspectors_2019_2020_full.pdf}
      {0.6em}

    \twopanels{2018 vrs (2019--2020)}
      {Output/scatter_share_v_Inspectors_2018_avg.pdf}
      {Output/scatter_share_v_Inspectors_2018_avg_full.pdf}
      {0em}
\end{figure}

%%%%% Share of below median evasion rate in t vrs share of execution on t+1

% share of all below median evasion rate cases  vrs total executed/assigned t+1
\begin{figure}[H]
    \centering
    \caption{Share of total $\text{below median evasion rate cases}_t$ vrs Share of total execution of $\text{assigned}_{t+1}$}

    \twopanels{2018 vrs 2019}
      {Output/scatter_share_med_er_all_2018_2019.pdf}
      {Output/scatter_share_med_er_all_2018_2019_full.pdf}
      {0.6em}

    \twopanels{2019 vrs 2020}
      {Output/scatter_share_med_er_all_2019_2020.pdf}
      {Output/scatter_share_med_er_all_2019_2020_full.pdf}
      {0.6em}

    \twopanels{2018 vrs (2019--2020)}
      {Output/scatter_share_med_er_all_2018_avg.pdf}
      {Output/scatter_share_med_er_all_2018_avg_full.pdf}
      {0em}
\end{figure}

% share of below median evasion rate algo cases vrs total executed/assigned t+1
\begin{figure}[H]
    \centering
    \caption{Share of below median evasion rate $\text{algo}_t$ vrs share of execution of $\text{algo}_{t+1}$}

    \twopanels{2018 vrs 2019}
      {Output/scatter_share_med_er_ALG_2018_2019.pdf}
      {Output/scatter_share_med_er_Algorithm_2018_2019_full.pdf}
      {0.35em}

    \twopanels{2019 vrs 2020}
      {Output/scatter_share_med_er_ALG_2019_2020.pdf}
      {Output/scatter_share_med_er_Algorithm_2019_2020_full.pdf}
      {0.35em}

    \twopanels{2018 vrs (2019--2020)}
      {Output/scatter_share_med_er_ALG_2018_avg.pdf}
      {Output/scatter_share_med_er_Algorithm_2018_avg_full.pdf}
      {0em}
\end{figure}

% share of below median evasion rate vrs share of inspector selected executions  t+1
\begin{figure}[H]
    \centering
    \caption{Share of below median evasion rate $\text{Inspector-selected}_t$ vrs share of execution of $\text{Inspector-selected}_{t+1}$}

    \twopanels{2018 vrs 2019}
      {Output/scatter_share_med_er_Inspectors_2018_2019.pdf}
      {Output/scatter_share_med_er_Inspectors_2018_2019_full.pdf}
      {0.6em}

    \twopanels{2019 vrs 2020}
      {Output/scatter_share_med_er_Inspectors_2019_2020.pdf}
      {Output/scatter_share_med_er_Inspectors_2019_2020_full.pdf}
      {0.6em}

    \twopanels{2018 vrs (2019--2020)}
      {Output/scatter_share_med_er_Inspectors_2018_avg.pdf}
      {Output/scatter_share_med_er_Inspectors_2018_avg_full.pdf}
      {0em}
\end{figure}

%%%%% Share of bottom quartile evasion rate in t vrs share of execution on t+1

% share of all bottom quartile evasion rate cases  vrs total executed/assigned t+1
\begin{figure}[H]
    \centering
    \caption{Share of total $\text{bottom quartile evasion rate cases}_t$ vrs Share of total execution of $\text{assigned}_{t+1}$}

    \twopanels{2018 vrs 2019}
      {Output/scatter_share_botq_er_all_2018_2019.pdf}
      {Output/scatter_share_botq_er_all_2018_2019_full.pdf}
      {0.6em}

    \twopanels{2019 vrs 2020}
      {Output/scatter_share_botq_er_all_2019_2020.pdf}
      {Output/scatter_share_botq_er_all_2019_2020_full.pdf}
      {0.6em}

    \twopanels{2018 vrs (2019--2020)}
      {Output/scatter_share_botq_er_all_2018_avg.pdf}
      {Output/scatter_share_botq_er_all_2018_avg_full.pdf}
      {0em}
\end{figure}

% share of bottom quartile evasion rate algo cases vrs total executed/assigned t+1
\begin{figure}[H]
    \centering
    \caption{Share of bottom quartile evasion rate $\text{algo}_t$ vrs share of execution of $\text{algo}_{t+1}$}

    \twopanels{2018 vrs 2019}
      {Output/scatter_share_botq_er_ALG_2018_2019.pdf}
      {Output/scatter_share_botq_er_Algorithm_2018_2019_full.pdf}
      {0.35em}

    \twopanels{2019 vrs 2020}
      {Output/scatter_share_botq_er_ALG_2019_2020.pdf}
      {Output/scatter_share_botq_er_Algorithm_2019_2020_full.pdf}
      {0.35em}

    \twopanels{2018 vrs (2019--2020)}
      {Output/scatter_share_botq_er_ALG_2018_avg.pdf}
      {Output/scatter_share_botq_er_Algorithm_2018_avg_full.pdf}
      {0em}
\end{figure}

% share of bottom quartile evasion rate vrs share of inspector selected executions  t+1
\begin{figure}[H]
    \centering
    \caption{Share of bottom quartile evasion rate $\text{Inspector-selected}_t$ vrs share of execution of $\text{Inspector-selected}_{t+1}$}

    \twopanels{2018 vrs 2019}
      {Output/scatter_share_botq_er_Inspectors_2018_2019.pdf}
      {Output/scatter_share_botq_er_Inspectors_2018_2019_full.pdf}
      {0.6em}

    \twopanels{2019 vrs 2020}
      {Output/scatter_share_botq_er_Inspectors_2019_2020.pdf}
      {Output/scatter_share_botq_er_Inspectors_2019_2020_full.pdf}
      {0.6em}

    \twopanels{2018 vrs (2019--2020)}
      {Output/scatter_share_botq_er_Inspectors_2018_avg.pdf}
      {Output/scatter_share_botq_er_Inspectors_2018_avg_full.pdf}
      {0em}
\end{figure}







\subsection{Approach 3: Top-listed high-potential algorithm-selected audits}

\textbf{Objective.} Using predicted evasion estimates generated by the Random Forest model, we test whether execution of algorithm-assigned cases depends on whether the "first-in-line" for execution algorithm case(s) on a given list are high- vs. low-potential. The "list" unit differs by audit type: for desk audits it is defined at the inspector $\times$ selection-year level, while for full audits it is defined at the bureau (detailed) $\times$ selection-year level. We compare execution rates of algorithm cases across these lists.

\textbf{Steps.}
\begin{enumerate}
    \item \textbf{Load predicted evasion estimates data which includes estimates for selected audits, and restrict by audit type (full or desk audit).} 
    We use the dataset containing predicted evasion (\texttt{yhatrf}) and execution outcomes, and split the analysis by audit type (\texttt{x2==0} desk; \texttt{x2==1} full).

    \item \textbf{Construct predicted-evasion quantile groups (audit-type specific).}
    We compute quintiles of predicted evasion \texttt{yhatrf} in the following way:
    \begin{itemize}
        \item \textbf{Desk audits:} quintiles within \texttt{verificateur\_selection $\times$ selectionyear}.
        \item \textbf{Full audits:} quintiles within \texttt{bureau\_detailed $\times$ selectionyear}.
    \end{itemize}

    \item \textbf{Restrict to algorithm-assigned cases.}
    We keep only algorithm-selected audits (\texttt{algorithm==1}) so that all outcomes and tags refer to the execution of algorithm-assigned cases.

    \item \textbf{Define "top of the list" using execution sequence among algorithm cases.}
    Within each list, we order algorithm cases by the sequencing variable (\texttt{sequencing}) to define early algorithm cases (e.g., Top 1/Top 2, first 10\% or first 20\% of the algorithm sequence).

    \item \textbf{Tag lists based on the potential of first-in-line algorithm cases.}
    For each list, we create indicators for whether an first-in-line algorithm case is:
    (i) \textbf{high-potential} (top quintile of \texttt{yhatrf}, \texttt{yhatrf\_q5==5}), or 
    (ii) \textbf{low-potential} (bottom quintile, \texttt{yhatrf\_q5==1}).
    This generates list-level groups such as "high-potential-first" vs "low-potential-first".

    \item \textbf{Collapse to the list level and compute execution rates.}
    We collapse the data to the list level and compute the mean execution rate of assigned algorithm cases within each list (mean of \texttt{y2} among \texttt{algorithm==1}).

    \item \textbf{Compare execution outcomes across list groups.}
    We present (i) boxplots of list-level algorithm execution rates (unweighted and weighted by the number of assigned algorithm cases), and (ii) regression estimates comparing execution rates across "high-potential-first" and "low-potential-first" lists (with fixed effects varying by audit type).
\end{enumerate}


\begin{figure}[H]
    \centering
    \caption{Execution rates when the first sequenced algorithm case has high potential evasion versus low potential evasion.}
    \label{fig:exec_rate_boxplots}

    % Panel A: Full audits
    \begin{subfigure}{\textwidth}
        \centering
        \caption{Full audits}
        \includegraphics[width=\textwidth]{Output/execution_rate_boxplot_Full.pdf}
        \label{fig:exec_rate_boxplot_full}
    \end{subfigure}

    \vspace{0.6em}

    % Panel B: Desk audits
    \begin{subfigure}{\textwidth}
        \centering
        \caption{Desk audits}
        \includegraphics[width=\textwidth]{Output/execution_rate_boxplot_Desk.pdf}
        \label{fig:exec_rate_boxplot_desk}
    \end{subfigure}
\end{figure}

\begin{landscape}
\subsubsection{Regression estimates}

    
%%%  Main table 
\begin{table}[H]\centering
\begin{threeparttable}
\caption{Execution of algorithm-assigned audits and high predicted evasion cases at the top of the list [Full audits]}
\label{tab:exec_rate_full_panels}

\footnotesize
\begin{tabular}{l*{9}{c}} \toprule
            &\multicolumn{3}{c}{No FE}                                        &\multicolumn{3}{c}{Selection year FE}                            &\multicolumn{3}{c}{Bureau FE}                                    \\\cmidrule(lr){2-4}\cmidrule(lr){5-7}\cmidrule(lr){8-10}
\multicolumn{1}{l}{\textbf{Outcome: Execution Rate of Algorithm-Selected Cases}} & Top 2 & First 10\% seq & First 20\% seq & Top 2 & First 10\% seq & First 20\% seq & Top 2 & First 10\% seq & First 20\% seq \\ \midrule \multicolumn{10}{l}{\textbf{Panel A: Discrete measures}}\\ \addlinespace
High potential evasion case in list =1&       0.218\sym{*}  &       0.289\sym{**} &       0.207\sym{*}  &       0.176         &       0.224\sym{**} &       0.199\sym{**} &       0.144         &       0.245         &       0.165         \\
            &     (0.124)         &     (0.123)         &     (0.108)         &     (0.105)         &     (0.102)         &     (0.087)         &     (0.185)         &     (0.173)         &     (0.153)         \\
\addlinespace
Low potential evasion case in list =1&      -0.079         &       0.044         &      -0.108         &      -0.079         &      -0.007         &      -0.072         &      -0.067         &       0.085         &      -0.049         \\
            &     (0.108)         &     (0.112)         &     (0.116)         &     (0.105)         &     (0.110)         &     (0.112)         &     (0.122)         &     (0.127)         &     (0.171)         \\
\midrule
Selection Year FE&          No         &          No         &          No         &         Yes         &         Yes         &         Yes         &          No         &          No         &          No         \\
Bureau FE   &          No         &          No         &          No         &          No         &          No         &          No         &         Yes         &         Yes         &         Yes         \\
Observations&          29         &          29         &          29         &          29         &          29         &          29         &          29         &          29         &          29         \\
R-squared   &       0.129         &       0.118         &       0.208         &       0.366         &       0.350         &       0.420         &       0.424         &       0.431         &       0.430         \\

\addlinespace \midrule \multicolumn{10}{l}{\textbf{Panel B: Continuous measure}}\\ \addlinespace

Share of high predicted evasion case at the top&       0.532\sym{**} &       0.530\sym{**} &       0.556\sym{***}&       0.457\sym{**} &       0.403         &       0.481\sym{**} &       0.398         &       0.492         &       0.444         \\
            &     (0.225)         &     (0.240)         &     (0.187)         &     (0.190)         &     (0.238)         &     (0.216)         &     (0.325)         &     (0.333)         &     (0.310)         \\
\midrule
Selection Year FE&          No         &          No         &          No         &         Yes         &         Yes         &         Yes         &          No         &          No         &          No         \\
Bureau FE   &          No         &          No         &          No         &          No         &          No         &          No         &         Yes         &         Yes         &         Yes         \\
Observations&          29         &          29         &          29         &          29         &          29         &          29         &          29         &          29         &          29         \\
R-squared   &       0.112         &       0.076         &       0.097         &       0.350         &       0.313         &       0.343         &       0.416         &       0.422         &       0.417         \\
\bottomrule \end{tabular}


\begin{tablenotes}[flushleft]
\footnotesize
\item Notes: Panel A reports discrete specifications where the dependent variable is the overall execution rate of algorithm cases (amongst the total assigned algorithm cases) and the  regressors are indicators for whether the top of the algorithm ranking contains a high-potential evasion case (and analogously a low-potential case), under alternative definitions of "top of the list" (Top 2, first 10\% of the sequence, first 20\% of the sequence). Panel B reports the continuous specification using the share of high-potential cases at the top of the list. The unit of observation is the audit list (Bureau (detailed)  $\times$ year). Standard errors are robust. Stars denote significance at $^{*}p<0.10$, $^{**}p<0.05$, $^{***}p<0.01$.
\end{tablenotes}
\end{threeparttable}
\end{table}

%%% Discrete Full Audits - Execution rates
\begin{table}[H]\centering
\begin{adjustbox}{max width=1.255\textwidth}
\begin{threeparttable}

\caption{Existence of High potential case vrs overall execution rate of algorithm cases [Full audits]}
\label{tab:inspector_effort_disc_full}

{
\def\sym#1{\ifmmode^{#1}\else\(^{#1}\)\fi}
\begin{tabular}{l*{9}{c}}
\toprule
            &\multicolumn{3}{c}{No FE}                                        &\multicolumn{3}{c}{Selection year FE}                            &\multicolumn{3}{c}{Bureau FE}                                    \\\cmidrule(lr){2-4}\cmidrule(lr){5-7}\cmidrule(lr){8-10}
            &\multicolumn{1}{c}{(1)}&\multicolumn{1}{c}{(2)}&\multicolumn{1}{c}{(3)}&\multicolumn{1}{c}{(4)}&\multicolumn{1}{c}{(5)}&\multicolumn{1}{c}{(6)}&\multicolumn{1}{c}{(7)}&\multicolumn{1}{c}{(8)}&\multicolumn{1}{c}{(9)}\\
            &\multicolumn{1}{c}{Top 2}&\multicolumn{1}{c}{First 10\% seq}&\multicolumn{1}{c}{First 20\% seq}&\multicolumn{1}{c}{Top 2}&\multicolumn{1}{c}{First 10\% seq}&\multicolumn{1}{c}{First 20\% seq}&\multicolumn{1}{c}{Top 2}&\multicolumn{1}{c}{First 10\% seq}&\multicolumn{1}{c}{First 20\% seq}\\
\midrule
\textbf{Dependent variable}: `Outcome: Execution Rate of Algorithm-Selected Cases&                     &                     &                     &                     &                     &                     &                     &                     &                     \\
\addlinespace
High potential evasion case in list =1&       0.218\sym{*}  &       0.289\sym{**} &       0.207\sym{*}  &       0.176         &       0.224\sym{**} &       0.199\sym{**} &       0.144         &       0.245         &       0.165         \\
            &     (0.124)         &     (0.123)         &     (0.108)         &     (0.105)         &     (0.102)         &     (0.087)         &     (0.185)         &     (0.173)         &     (0.153)         \\
\addlinespace
Low potential evasion case in list =1&      -0.079         &       0.044         &      -0.108         &      -0.079         &      -0.007         &      -0.072         &      -0.067         &       0.085         &      -0.049         \\
            &     (0.108)         &     (0.112)         &     (0.116)         &     (0.105)         &     (0.110)         &     (0.112)         &     (0.122)         &     (0.127)         &     (0.171)         \\
\midrule
Selection Year FE&          No         &          No         &          No         &         Yes         &         Yes         &         Yes         &          No         &          No         &          No         \\
Bureau FE   &          No         &          No         &          No         &          No         &          No         &          No         &         Yes         &         Yes         &         Yes         \\
Observations&          29         &          29         &          29         &          29         &          29         &          29         &          29         &          29         &          29         \\
R-squared   &       0.129         &       0.118         &       0.208         &       0.366         &       0.350         &       0.420         &       0.424         &       0.431         &       0.430         \\
\bottomrule
\end{tabular}
}


\begin{tablenotes}[flushleft]
\footnotesize
\item Notes: The regression was estimatated at the list level (Tax office x Selection year), considering only the lists that had assigned algorithm cases in any of the years. The dependent variable is the execution rate of algorithm-selected cases (for each list). All specifications include Selection Year and Bureau fixed effects as indicated. Standard errors in parentheses. 
$^{*}p<0.10$, $^{**}p<0.05$, $^{***}p<0.01$.
\end{tablenotes}
\end{threeparttable}
\end{adjustbox}
\end{table}

%%% Discrete Full Audits - Execution rates of bottom 20% over assigned bottom 20%
\begin{table}[H]\centering
\begin{adjustbox}{max width=1.255\textwidth}
\begin{threeparttable}

\caption{Existence of High potential case vrs execution rate of bottom 20\% cases [Full audits]}
\label{tab:inspector_effort_disc_full}

{
\def\sym#1{\ifmmode^{#1}\else\(^{#1}\)\fi}
\begin{tabular}{l*{9}{c}}
\toprule
            &\multicolumn{3}{c}{No FE}                                        &\multicolumn{3}{c}{Selection year FE}                            &\multicolumn{3}{c}{Bureau FE}                                    \\\cmidrule(lr){2-4}\cmidrule(lr){5-7}\cmidrule(lr){8-10}
            &\multicolumn{1}{c}{(1)}&\multicolumn{1}{c}{(2)}&\multicolumn{1}{c}{(3)}&\multicolumn{1}{c}{(4)}&\multicolumn{1}{c}{(5)}&\multicolumn{1}{c}{(6)}&\multicolumn{1}{c}{(7)}&\multicolumn{1}{c}{(8)}&\multicolumn{1}{c}{(9)}\\
            &\multicolumn{1}{c}{Top 2}&\multicolumn{1}{c}{First 10\% seq}&\multicolumn{1}{c}{First 20\% seq}&\multicolumn{1}{c}{Top 2}&\multicolumn{1}{c}{First 10\% seq}&\multicolumn{1}{c}{First 20\% seq}&\multicolumn{1}{c}{Top 2}&\multicolumn{1}{c}{First 10\% seq}&\multicolumn{1}{c}{First 20\% seq}\\
\midrule
\textbf{Dependent variable}: `Execution rate of bottom-20\% cases&                     &                     &                     &                     &                     &                     &                     &                     &                     \\
\addlinespace
High potential evasion case in list =1&       0.169         &       0.331\sym{*}  &       0.313         &       0.112         &       0.246         &       0.320\sym{*}  &       0.070         &       0.240         &       0.390         \\
            &     (0.207)         &     (0.182)         &     (0.196)         &     (0.188)         &     (0.154)         &     (0.169)         &     (0.363)         &     (0.325)         &     (0.285)         \\
\addlinespace
Low potential evasion case in list =1&      -0.081         &       0.193         &      -0.079         &      -0.077         &       0.128         &      -0.014         &      -0.034         &       0.262         &      -0.007         \\
            &     (0.198)         &     (0.179)         &     (0.209)         &     (0.191)         &     (0.184)         &     (0.180)         &     (0.249)         &     (0.208)         &     (0.250)         \\
\midrule
Selection Year FE&          No         &          No         &          No         &         Yes         &         Yes         &         Yes         &          No         &          No         &          No         \\
Bureau FE   &          No         &          No         &          No         &          No         &          No         &          No         &         Yes         &         Yes         &         Yes         \\
Observations&          27         &          27         &          27         &          27         &          27         &          27         &          27         &          27         &          27         \\
R-squared   &       0.047         &       0.088         &       0.168         &       0.242         &       0.255         &       0.339         &       0.408         &       0.476         &       0.494         \\
\bottomrule
\end{tabular}
}


\begin{tablenotes}[flushleft]
\footnotesize
\item Notes: The regression was estimatated at the list level (Tax office x Selection year), considering only the lists that had assigned algorithm cases in any of the years. The dependent variable is the share of algorithm-selected cases at the bottom 20\% of the predicted evasion distribution over the number of assigned bad cases (for each list). All specifications include Selection Year and Bureau fixed effects as indicated. Standard errors in parentheses. 
$^{*}p<0.10$, $^{**}p<0.05$, $^{***}p<0.01$.
\end{tablenotes}
\end{threeparttable}
\end{adjustbox}
\end{table}

%%% Discrete Full Audits - Execution rates bottom 20%
\begin{table}[H]\centering
\begin{adjustbox}{max width=1.25\textwidth}
\begin{threeparttable}

\caption{Existence of High potential case vrs Executed bottom-20\% \/ assigned (all)  [Full audits]}
\label{tab:inspector_effort_disc_full}

{
\def\sym#1{\ifmmode^{#1}\else\(^{#1}\)\fi}
\begin{tabular}{l*{9}{c}}
\toprule
            &\multicolumn{3}{c}{No FE}                                        &\multicolumn{3}{c}{Selection year FE}                            &\multicolumn{3}{c}{Bureau FE}                                    \\\cmidrule(lr){2-4}\cmidrule(lr){5-7}\cmidrule(lr){8-10}
            &\multicolumn{1}{c}{(1)}&\multicolumn{1}{c}{(2)}&\multicolumn{1}{c}{(3)}&\multicolumn{1}{c}{(4)}&\multicolumn{1}{c}{(5)}&\multicolumn{1}{c}{(6)}&\multicolumn{1}{c}{(7)}&\multicolumn{1}{c}{(8)}&\multicolumn{1}{c}{(9)}\\
            &\multicolumn{1}{c}{Top 2}&\multicolumn{1}{c}{First 10\% seq}&\multicolumn{1}{c}{First 20\% seq}&\multicolumn{1}{c}{Top 2}&\multicolumn{1}{c}{First 10\% seq}&\multicolumn{1}{c}{First 20\% seq}&\multicolumn{1}{c}{Top 2}&\multicolumn{1}{c}{First 10\% seq}&\multicolumn{1}{c}{First 20\% seq}\\
\midrule
\textbf{Dependent variable}: `Executed bottom-20\% \/ assigned (all)&                     &                     &                     &                     &                     &                     &                     &                     &                     \\
\addlinespace
High potential evasion case in list =1&       0.086         &       0.107\sym{**} &       0.127\sym{**} &       0.061         &       0.077         &       0.120\sym{**} &       0.047         &       0.083         &       0.166\sym{**} \\
            &     (0.057)         &     (0.050)         &     (0.050)         &     (0.059)         &     (0.053)         &     (0.056)         &     (0.080)         &     (0.072)         &     (0.075)         \\
\addlinespace
Low potential evasion case in list =1&       0.041         &       0.089         &       0.047         &       0.037         &       0.067         &       0.060         &       0.048         &       0.138\sym{**} &       0.113         \\
            &     (0.058)         &     (0.056)         &     (0.062)         &     (0.063)         &     (0.061)         &     (0.066)         &     (0.068)         &     (0.063)         &     (0.083)         \\
\midrule
Selection Year FE&          No         &          No         &          No         &         Yes         &         Yes         &         Yes         &          No         &          No         &          No         \\
Bureau FE   &          No         &          No         &          No         &          No         &          No         &          No         &         Yes         &         Yes         &         Yes         \\
Observations&          29         &          29         &          29         &          29         &          29         &          29         &          29         &          29         &          29         \\
R-squared   &       0.043         &       0.115         &       0.119         &       0.202         &       0.236         &       0.276         &       0.443         &       0.589         &       0.556         \\
\bottomrule
\end{tabular}
}


\begin{tablenotes}[flushleft]
\footnotesize
\item Notes: The regression was estimatated at the list level (Tax office x Selection year), considering only the lists that had assigned algorithm cases in any of the years. The dependent variable is the execution rate of algorithm-executed cases in the bottom 20\% of predicted evasion distribution over the total number of algo cases assigned (for each list). All specifications include Selection Year and Bureau fixed effects as indicated. Standard errors in parentheses. 
$^{*}p<0.10$, $^{**}p<0.05$, $^{***}p<0.01$.
\end{tablenotes}
\end{threeparttable}
\end{adjustbox}
\end{table}


%%% Continuous Full Audits - overall execution rate
\begin{table}[H]\centering
\begin{adjustbox}{max width=1.255\textwidth}
\begin{threeparttable}
\caption{Share of High potential case vrs overall execution rate [Full audits]}
\label{tab:inspector_effort_disc_full}
\footnotesize
{
\def\sym#1{\ifmmode^{#1}\else\(^{#1}\)\fi}
\begin{tabular}{l*{9}{c}}
\toprule
            &\multicolumn{3}{c}{No FE}                                        &\multicolumn{3}{c}{Selection year FE}                            &\multicolumn{3}{c}{Bureau FE}                                    \\\cmidrule(lr){2-4}\cmidrule(lr){5-7}\cmidrule(lr){8-10}
            &\multicolumn{1}{c}{(1)}&\multicolumn{1}{c}{(2)}&\multicolumn{1}{c}{(3)}&\multicolumn{1}{c}{(4)}&\multicolumn{1}{c}{(5)}&\multicolumn{1}{c}{(6)}&\multicolumn{1}{c}{(7)}&\multicolumn{1}{c}{(8)}&\multicolumn{1}{c}{(9)}\\
            &\multicolumn{1}{c}{Top 2}&\multicolumn{1}{c}{First 10\% seq}&\multicolumn{1}{c}{First 20\% seq}&\multicolumn{1}{c}{Top 2}&\multicolumn{1}{c}{First 10\% seq}&\multicolumn{1}{c}{First 20\% seq}&\multicolumn{1}{c}{Top 2}&\multicolumn{1}{c}{First 10\% seq}&\multicolumn{1}{c}{First 20\% seq}\\
\midrule
\textbf{Dependent variable}: `Outcome: Execution Rate of Algorithm-Selected Cases&                     &                     &                     &                     &                     &                     &                     &                     &                     \\
\addlinespace
Share of potential evasion case at the top&       0.532\sym{**} &       0.530\sym{**} &       0.556\sym{***}&       0.457\sym{**} &       0.403         &       0.481\sym{**} &       0.398         &       0.492         &       0.444         \\
            &     (0.225)         &     (0.240)         &     (0.187)         &     (0.190)         &     (0.238)         &     (0.216)         &     (0.325)         &     (0.333)         &     (0.310)         \\
\midrule
Selection Year FE&          No         &          No         &          No         &         Yes         &         Yes         &         Yes         &          No         &          No         &          No         \\
Bureau FE   &          No         &          No         &          No         &          No         &          No         &          No         &         Yes         &         Yes         &         Yes         \\
Observations&          29         &          29         &          29         &          29         &          29         &          29         &          29         &          29         &          29         \\
R-squared   &       0.112         &       0.076         &       0.097         &       0.350         &       0.313         &       0.343         &       0.416         &       0.422         &       0.417         \\
\bottomrule
\end{tabular}
}


\begin{tablenotes}[flushleft]
\footnotesize
\item Notes: The regression was estimatated at the list level (Tax office x Selection year), considering only the lists that had assigned algorithm cases in any of the years. The dependent variable is the execution rate of algorithm-selected cases (for each list). All specifications include Selection Year and Bureau fixed effects as indicated. Standard errors in parentheses. 
$^{*}p<0.10$, $^{**}p<0.05$, $^{***}p<0.01$.
\end{tablenotes}
\end{threeparttable}
\end{adjustbox}
\end{table}

%%% continuos measure full audits  b20
\begin{table}[H]\centering
\begin{adjustbox}{max width=1.255\textwidth}
\begin{threeparttable}
\caption{Share of High potential case vrs execution rate of bottom 20\% cases [Full audits]}
\label{tab:inspector_effort_disc_full}
\footnotesize
{
\def\sym#1{\ifmmode^{#1}\else\(^{#1}\)\fi}
\begin{tabular}{l*{9}{c}}
\toprule
            &\multicolumn{3}{c}{No FE}                                        &\multicolumn{3}{c}{Selection year FE}                            &\multicolumn{3}{c}{Bureau FE}                                    \\\cmidrule(lr){2-4}\cmidrule(lr){5-7}\cmidrule(lr){8-10}
            &\multicolumn{1}{c}{(1)}&\multicolumn{1}{c}{(2)}&\multicolumn{1}{c}{(3)}&\multicolumn{1}{c}{(4)}&\multicolumn{1}{c}{(5)}&\multicolumn{1}{c}{(6)}&\multicolumn{1}{c}{(7)}&\multicolumn{1}{c}{(8)}&\multicolumn{1}{c}{(9)}\\
            &\multicolumn{1}{c}{Top 2}&\multicolumn{1}{c}{First 10\% seq}&\multicolumn{1}{c}{First 20\% seq}&\multicolumn{1}{c}{Top 2}&\multicolumn{1}{c}{First 10\% seq}&\multicolumn{1}{c}{First 20\% seq}&\multicolumn{1}{c}{Top 2}&\multicolumn{1}{c}{First 10\% seq}&\multicolumn{1}{c}{First 20\% seq}\\
\midrule
\textbf{Dependent variable}: `Execution rate of bottom-20\% cases&                     &                     &                     &                     &                     &                     &                     &                     &                     \\
\addlinespace
Share of potential evasion case at the top&       0.443         &       0.501         &       0.796\sym{**} &       0.331         &       0.329         &       0.708\sym{**} &       0.194         &       0.321         &       0.717         \\
            &     (0.298)         &     (0.371)         &     (0.295)         &     (0.247)         &     (0.359)         &     (0.293)         &     (0.575)         &     (0.672)         &     (0.492)         \\
\midrule
Selection Year FE&          No         &          No         &          No         &         Yes         &         Yes         &         Yes         &          No         &          No         &          No         \\
Bureau FE   &          No         &          No         &          No         &          No         &          No         &          No         &         Yes         &         Yes         &         Yes         \\
Observations&          27         &          27         &          27         &          27         &          27         &          27         &          27         &          27         &          27         \\
R-squared   &       0.039         &       0.034         &       0.099         &       0.235         &       0.228         &       0.292         &       0.407         &       0.412         &       0.449         \\
\bottomrule
\end{tabular}
}


\begin{tablenotes}[flushleft]
\footnotesize
\item Notes: The regression was estimatated at the list level (Tax office x Selection year), considering only the lists that had assigned algorithm cases in any of the years. The dependent variable is share of algorithm-selected cases at the bottom 20\% of the predicted evasion distribution over the number of assigned bad cases (for each list). All specifications include Selection Year and Bureau fixed effects as indicated. Standard errors in parentheses. 
$^{*}p<0.10$, $^{**}p<0.05$, $^{***}p<0.01$.
\end{tablenotes}
\end{threeparttable}
\end{adjustbox}
\end{table}


\begin{table}[H]\centering
\begin{adjustbox}{max width=1.255\textwidth}
\begin{threeparttable}
\caption{Share of High potential case vrs Executed bottom-20\% \/ assigned (all) [Full audits]}
\label{tab:inspector_effort_disc_full}
\footnotesize
{
\def\sym#1{\ifmmode^{#1}\else\(^{#1}\)\fi}
\begin{tabular}{l*{9}{c}}
\toprule
            &\multicolumn{3}{c}{No FE}                                        &\multicolumn{3}{c}{Selection year FE}                            &\multicolumn{3}{c}{Bureau FE}                                    \\\cmidrule(lr){2-4}\cmidrule(lr){5-7}\cmidrule(lr){8-10}
            &\multicolumn{1}{c}{(1)}&\multicolumn{1}{c}{(2)}&\multicolumn{1}{c}{(3)}&\multicolumn{1}{c}{(4)}&\multicolumn{1}{c}{(5)}&\multicolumn{1}{c}{(6)}&\multicolumn{1}{c}{(7)}&\multicolumn{1}{c}{(8)}&\multicolumn{1}{c}{(9)}\\
            &\multicolumn{1}{c}{Top 2}&\multicolumn{1}{c}{First 10\% seq}&\multicolumn{1}{c}{First 20\% seq}&\multicolumn{1}{c}{Top 2}&\multicolumn{1}{c}{First 10\% seq}&\multicolumn{1}{c}{First 20\% seq}&\multicolumn{1}{c}{Top 2}&\multicolumn{1}{c}{First 10\% seq}&\multicolumn{1}{c}{First 20\% seq}\\
\midrule
\textbf{Dependent variable}: `Executed bottom-20\% \/ assigned (all)&                     &                     &                     &                     &                     &                     &                     &                     &                     \\
\addlinespace
Share of potential evasion case at the top&       0.122         &       0.100         &       0.156         &       0.072         &       0.032         &       0.123         &       0.015         &       0.022         &       0.089         \\
            &     (0.096)         &     (0.099)         &     (0.101)         &     (0.086)         &     (0.106)         &     (0.116)         &     (0.101)         &     (0.113)         &     (0.104)         \\
\midrule
Selection Year FE&          No         &          No         &          No         &         Yes         &         Yes         &         Yes         &          No         &          No         &          No         \\
Bureau FE   &          No         &          No         &          No         &          No         &          No         &          No         &         Yes         &         Yes         &         Yes         \\
Observations&          29         &          29         &          29         &          29         &          29         &          29         &          29         &          29         &          29         \\
R-squared   &       0.024         &       0.011         &       0.030         &       0.187         &       0.180         &       0.198         &       0.426         &       0.427         &       0.432         \\
\bottomrule
\end{tabular}
}


\begin{tablenotes}[flushleft]
\footnotesize
\item Notes: The regression was estimatated at the list level (Tax office x Selection year), considering only the lists that had assigned algorithm cases in any of the years. The dependent variable is the execution rate of algorithm-executed cases in the bottom 20\% of predicted evasion distribution over the total number of algo cases assigned (for each list). All specifications include Selection Year and Bureau fixed effects as indicated. Standard errors in parentheses. 
$^{*}p<0.10$, $^{**}p<0.05$, $^{***}p<0.01$.
\end{tablenotes}
\end{threeparttable}
\end{adjustbox}
\end{table}



\end{landscape}
%%% Discrete Desk Audits - exec rate
\begin{landscape}
    

\begin{table}[!htbp]\centering
\begin{threeparttable}
\caption{Existence of High potential case vrs overall execution rate [Desk audits]}
\label{tab:inspector_effort_disc_Desk}
\footnotesize
{
\def\sym#1{\ifmmode^{#1}\else\(^{#1}\)\fi}
\begin{tabular}{l*{4}{c}}
\toprule
            &\multicolumn{1}{c}{(1)}&\multicolumn{1}{c}{(2)}&\multicolumn{1}{c}{(3)}&\multicolumn{1}{c}{(4)}\\
            &\multicolumn{1}{c}{Top 1}&\multicolumn{1}{c}{Top 2}&\multicolumn{1}{c}{First 10\% seq}&\multicolumn{1}{c}{First 20\% seq}\\
\midrule
\textbf{Dependent variable}: `Outcome: Execution Rate of Algorithm-Selected Cases&                     &                     &                     &                     \\
\addlinespace
High potential evasion case in list =1&       0.042         &       0.057         &       0.055         &       0.088\sym{*}  \\
            &     (0.049)         &     (0.048)         &     (0.046)         &     (0.045)         \\
\addlinespace
Low potential evasion case in list =1&       0.018         &       0.071\sym{*}  &       0.068         &       0.064         \\
            &     (0.044)         &     (0.041)         &     (0.045)         &     (0.042)         \\
\midrule
Selection Year FE&         Yes         &         Yes         &         Yes         &         Yes         \\
Bureau FE   &         Yes         &         Yes         &         Yes         &         Yes         \\
Observations&         165         &         165         &         165         &         165         \\
R-squared   &       0.521         &       0.528         &       0.527         &       0.531         \\
\bottomrule
\end{tabular}
}


\begin{tablenotes}[flushleft]
\footnotesize
\item Notes: The regression was estimatated at the list level (Inspector x Selection year), considering only the lists that had assigned algorithm cases in any of the years. The dependent variable is the execution rate of algorithm-selected cases (for each list). All specifications include Selection Year and Bureau fixed effects as indicated. Standard errors in parentheses. 
$^{*}p<0.10$, $^{**}p<0.05$, $^{***}p<0.01$.
\end{tablenotes}
\end{threeparttable}
\end{table}

%%% Discrete Desk Audits -b20


\begin{table}[H]\centering
\begin{threeparttable}
\caption{Existence of High potential case vrs execution rate of bottom 20\% cases [Desk audits]}
\label{tab:inspector_effort_disc_Desk}
\footnotesize
{
\def\sym#1{\ifmmode^{#1}\else\(^{#1}\)\fi}
\begin{tabular}{l*{4}{c}}
\toprule
            &\multicolumn{1}{c}{(1)}&\multicolumn{1}{c}{(2)}&\multicolumn{1}{c}{(3)}&\multicolumn{1}{c}{(4)}\\
            &\multicolumn{1}{c}{Top 1}&\multicolumn{1}{c}{Top 2}&\multicolumn{1}{c}{First 10\% seq}&\multicolumn{1}{c}{First 20\% seq}\\
\midrule
\textbf{Dependent variable}: `Execution rate of bottom-20\% cases&                     &                     &                     &                     \\
\addlinespace
High potential evasion case in list =1&       0.113         &       0.183\sym{**} &       0.136         &       0.167\sym{**} \\
            &     (0.095)         &     (0.073)         &     (0.091)         &     (0.080)         \\
\addlinespace
Low potential evasion case in list =1&       0.088         &       0.150\sym{**} &       0.155\sym{**} &       0.108         \\
            &     (0.072)         &     (0.068)         &     (0.076)         &     (0.076)         \\
\midrule
Selection Year FE&         Yes         &         Yes         &         Yes         &         Yes         \\
Bureau FE   &         Yes         &         Yes         &         Yes         &         Yes         \\
Observations&         144         &         144         &         144         &         144         \\
R-squared   &       0.422         &       0.441         &       0.436         &       0.431         \\
\bottomrule
\end{tabular}
}


\begin{tablenotes}[flushleft]
\footnotesize
\item Notes: The regression was estimatated at the list level (Inspector x Selection year), considering only the lists that had assigned algorithm cases in any of the years. The dependent variable is the share of algorithm-selected
cases at the bottom 20\% of the predicted evasion distribution over the number of assigned bad cases (for each list). All specifications include Selection Year and Bureau fixed effects as indicated. Standard errors in parentheses. 
$^{*}p<0.10$, $^{**}p<0.05$, $^{***}p<0.01$.
\end{tablenotes}
\end{threeparttable}
\end{table}

%%% Discrete Desk Audits -b20t
    

\begin{table}[H]\centering
\begin{threeparttable}
\caption{Existence of High potential case vrs Executed bottom-20\% assigned (all)  [Desk audits]}
\label{tab:inspector_effort_disc_Desk}
\footnotesize
{
\def\sym#1{\ifmmode^{#1}\else\(^{#1}\)\fi}
\begin{tabular}{l*{4}{c}}
\toprule
            &\multicolumn{1}{c}{(1)}&\multicolumn{1}{c}{(2)}&\multicolumn{1}{c}{(3)}&\multicolumn{1}{c}{(4)}\\
            &\multicolumn{1}{c}{Top 1}&\multicolumn{1}{c}{Top 2}&\multicolumn{1}{c}{First 10\% seq}&\multicolumn{1}{c}{First 20\% seq}\\
\midrule
\textbf{Dependent variable}: `Executed bottom-20\% \/ assigned (all)&                     &                     &                     &                     \\
\addlinespace
High potential evasion case in list =1&       0.027         &       0.059\sym{***}&       0.057\sym{**} &       0.052\sym{**} \\
            &     (0.022)         &     (0.022)         &     (0.025)         &     (0.024)         \\
\addlinespace
Low potential evasion case in list =1&       0.076\sym{**} &       0.081\sym{***}&       0.104\sym{***}&       0.068\sym{**} \\
            &     (0.032)         &     (0.025)         &     (0.030)         &     (0.027)         \\
\midrule
Selection Year FE&         Yes         &         Yes         &         Yes         &         Yes         \\
Bureau FE   &         Yes         &         Yes         &         Yes         &         Yes         \\
Observations&         165         &         165         &         165         &         165         \\
R-squared   &       0.491         &       0.495         &       0.513         &       0.485         \\
\bottomrule
\end{tabular}
}


\begin{tablenotes}[flushleft]
\footnotesize
\item Notes: The regression was estimatated at the list level (Inspector x Selection year), considering only the lists that had assigned algorithm cases in any of the years.  The dependent variable is the execution rate of algorithm-executed cases in the bottom 20\% of predicted evasion distribution over the total number of algo cases assigned (for each list). All specifications include Selection Year and Bureau fixed effects as indicated. Standard errors in parentheses. 
$^{*}p<0.10$, $^{**}p<0.05$, $^{***}p<0.01$.
\end{tablenotes}
\end{threeparttable}
\end{table}


%%% Continuous Desk Audits

\begin{table}[!htbp]\centering
\begin{threeparttable}
\caption{Share of High potential case  vrs overall execution rate [Desk audits]}
\label{tab:inspector_effort_share_Desk}
\footnotesize
{
\def\sym#1{\ifmmode^{#1}\else\(^{#1}\)\fi}
\begin{tabular}{l*{4}{c}}
\toprule
            &\multicolumn{1}{c}{(1)}&\multicolumn{1}{c}{(2)}&\multicolumn{1}{c}{(3)}&\multicolumn{1}{c}{(4)}\\
            &\multicolumn{1}{c}{Top 1}&\multicolumn{1}{c}{Top 2}&\multicolumn{1}{c}{First 10\% seq}&\multicolumn{1}{c}{First 20\% seq}\\
\midrule
\textbf{Dependent variable}: `Outcome: Execution Rate of Algorithm-Selected Cases&                     &                     &                     &                     \\
\addlinespace
Share of high predicted evasion case at the top&       0.038         &       0.008         &       0.046         &       0.077         \\
            &     (0.048)         &     (0.064)         &     (0.048)         &     (0.062)         \\
\midrule
Selection Year FE&         Yes         &         Yes         &         Yes         &         Yes         \\
Bureau FE   &         Yes         &         Yes         &         Yes         &         Yes         \\
Observations&         165         &         165         &         165         &         165         \\
R-squared   &       0.521         &       0.519         &       0.521         &       0.524         \\
\bottomrule
\end{tabular}
}


\begin{tablenotes}[flushleft]
\footnotesize
\item Notes: The regression was estimatated at the list level (Inspector x Selection year), considering only the lists that had assigned algorithm cases in any of the years. The dependent variable is the execution rate of algorithm-selected cases (for each list). All specifications include Selection Year and Bureau fixed effects as indicated. Standard errors in parentheses. 
$^{*}p<0.10$, $^{**}p<0.05$, $^{***}p<0.01$.
\end{tablenotes}
\end{threeparttable}
\end{table}

%%% Discrete Desk Audits -b20


\begin{table}[H]\centering
\begin{threeparttable}
\caption{Share of High potential case vrs execution rate of bottom 20\% cases [Desk audits]}
\label{tab:inspector_effort_share_Desk}
\footnotesize
{
\def\sym#1{\ifmmode^{#1}\else\(^{#1}\)\fi}
\begin{tabular}{l*{4}{c}}
\toprule
            &\multicolumn{1}{c}{(1)}&\multicolumn{1}{c}{(2)}&\multicolumn{1}{c}{(3)}&\multicolumn{1}{c}{(4)}\\
            &\multicolumn{1}{c}{Top 1}&\multicolumn{1}{c}{Top 2}&\multicolumn{1}{c}{First 10\% seq}&\multicolumn{1}{c}{First 20\% seq}\\
\midrule
\textbf{Dependent variable}: `Execution rate of bottom-20\% cases&                     &                     &                     &                     \\
\addlinespace
Share of high predicted evasion case at the top&       0.081         &       0.140         &       0.095         &       0.151         \\
            &     (0.094)         &     (0.113)         &     (0.099)         &     (0.120)         \\
\midrule
Selection Year FE&         Yes         &         Yes         &         Yes         &         Yes         \\
Bureau FE   &         Yes         &         Yes         &         Yes         &         Yes         \\
Observations&         144         &         144         &         144         &         144         \\
R-squared   &       0.415         &       0.417         &       0.416         &       0.419         \\
\bottomrule
\end{tabular}
}


\begin{tablenotes}[flushleft]
\footnotesize
\item Notes: The regression was estimatated at the list level (Inspector x Selection year), considering only the lists that had assigned algorithm cases in any of the years. The dependent variable is the share of algorithm-selected
cases at the bottom 20\% of the predicted evasion distribution over the number of assigned bad cases (for each list). All specifications include Selection Year and Bureau fixed effects as indicated. Standard errors in parentheses. 
$^{*}p<0.10$, $^{**}p<0.05$, $^{***}p<0.01$.
\end{tablenotes}
\end{threeparttable}
\end{table}

%%% Discrete Desk Audits -b20t
    

\begin{table}[H]\centering
\begin{threeparttable}
\caption{Share of High potential case vrs Executed bottom-20\% assigned (all)  [Desk audits]}
\label{tab:inspector_effort_share_Desk}
\footnotesize
{
\def\sym#1{\ifmmode^{#1}\else\(^{#1}\)\fi}
\begin{tabular}{l*{4}{c}}
\toprule
            &\multicolumn{1}{c}{(1)}&\multicolumn{1}{c}{(2)}&\multicolumn{1}{c}{(3)}&\multicolumn{1}{c}{(4)}\\
            &\multicolumn{1}{c}{Top 1}&\multicolumn{1}{c}{Top 2}&\multicolumn{1}{c}{First 10\% seq}&\multicolumn{1}{c}{First 20\% seq}\\
\midrule
\textbf{Dependent variable}: `Executed bottom-20\% \/ assigned (all)&                     &                     &                     &                     \\
\addlinespace
Share of high predicted evasion case at the top&       0.006         &       0.020         &       0.024         &       0.033         \\
            &     (0.022)         &     (0.036)         &     (0.024)         &     (0.031)         \\
\midrule
Selection Year FE&         Yes         &         Yes         &         Yes         &         Yes         \\
Bureau FE   &         Yes         &         Yes         &         Yes         &         Yes         \\
Observations&         165         &         165         &         165         &         165         \\
R-squared   &       0.464         &       0.465         &       0.466         &       0.467         \\
\bottomrule
\end{tabular}
}


\begin{tablenotes}[flushleft]
\footnotesize
\item Notes: The regression was estimatated at the list level (Inspector x Selection year), considering only the lists that had assigned algorithm cases in any of the years.  The dependent variable is the execution rate of algorithm-executed cases in the bottom 20\% of predicted evasion distribution over the total number of algo cases assigned (for each list). All specifications include Selection Year and Bureau fixed effects as indicated. Standard errors in parentheses. 
$^{*}p<0.10$, $^{**}p<0.05$, $^{***}p<0.01$.
\end{tablenotes}
\end{threeparttable}
\end{table}




\end{landscape}







%%% Probability of execution according to ranking sequence
\begin{landscape}
\begin{table}[!htbp]\centering
\begin{threeparttable}
\caption{Probability of execution according to ranking sequence}
\label{tab:inspector_effort_disc_Desk}
\footnotesize
{
\def\sym#1{\ifmmode^{#1}\else\(^{#1}\)\fi}
\begin{tabular}{l*{9}{c}}
\toprule
                &\multicolumn{3}{c}{Sequence}                            &\multicolumn{3}{c}{Rank/ Total assigned ratio}          &\multicolumn{3}{c}{Top 5 in sequence=1}                 \\\cmidrule(lr){2-4}\cmidrule(lr){5-7}\cmidrule(lr){8-10}
                &\multicolumn{1}{c}{(1)}&\multicolumn{1}{c}{(2)}&\multicolumn{1}{c}{(3)}&\multicolumn{1}{c}{(4)}&\multicolumn{1}{c}{(5)}&\multicolumn{1}{c}{(6)}&\multicolumn{1}{c}{(7)}&\multicolumn{1}{c}{(8)}&\multicolumn{1}{c}{(9)}\\
                &\multicolumn{1}{c}{Full audits}&\multicolumn{1}{c}{Desk audits}&\multicolumn{1}{c}{Desk audits}&\multicolumn{1}{c}{Full audits}&\multicolumn{1}{c}{Desk audits }&\multicolumn{1}{c}{Desk audits}&\multicolumn{1}{c}{Full audits}&\multicolumn{1}{c}{Desk audits }&\multicolumn{1}{c}{Desk audits}\\
\midrule
Rank in list (seq)&   -0.002\sym{**} &   -0.001         &   -0.001         &                  &                  &                  &                  &                  &                  \\
                &  (0.001)         &  (0.001)         &  (0.001)         &                  &                  &                  &                  &                  &                  \\
\addlinespace
Rank / total assigned&                  &                  &                  &   -0.109\sym{**} &   -0.029         &   -0.032         &                  &                  &                  \\
                &                  &                  &                  &  (0.049)         &  (0.033)         &  (0.032)         &                  &                  &                  \\
\addlinespace
Seq $\le$ 5=1   &                  &                  &                  &                  &                  &                  &    0.057         &    0.011         &    0.008         \\
                &                  &                  &                  &                  &                  &                  &  (0.042)         &  (0.019)         &  (0.018)         \\
\addlinespace
Algorithm       &   -0.168\sym{***}&   -0.045\sym{**} &   -0.043\sym{**} &   -0.167\sym{***}&   -0.046\sym{**} &   -0.044\sym{**} &   -0.169\sym{***}&   -0.046\sym{**} &   -0.044\sym{**} \\
                &  (0.030)         &  (0.018)         &  (0.017)         &  (0.030)         &  (0.018)         &  (0.017)         &  (0.030)         &  (0.018)         &  (0.017)         \\
\addlinespace
Overlap         &    0.157\sym{**} &    0.049         &    0.047         &    0.156\sym{**} &    0.049         &    0.048         &    0.152\sym{**} &    0.049         &    0.047         \\
                &  (0.067)         &  (0.035)         &  (0.036)         &  (0.068)         &  (0.035)         &  (0.036)         &  (0.068)         &  (0.035)         &  (0.036)         \\
\addlinespace
Random          &                  &   -0.001         &   -0.000         &                  &   -0.001         &   -0.000         &                  &   -0.001         &   -0.000         \\
                &                  &  (0.029)         &  (0.028)         &                  &  (0.029)         &  (0.028)         &                  &  (0.029)         &  (0.028)         \\
\midrule
Mean of dep. var.&     0.53         &     0.37         &     0.37         &     0.53         &     0.37         &     0.37         &     0.53         &     0.37         &     0.37         \\
N               &      944         &     2731         &     2731         &      944         &     2731         &     2731         &      944         &     2731         &     2731         \\
Tax center × year FE&      Yes         &      Yes         &      Yes         &      Yes         &      Yes         &      Yes         &      Yes         &      Yes         &      Yes         \\
Inspector × year FE&       No         &       No         &      Yes         &       No         &       No         &      Yes         &       No         &       No         &      Yes         \\
\bottomrule
\end{tabular}
}


\begin{tablenotes}[flushleft]
\footnotesize
\item Notes: $^{*}p<0.10$, $^{**}p<0.05$, $^{***}p<0.01$. OLS results of a regression of the main audit outcomes on the selection method of the case, for full and desk
audits separately. Specifications for full audits control for tax office x year fixed effects, which corresponds to the selection list for full audits. For desk audits,
we show specifications with tax office x year fixed effects, and inspector x year fixed effects, the latter corresponding to the list for desk audits  

\end{tablenotes}
\end{threeparttable}
\end{table}
\end{landscape}

\section{Measures of inspector effort by predicted evasion}

\subsection{Whole Sample Bins}
The bin-plot analysis uses the Random Forest predicted evasion score \texttt{yhatrf}, generated from the replication R pipeline and stored in \texttt{fullaudits\_predicted.dta} / \texttt{deskaudits\_predicted.dta}. The model is trained on executed audits with observed realized evasion (\texttt{y2==1} and non-missing \texttt{y4}) and then applied to the full selected sample.

\begin{itemize}
\item The two prediction files (\texttt{fullaudits\_predicted.dta} and \texttt{deskaudits\_predicted.dta}) are merged, and all steps are run separately by audit type (\texttt{x2==1} full, \texttt{x2==0} desk).
\item Admin duration is cleaned by setting \texttt{y19=.} when \texttt{y19==0}.
\item Firm size is constructed as mean turnover over 2014--2020 (\texttt{turnover\_mean}).
\item For each x-variable (\texttt{yhatrf}, \texttt{turnover\_mean}), bins are created on the pooled selected sample using quantile bins (quintiles and deciles) and fixed-width bins (\texttt{w1/w2/w4/w5}).
 Bin assignment is done before \texttt{keep if y2==1}, so executed-only filtering is applied later for number of inspectors and duration outcomes.
\item For number of inspectors and duration plots, only executed audits are kept (\texttt{y2==1}), data are collapsed by method and bin for \texttt{numberagents}, \texttt{y19}, \texttt{q30}, and \texttt{y8}, and 95\% confidence intervals are computed using outcome-specific non-missing counts.
\item Algorithm and inspector series are plotted with confidence bands, and one PDF is exported per outcome, x-variable, audit type, bin definition, and bin scope.
\end{itemize}

\subsubsection{Predicted evasion}

\BinCompareQuantilePanelsOutcomeAvailFigure
{bin_ninspectors_yhatrf_bins_outcomeavail}
{Average number of inspectors by predicted evasion bins (outcome-specific executed-sample quantile binning).}
{avg_ninspectors_binplot_yhatrf}

\BinCompareQuantilePanelsOutcomeAvailFigure
{bin_duration_admin_yhatrf_bins_outcomeavail}
{Average duration (admin data) by predicted evasion bins (outcome-specific executed-sample quantile binning).}
{avg_duration_admin_binplot_yhatrf}

\BinCompareQuantilePanelsOutcomeAvailFigure
{bin_duration_tp_yhatrf_bins_outcomeavail}
{Average duration (taxpayer survey) by predicted evasion bins (outcome-specific executed-sample quantile binning).}
{avg_duration_tp_binplot_yhatrf}

\BinCompareQuantilePanelsOutcomeAvailFigure
{bin_duration_self_yhatrf_bins_outcomeavail}
{Average duration (self-reported) by predicted evasion bins (outcome-specific executed-sample quantile binning).}
{avg_duration_self_binplot_yhatrf}

\BinCompareBinsPanelsWholeSampleFigure
{bin_ninspectors_yhatrf_bins_whole_sample}
{Average number of inspectors by predicted evasion bins (Executed-sample binning).}
{avg_ninspectors_binplot_yhatrf}

\BinCompareBinsPanelsWholeSampleFigure
{bin_duration_admin_yhatrf_bins_whole_sample}
{Average duration (admin data) by predicted evasion bins (Executed-sample binning).}
{avg_duration_admin_binplot_yhatrf}

\BinCompareBinsPanelsWholeSampleFigure
{bin_duration_tp_yhatrf_bins_whole_sample}
{Average duration (taxpayer survey) by predicted evasion bins (Executed-sample binning).}
{avg_duration_tp_binplot_yhatrf}

\BinCompareBinsPanelsWholeSampleFigure
{bin_duration_self_yhatrf_bins_whole_sample}
{Average duration (self-reported) by predicted evasion bins (Executed-sample binning).}
{avg_duration_self_binplot_yhatrf}

\begin{landscape}

\begin{table}[!htbp]
\centering
\resizebox{\linewidth}{!}{%
\begin{threeparttable}

\caption{Support for whole-sample predicted-evasion decile effort plots, full audits}
\label{tab:bin_support_yhatrf_full_decile_whole_sample}

\begin{tabular}{llrr*{4}{rrrrrr}}
\toprule
 & & & & \multicolumn{6}{c}{Number of inspectors} & \multicolumn{6}{c}{Admin duration} & \multicolumn{6}{c}{Taxpayer duration} & \multicolumn{6}{c}{Self-reported duration} \\\\
\cmidrule(lr){5-10}\cmidrule(lr){11-16}\cmidrule(lr){17-22}\cmidrule(lr){23-28}
Decile & Method & Selected N & Executed N & N & Mean & SD & SE & CI low & CI high & N & Mean & SD & SE & CI low & CI high & N & Mean & SD & SE & CI low & CI high & N & Mean & SD & SE & CI low & CI high \\\\
\midrule
1 & Algorithm & 40 & 18 & 18 & 1.56 & 0.78 & 0.18 & 1.19 & 1.92 & 7 & 110.14 & 76.52 & 28.92 & 53.45 & 166.83 & 1 & 7.00 & 0.00 & 0.00 & 7.00 & 7.00 & 4 & 41.50 & 24.13 & 12.07 & 17.85 & 65.15 \\\\
1 & Inspector & 55 & 35 & 35 & 1.71 & 0.93 & 0.16 & 1.41 & 2.02 & 17 & 135.88 & 106.14 & 25.74 & 85.42 & 186.34 & 5 & 25.20 & 10.62 & 4.75 & 15.89 & 34.51 & 7 & 59.86 & 21.31 & 8.05 & 44.07 & 75.64 \\\\
2 & Algorithm & 60 & 3 & 3 & 2.67 & 0.58 & 0.33 & 2.01 & 3.32 & 0 &  &  &  &  &  & 0 &  &  &  &  &  & 0 &  &  &  &  &  \\\\
2 & Inspector & 34 & 2 & 2 & 2.50 & 0.71 & 0.50 & 1.52 & 3.48 & 0 &  &  &  &  &  & 1 & 28.00 & 0.00 & 0.00 & 28.00 & 28.00 & 0 &  &  &  &  &  \\\\
3 & Algorithm & 51 & 14 & 14 & 2.43 & 0.51 & 0.14 & 2.16 & 2.70 & 3 & 277.00 & 88.36 & 51.01 & 177.01 & 376.99 & 2 & 45.50 & 24.75 & 17.50 & 11.20 & 79.80 & 0 &  &  &  &  &  \\\\
3 & Inspector & 44 & 12 & 12 & 2.92 & 0.67 & 0.19 & 2.54 & 3.29 & 4 & 180.75 & 204.94 & 102.47 & 0.00 & 381.59 & 2 & 31.50 & 34.65 & 24.50 & 0.00 & 79.52 & 1 & 120.00 & 0.00 & 0.00 & 120.00 & 120.00 \\\\
4 & Algorithm & 54 & 24 & 24 & 2.92 & 0.65 & 0.13 & 2.66 & 3.18 & 5 & 173.40 & 160.22 & 71.65 & 32.96 & 313.84 & 6 & 24.50 & 23.74 & 9.69 & 5.51 & 43.49 & 3 & 233.33 & 128.97 & 74.46 & 87.39 & 379.28 \\\\
4 & Inspector & 40 & 15 & 15 & 2.80 & 0.56 & 0.14 & 2.52 & 3.08 & 4 & 179.25 & 90.92 & 45.46 & 90.15 & 268.35 & 4 & 42.00 & 29.70 & 14.85 & 12.90 & 71.10 & 2 & 62.00 & 82.02 & 58.00 & 0.00 & 175.68 \\\\
5 & Algorithm & 55 & 27 & 27 & 2.85 & 0.60 & 0.12 & 2.62 & 3.08 & 9 & 89.78 & 41.57 & 13.86 & 62.62 & 116.94 & 7 & 22.00 & 20.38 & 7.70 & 6.90 & 37.10 & 5 & 246.00 & 120.75 & 54.00 & 140.16 & 351.84 \\\\
5 & Inspector & 39 & 23 & 23 & 3.00 & 0.74 & 0.15 & 2.70 & 3.30 & 10 & 160.50 & 87.47 & 27.66 & 106.28 & 214.72 & 6 & 64.17 & 57.11 & 23.32 & 18.47 & 109.87 & 4 & 163.75 & 96.21 & 48.10 & 69.46 & 258.04 \\\\
6 & Algorithm & 39 & 26 & 26 & 2.58 & 0.70 & 0.14 & 2.31 & 2.85 & 10 & 89.10 & 63.86 & 20.19 & 49.52 & 128.68 & 5 & 29.40 & 26.84 & 12.00 & 5.88 & 52.92 & 2 & 179.00 & 154.15 & 109.00 & 0.00 & 392.64 \\\\
6 & Inspector & 56 & 34 & 34 & 3.00 & 0.60 & 0.10 & 2.80 & 3.20 & 18 & 187.06 & 133.23 & 31.40 & 125.51 & 248.60 & 8 & 33.25 & 25.31 & 8.95 & 15.71 & 50.79 & 5 & 178.00 & 60.17 & 26.91 & 125.26 & 230.74 \\\\
7 & Algorithm & 40 & 26 & 26 & 2.65 & 0.63 & 0.12 & 2.41 & 2.90 & 16 & 140.06 & 93.49 & 23.37 & 94.25 & 185.87 & 11 & 25.45 & 31.81 & 9.59 & 6.65 & 44.25 & 1 & 189.00 & 0.00 & 0.00 & 189.00 & 189.00 \\\\
7 & Inspector & 54 & 42 & 42 & 3.00 & 0.66 & 0.10 & 2.80 & 3.20 & 24 & 124.96 & 96.97 & 19.79 & 86.16 & 163.75 & 13 & 47.38 & 47.06 & 13.05 & 21.80 & 72.97 & 7 & 265.71 & 125.41 & 47.40 & 172.81 & 358.62 \\\\
8 & Algorithm & 18 & 10 & 10 & 3.00 & 0.47 & 0.15 & 2.71 & 3.29 & 7 & 146.43 & 101.76 & 38.46 & 71.04 & 221.82 & 1 & 21.00 & 0.00 & 0.00 & 21.00 & 21.00 & 0 &  &  &  &  &  \\\\
8 & Inspector & 77 & 52 & 52 & 3.12 & 0.78 & 0.11 & 2.90 & 3.33 & 32 & 150.81 & 86.03 & 15.21 & 121.00 & 180.62 & 8 & 23.62 & 14.94 & 5.28 & 13.27 & 33.98 & 6 & 218.50 & 103.43 & 42.22 & 135.74 & 301.26 \\\\
9 & Algorithm & 30 & 20 & 20 & 3.10 & 1.02 & 0.23 & 2.65 & 3.55 & 11 & 200.18 & 154.50 & 46.58 & 108.88 & 291.49 & 3 & 32.67 & 44.46 & 25.67 & 0.00 & 82.97 & 2 & 126.00 & 8.49 & 6.00 & 114.24 & 137.76 \\\\
9 & Inspector & 64 & 46 & 46 & 3.07 & 0.65 & 0.10 & 2.88 & 3.25 & 34 & 158.88 & 86.83 & 14.89 & 129.69 & 188.07 & 13 & 21.54 & 15.77 & 4.37 & 12.96 & 30.11 & 3 & 256.00 & 109.50 & 63.22 & 132.09 & 379.91 \\\\
10 & Algorithm & 20 & 17 & 17 & 3.24 & 0.83 & 0.20 & 2.84 & 3.63 & 11 & 181.45 & 138.08 & 41.63 & 99.85 & 263.06 & 1 & 14.00 & 0.00 & 0.00 & 14.00 & 14.00 & 0 &  &  &  &  &  \\\\
10 & Inspector & 74 & 62 & 62 & 3.53 & 0.90 & 0.11 & 3.31 & 3.76 & 55 & 211.38 & 130.99 & 17.66 & 176.76 & 246.00 & 11 & 26.09 & 15.37 & 4.63 & 17.01 & 35.17 & 2 & 154.00 & 98.99 & 70.00 & 16.80 & 291.20 \\\\
\bottomrule
\end{tabular}


\begin{tablenotes}[flushleft]
\footnotesize
\item Notes: Predicted-evasion deciles are formed on the pooled sample of selected full audits. Outcome statistics are computed for executed cases only (\texttt{y2==1}). \texttt{Selected N} reports all selected cases in each method-by-decile cell, while \texttt{Executed N} and all outcome moments are based on executed cases; outcome-specific columns further restrict to non-missing observations for that outcome among executed cases only. Confidence intervals are 95\% normal-approximation intervals constructed to match the corresponding bin plots.
\end{tablenotes}
\end{threeparttable}
}
\end{table}

\end{landscape}

\subsubsection{Firm turnover}
\BinCompareQuantilePanelsOutcomeAvailFigure
{bin_ninspectors_turnover_bins_outcomeavail}
{Average number of inspectors by firm turnover bins (outcome-specific executed-sample quantile binning).}
{avg_ninspectors_binplot_turnover_mean}

\BinCompareQuantilePanelsOutcomeAvailFigure
{bin_duration_admin_turnover_bins_outcomeavail}
{Average duration (admin data) by firm turnover bins (outcome-specific executed-sample quantile binning).}
{avg_duration_admin_binplot_turnover_mean}

\BinCompareQuantilePanelsOutcomeAvailFigure
{bin_duration_tp_turnover_bins_outcomeavail}
{Average duration (taxpayer survey) by firm turnover bins (outcome-specific executed-sample quantile binning).}
{avg_duration_tp_binplot_turnover_mean}

\BinCompareQuantilePanelsOutcomeAvailFigure
{bin_duration_self_turnover_bins_outcomeavail}
{Average duration (self-reported) by firm turnover bins (outcome-specific executed-sample quantile binning).}
{avg_duration_self_binplot_turnover_mean}

\BinCompareBinsPanelsWholeSampleFigure
{bin_ninspectors_turnover_bins_whole_sample}
{Average number of inspectors by firm turnover bins (Executed-sample binning).}
{avg_ninspectors_binplot_turnover_mean}

\BinCompareBinsPanelsWholeSampleFigure
{bin_duration_admin_turnover_bins_whole_sample}
{Average duration (admin data) by firm turnover bins (Executed-sample binning).}
{avg_duration_admin_binplot_turnover_mean}

\BinCompareBinsPanelsWholeSampleFigure
{bin_duration_tp_turnover_bins_whole_sample}
{Average duration (taxpayer survey) by firm turnover bins (Executed-sample binning).}
{avg_duration_tp_binplot_turnover_mean}

\BinCompareBinsPanelsWholeSampleFigure
{bin_duration_self_turnover_bins_whole_sample}
{Average duration (self-reported) by firm turnover bins (Executed-sample binning).}
{avg_duration_self_binplot_turnover_mean}

\subsection{Within Algo/Inspector Estimated Bins}

\subsubsection{Predicted evasion}

\BinCompareBinsPanelsWithinMethodFigure
{bin_ninspectors_yhatrf_bins_within_method}
{Average number of inspectors by predicted evasion bins (within algorithm/inspector estimated bins).}
{avg_ninspectors_binplot_yhatrf}

\BinCompareBinsPanelsWithinMethodFigure
{bin_duration_admin_yhatrf_bins_within_method}
{Average duration (admin data) by predicted evasion bins (within algorithm/inspector estimated bins).}
{avg_duration_admin_binplot_yhatrf}

\BinCompareBinsPanelsWithinMethodFigure
{bin_duration_tp_yhatrf_bins_within_method}
{Average duration (taxpayer survey) by predicted evasion bins (within algorithm/inspector estimated bins).}
{avg_duration_tp_binplot_yhatrf}

\BinCompareBinsPanelsWithinMethodFigure
{bin_duration_self_yhatrf_bins_within_method}
{Average duration (self-reported) by predicted evasion bins (within algorithm/inspector estimated bins).}
{avg_duration_self_binplot_yhatrf}

\subsubsection{Firm turnover}

\BinCompareBinsPanelsWithinMethodFigure
{bin_ninspectors_turnover_bins_within_method}
{Average number of inspectors by firm turnover bins (within algorithm/inspector estimated bins).}
{avg_ninspectors_binplot_turnover_mean}

\BinCompareBinsPanelsWithinMethodFigure
{bin_duration_admin_turnover_bins_within_method}
{Average duration (admin data) by firm turnover bins (within algorithm/inspector estimated bins).}
{avg_duration_admin_binplot_turnover_mean}

\BinCompareBinsPanelsWithinMethodFigure
{bin_duration_tp_turnover_bins_within_method}
{Average duration (taxpayer survey) by firm turnover bins (within algorithm/inspector estimated bins).}
{avg_duration_tp_binplot_turnover_mean}

\BinCompareBinsPanelsWithinMethodFigure
{bin_duration_self_turnover_bins_within_method}
{Average duration (self-reported) by firm turnover bins (within algorithm/inspector estimated bins).}
{avg_duration_self_binplot_turnover_mean}


%%% inspector characteristics selection desk audits
\begin{landscape}
\subsection{More on inspector characteristics performance and firm size }
\begin{table}[!htbp]\centering
\begin{adjustbox}{max width=1.4\textwidth}
\begin{threeparttable}
\caption{Inspector Characteristics and firm size, Desk Audits}
\label{tab:inspector_effort_disc_Desk}

{
\def\sym#1{\ifmmode^{#1}\else\(^{#1}\)\fi}
\begin{tabular}{l*{12}{c}}
\toprule
            &\multicolumn{4}{c}{All}                                                                &\multicolumn{4}{c}{Algorithm cases}                                                    &\multicolumn{4}{c}{Inspector cases}                                                    \\\cmidrule(lr){2-5}\cmidrule(lr){6-9}\cmidrule(lr){10-13}
            &\multicolumn{1}{c}{\shortstack{log(Turnover)\\(Mean)}}&\multicolumn{1}{c}{\shortstack{log(N. Employees)\\(Mean)}}&\multicolumn{1}{c}{\shortstack{log(Firm age)\\(Mean)}}&\multicolumn{1}{c}{\shortstack{Risk score\\(Mean)}}&\multicolumn{1}{c}{\shortstack{log(Turnover)\\(Mean)}}&\multicolumn{1}{c}{\shortstack{log(N. Employees)\\(Mean)}}&\multicolumn{1}{c}{\shortstack{log(Firm age)\\(Mean)}}&\multicolumn{1}{c}{\shortstack{Risk score\\(Mean)}}&\multicolumn{1}{c}{\shortstack{log(Turnover)\\(Mean)}}&\multicolumn{1}{c}{\shortstack{log(N. Employees)\\(Mean)}}&\multicolumn{1}{c}{\shortstack{log(Firm age)\\(Mean)}}&\multicolumn{1}{c}{\shortstack{Risk score\\(Mean)}}\\
\midrule
Masters/PhD &       -0.06         &        0.14         &        0.03         &        0.05         &       -0.37         &        0.27         &        0.10         &       -0.01         &        0.23         &       -0.11         &       -0.01         &        0.15         \\
            &      (0.29)         &      (0.28)         &      (0.04)         &      (0.08)         &      (0.28)         &      (0.31)         &      (0.09)         &      (0.12)         &      (0.65)         &      (0.45)         &      (0.06)         &      (0.09)         \\
\addlinespace
Above median age&       -0.49         &        0.03         &        0.02         &       -0.02         &       -0.96\sym{*}  &       -0.09         &        0.01         &        0.09         &        0.04         &        0.48         &        0.06         &       -0.08         \\
            &      (0.35)         &      (0.21)         &      (0.04)         &      (0.07)         &      (0.49)         &      (0.32)         &      (0.06)         &      (0.13)         &      (0.49)         &      (0.31)         &      (0.07)         &      (0.09)         \\
\addlinespace
Above median experience&        0.40         &        0.15         &        0.02         &       -0.06         &        0.48         &        0.25         &        0.12\sym{**} &       -0.07         &        0.09         &        0.01         &       -0.07         &       -0.10         \\
            &      (0.41)         &      (0.19)         &      (0.04)         &      (0.06)         &      (0.54)         &      (0.31)         &      (0.06)         &      (0.13)         &      (0.51)         &      (0.28)         &      (0.06)         &      (0.09)         \\
\midrule
N           &       86.00         &       86.00         &       86.00         &       86.00         &       86.00         &       86.00         &       86.00         &       86.00         &       86.00         &       86.00         &       86.00         &       86.00         \\
R2          &        0.63         &        0.61         &        0.05         &        0.78         &        0.43         &        0.50         &        0.11         &        0.70         &        0.66         &        0.42         &        0.10         &        0.54         \\
Mean outcome&       20.07         &        4.43         &        1.09         &       -0.28         &       20.38         &        4.43         &        1.10         &        0.19         &       19.83         &        4.61         &        1.07         &       -0.64         \\
\bottomrule
\multicolumn{13}{l}{\footnotesize Standard errors in parentheses}\\
\multicolumn{13}{l}{\footnotesize \sym{*} \(p<0.10\), \sym{**} \(p<0.05\), \sym{***} \(p<0.01\)}\\
\end{tabular}
}


\begin{tablenotes}[flushleft]
\footnotesize \item Notes: * 0.10 ** 0.05 *** 0.01 significance levels. This table shows the estimation of the OLS regressions of outcomes of inspectors on inspector characteristics. The analysis is done at the list (inspector x year) level and only includes the lists of desk audits. These lists were assigned individually to inspectors. The regressors are three individual characteristics of the inspectors collected using the inspector survey (whether the inspector had a Masters/PhD and whether they had higher than median experience working at the tax authority) and administrative data (age). Column 1's outcome is the total started cases for each list, Columns 2-4 are averages within list of the main outcomes , Column 5 uses the number of started algorithm cases, Column 6 the share of started algorithm cases among algorithm-selected cases, Column 7 uses an indicator of whether the number of started algorithm cases was larger than the number of started inspector cases in the list, and Column 8 uses the response of the inspector to whether they would favor the automation of case selection using an algorithm. The regressions control for fixed effects at the tax office-year level. Robust standard errors are shown in parentheses (Huber-White formula).
\end{tablenotes}
\end{threeparttable}
\end{adjustbox}
\end{table}
\end{landscape}

%%% table for full audits:
\begin{landscape}
\begin{table}[!htbp]\centering
\begin{adjustbox}{max width=1.4\textwidth}
\begin{threeparttable}
\caption{Inspector Characteristics and Performance, Full Audits [Executed]}
\label{tab:inspector_effort_disc_Desk}

{
\def\sym#1{\ifmmode^{#1}\else\(^{#1}\)\fi}
\begin{tabular}{l*{15}{c}}
\toprule
            &\multicolumn{5}{c}{All}                                                                                      &\multicolumn{5}{c}{Algorithm cases}                                                                          &\multicolumn{5}{c}{Inspector cases}                                                                          \\\cmidrule(lr){2-6}\cmidrule(lr){7-11}\cmidrule(lr){12-16}
            &\multicolumn{1}{c}{\shortstack{N Cases\\(Unw. Avg.)}}&\multicolumn{1}{c}{\shortstack{log(Turnover)\\(W. Avg.)}}&\multicolumn{1}{c}{\shortstack{log(Employees)\\(W. Avg.)}}&\multicolumn{1}{c}{\shortstack{log(Firm age)\\(W. Avg.)}}&\multicolumn{1}{c}{\shortstack{Risk score\\(W. Avg.)}}&\multicolumn{1}{c}{\shortstack{N Cases\\(Unw. Avg.)}}&\multicolumn{1}{c}{\shortstack{log(Turnover)\\(W. Avg.)}}&\multicolumn{1}{c}{\shortstack{log(Employees)\\(W. Avg.)}}&\multicolumn{1}{c}{\shortstack{log(Firm age)\\(W. Avg.)}}&\multicolumn{1}{c}{\shortstack{Risk score\\(W. Avg.)}}&\multicolumn{1}{c}{\shortstack{N Cases\\(Unw. Avg.)}}&\multicolumn{1}{c}{\shortstack{log(Turnover)\\(W. Avg.)}}&\multicolumn{1}{c}{\shortstack{log(Employees)\\(W. Avg.)}}&\multicolumn{1}{c}{\shortstack{log(Firm age)\\(W. Avg.)}}&\multicolumn{1}{c}{\shortstack{Risk score\\(W. Avg.)}}\\
\midrule
Masters/PhD &       -0.67         &        0.46         &        0.10         &       -0.12\sym{*}  &        0.04         &        0.37         &        0.08         &       -0.23         &       -0.07         &        0.03         &       -0.37         &        0.41         &        0.32         &       -0.11\sym{*}  &        0.07         \\
            &      (1.51)         &      (0.49)         &      (0.31)         &      (0.07)         &      (0.28)         &      (0.83)         &      (0.51)         &      (0.42)         &      (0.13)         &      (0.28)         &      (0.99)         &      (0.57)         &      (0.35)         &      (0.07)         &      (0.17)         \\
\addlinespace
Above median age&        3.51\sym{***}&        0.57         &        0.06         &        0.03         &        0.05         &        1.71\sym{**} &        0.02         &       -0.02         &       -0.15         &        0.14         &        2.20\sym{***}&        0.66         &        0.01         &        0.07         &        0.08         \\
            &      (1.17)         &      (0.39)         &      (0.22)         &      (0.06)         &      (0.18)         &      (0.66)         &      (0.58)         &      (0.37)         &      (0.12)         &      (0.19)         &      (0.77)         &      (0.47)         &      (0.27)         &      (0.06)         &      (0.16)         \\
\addlinespace
Above median experience&        0.94         &        0.33         &        0.34         &       -0.07         &       -0.00         &        0.42         &        0.40         &        0.31         &       -0.07         &       -0.26         &        0.65         &        0.39         &        0.37         &       -0.04         &       -0.02         \\
            &      (0.80)         &      (0.50)         &      (0.21)         &      (0.05)         &      (0.18)         &      (0.47)         &      (0.57)         &      (0.36)         &      (0.11)         &      (0.21)         &      (0.57)         &      (0.54)         &      (0.24)         &      (0.05)         &      (0.15)         \\
\midrule
N           &      152.00         &      150.00         &      150.00         &      152.00         &      152.00         &       99.00         &       98.00         &       97.00         &       99.00         &       99.00         &      140.00         &      138.00         &      138.00         &      140.00         &      140.00         \\
R2          &        0.26         &        0.28         &        0.42         &        0.12         &        0.46         &        0.28         &        0.28         &        0.48         &        0.15         &        0.64         &        0.24         &        0.27         &        0.25         &        0.11         &        0.53         \\
Mean outcome&        5.53         &       20.84         &        5.11         &        1.12         &        0.21         &        2.71         &       20.78         &        4.61         &        1.03         &        1.92         &        4.09         &       21.01         &        5.35         &        1.14         &       -0.59         \\
\bottomrule
\multicolumn{16}{l}{\footnotesize Standard errors in parentheses}\\
\multicolumn{16}{l}{\footnotesize \sym{*} \(p<0.10\), \sym{**} \(p<0.05\), \sym{***} \(p<0.01\)}\\
\end{tabular}
}


\begin{tablenotes}[flushleft]
\footnotesize \item Notes: * 0.10 ** 0.05 *** 0.01 significance levels. This table shows the estimation of OLS regressions of outcomes on inspector characteristics. The analysis is done at the inspector-year level and focuses on full audits. Unlike desk audits, full audits may involve teams: for team cases, we allocate each executed case equally across participating inspectors, assigning weight $1/k$ to each inspector, where k is the number of inspectors recorded on the case. Column 1's outcome is the team-adjusted number of executed cases for the inspector-year, Column 2 is the team-adjusted number of executed algorithm cases, Column 3 is the share of executed algorithm cases among executed cases (constructed from the team-adjusted counts), Column 4 is an indicator for whether the share of executed algorithm cases exceeds 0.5, Columns 5-6 report averages among executed cases of the main outcomes (detection and log evasion), and Columns 7-10 report averages among executed cases of firm characteristics (log turnover, log number of employees, log firm age, and the risk score). The regressors are individual characteristics of inspectors collected using the inspector survey (e.g., whether the inspector had a Masters/PhD and whether they had higher than median experience working at the tax authority) and administrative data (age). The regressions control for fixed effects at the tax office-year level. Robust standard errors are shown in parentheses (Huber-White formula).
\end{tablenotes}
\end{threeparttable}
\end{adjustbox}
\end{table}
\end{landscape}

\section{Audit Discrepancies and taxpayer dissatisfaction with the audit process.}
\label{sec:audit_discrepancy_taxpayer}
\providecommand{\sym}[1]{\ifmmode^{#1}\else\(^{#1}\)\fi}

\makeatletter
\@ifundefined{landscape}{\newenvironment{landscape}{}{}}{}
\makeatother

This section studies whether larger discrepancies between notified and confirmed audit amounts are associated with taxpayer-reported efficiency and corruption perceptions.

\subsection{Current empirical steps (updated main specification)}
\begin{enumerate}
\item Start from \texttt{datasetforanalysis.dta}.
\item Keep the broad analysis sample by keeping safeties and non-selected, and account for them in the regression controls.
\item Recreate cluster groups used in fixed effects:
\begin{verbatim}
capture drop clusterid
egen clusterid = group(inspectorclusteryear)
\end{verbatim}
\item Recreate Table C5 dispute variables (\texttt{d1}--\texttt{d21}) exactly as in the main analysis script.
\item Construct the baseline discrepancy regressor:
\[
W_{\text{main}}=\frac{\text{notificationvalue}-\text{confirmationvalue}}{\text{notificationvalue}}
\]
\item Cap negative baseline values to zero:
\begin{verbatim}
replace W_main = 0 if W_main < 0 & W_main != .
\end{verbatim}
\item Build robustness discrepancy regressors:
\begin{verbatim}
W_nonzero = 1{W_main != 0}
W_abovemed = 1{W_main > median(W_main)}
W_topquart = 1{W_main >= p75(W_main)}
\end{verbatim}
where \(\text{median}(W_{\text{main}})\) and \(p75(W_{\text{main}})\) are single pooled moments from non-missing \(W_{\text{main}}\) in the regression-eligible sample.
\item In exported regression tables, coefficient labels are:
\textit{Downward revision (\% of notification)},
\textit{Downward revision (\% of notification) $\neq 0$},
\textit{Downward revision (\% of notification) $>$ median}, and
\textit{Downward revision (\% of notification) in top quartile}.
\item Build panel indicators:
\begin{verbatim}
selfreported_audit = 1 if q28a in [2018,2021] or q29a in [2018,2021]
\end{verbatim}
Panel A uses \texttt{q1 $\neq$ . AND selfreported  = 1}; Panel B uses \texttt{q1 $\neq$ .}.
\item Build index outcomes as in taxpayer-survey code:
\begin{verbatim}
index_efficiency  from q31 q33 q41
index_corruption  from q32_inverted q42 q34   (swindex)
\end{verbatim}
\item Estimate Table 6 / D7-style columns by audit type (full, desk, all), using selection method and the following controls controls:
\begin{verbatim}
W_main algorithm overlap random safeties horsprogramme
\end{verbatim}
Index outcomes use \texttt{a(selectionyear center)}; question outcomes (\texttt{q34 q35 q32 q42}) use \texttt{a(inspectorclusteryear)}. The same backbone is used for the three robustness \(W\) variants.
\end{enumerate}

\subsection{Regression equation and notation}

Our empirical exercise studies whether larger discrepancies between the notified and confirmed audit amounts are associated with taxpayer-reported perceptions of the audit process. Concretely, for each taxpayer-survey outcome, we estimate the following model standard OLS regression:
\[
y_{ig}
=
\beta_0
+ \beta_1 W_i
+ \beta_2 \text{Algorithm}_i
+ \beta_3 \text{Overlap}_i
+ \beta_4 \text{Random}_i
+ \beta_5 \text{Safeties}_i
+ \beta_6 \text{Horsprogramme}_i
+ \gamma_g
+ \varepsilon_{ig},
\]
where \(y_{ig}\) is the taxpayer-survey outcome for firm-audit observation \(i\), and \(\gamma_g\) denotes the relevant fixed effects. The unit of observation is a firm-audit observation in the merged administrative audit and taxpayer-survey data, so a given firm may appear more than once if it is associated with multiple audit cases.

Our main discrepancy regressor is
\[
W_i
=
\max\left\{
\frac{\text{Notification}_i-\text{Confirmation}_i}{\text{Notification}_i},
\,0
\right\},
\]
defined for observations with strictly positive notified and confirmed amounts and an executed audit. Thus, \(W_i=0\) both when there is no positive discrepancy and when the confirmed amount is at least as large as the notified amount.

The indicators \(\text{Algorithm}_i\), \(\text{Overlap}_i\), \(\text{Random}_i\), \(\text{Safeties}_i\), and \(\text{Horsprogramme}_i\) capture the audit-case classification. The dummy for random selection is applicable only to desk-audit specifications. The coefficient of interest is \(\beta_1\), which captures the association between the capped discrepancy measure and the taxpayer-reported outcome.

We estimate this equation separately by panel and audit type. Panel A restricts the sample to firms that report having been audited recently,
\[
\texttt{q1 != . \& selfreported\_audit == 1},
\]
while Panel B includes all surveyed firms with non-missing \texttt{q1},
\[
\texttt{q1 != .}.
\]
For audit type, we estimate the specification separately for full audits (\texttt{x2 == 1}), desk audits (\texttt{x2 == 0}), and the pooled sample with no restriction on \texttt{x2}.

The outcomes are the taxpayer-survey indices
\[
\texttt{index\_efficiency}, \quad \texttt{index\_corruption},
\]
and the individual survey questions
\[
\texttt{q34}, \quad \texttt{q35}, \quad \texttt{q32}, \quad \texttt{q42}.
\]
For index outcomes, the fixed effects \(\gamma_g\) correspond to \texttt{selectionyear} \(\times\) \texttt{center}. For question-level outcomes, \(\gamma_g\) corresponds to \texttt{inspectorclusteryear}. We use robust standard errors to perform inference.

In robustness checks, we also replace \(W_i\) with indicator versions:
\[
W_i^{\neq 0}=\mathbf{1}\{W_i\neq 0\},
\qquad
W_i^{>\mathrm{med}}=\mathbf{1}\{W_i>\widetilde{W}\},
\qquad
W_i^{\ge q75}=\mathbf{1}\{W_i\ge\widetilde{W}_{75}\},
\]
where \(\widetilde{W}\) is the pooled median and \(\widetilde{W}_{75}\) is the pooled 75th percentile of non-missing \(W_i\) in the relevant regression sample.

\subsection{Coverage and \(W\) diagnostics}

\begin{table}[!htbp]
\centering
\resizebox{\textwidth}{!}{%
\begin{threeparttable}
\caption{Informative regression-sample coverage by panel and audit type (counts and \% of base)}
\begin{tabular}{l*{6}{c}}

                    &\multicolumn{3}{c}{Panel A: Recent audit}&\multicolumn{3}{c}{Panel B: All surveyed}\\
                    &\multicolumn{1}{c}{Full audits}&\multicolumn{1}{c}{Desk audits}&\multicolumn{1}{c}{All audits}&\multicolumn{1}{c}{Full audits}&\multicolumn{1}{c}{Desk audits}&\multicolumn{1}{c}{All audits}\\
\hline
Eligible surveyed firms (q1 observed)&         227&         413&         640&         262&         501&         763\\
Executed cases (y2=1)&         107&         181&         288&         120&         229&         349\\
Non-missing W       &          54&          62&         116&          59&          76&         135\\
Non-missing W + method controls&          54&          62&         116&          59&          76&         135\\
Regression sample: Efficiency index&          54&          59&         113&          58&          70&         128\\
Regression sample: Corruption index&          53&          59&         112&          57&          71&         128\\
Regression sample: Corruption in General (q34)&          40&          36&          76&          43&          44&          87\\
Regression sample: Corruption Experience (q35)&          54&          62&         116&          59&          75&         134\\
Regression sample: Grade on Honesty (q32)&          49&          52&         101&          52&          59&         111\\
Regression sample: Friend at tax Auth. (q42)&          50&          51&         101&          52&          61&         113\\
\hline

\end{tabular}
\begin{tablenotes}[flushleft]
\footnotesize
\item Notes: Coverage rows report counts and percentages relative to the base surveyed sample in each panel-column cell (base \(=\) firms with \(q1\) observed). Panels are defined as in the regression tables.
\end{tablenotes}
\end{threeparttable}
}
\end{table}

\begin{figure}[!htbp]
\centering
\includegraphics[width=0.82\textwidth]{Output/w_main_scatter_notification.pdf}
\caption{Scatter and linear fit: \(W_{\text{main}}\) vs notification value}
\end{figure}

\begin{figure}[!htbp]
\centering
\includegraphics[width=0.75\textwidth]{Output/w_main_density.pdf}
\caption{Density of \(W_{\text{main}}\)}
\end{figure}

\subsection{Main regression tables (current outputs)}

\begin{table}[!htbp]
\centering
\begin{adjustbox}{max width=1\textwidth}
\begin{threeparttable}
\caption{Selection method, capped \(W\), and survey indexes: two panels (Recent audit / All surveyed firms)}
\begin{tabular}{l*{6}{c}}
\textbf{A: Surveyed Firms (Self-Reporting a Recent Audit)}
                    &\multicolumn{3}{c}{Efficiency Index}                             &\multicolumn{3}{c}{Corruption Index}                             \\
                    &\multicolumn{1}{c}{(1)}&\multicolumn{1}{c}{(2)}&\multicolumn{1}{c}{(3)}&\multicolumn{1}{c}{(4)}&\multicolumn{1}{c}{(5)}&\multicolumn{1}{c}{(6)}\\
                    &\multicolumn{1}{c}{Full audits}&\multicolumn{1}{c}{Desk audits}&\multicolumn{1}{c}{All audits}&\multicolumn{1}{c}{Full audits}&\multicolumn{1}{c}{Desk audits}&\multicolumn{1}{c}{All audits}\\
\hline
W                   &        0.20         &       -0.01         &        0.06         &        0.30         &       -0.00         &       -0.02         \\
                    &      (0.43)         &      (0.05)         &      (0.05)         &      (0.37)         &      (0.04)         &      (0.03)         \\
[1em]
Algorithm           &       -0.53         &        0.44         &       -0.13         &        0.30         &       -0.38         &       -0.03         \\
                    &      (0.35)         &      (0.33)         &      (0.24)         &      (0.28)         &      (0.29)         &      (0.22)         \\
[1em]
Inspectors x Overlap&        0.17         &       -0.76         &       -0.52         &        0.65\sym{*}  &       -0.41         &        0.30         \\
                    &      (0.39)         &      (0.53)         &      (0.46)         &      (0.37)         &      (0.46)         &      (0.30)         \\
[1em]
Algorithm x Random  &                     &       -0.17         &       -0.01         &                     &       -0.15         &       -0.13         \\
                    &                     &      (0.30)         &      (0.30)         &                     &      (0.35)         &      (0.26)         \\
\hline
N                   &          51         &          59         &         113         &          50         &          59         &         112         \\
R2                  &        0.11         &        0.32         &        0.12         &        0.27         &        0.39         &        0.17         \\
Mean outcome        &        0.00         &       -0.06         &       -0.02         &       -0.15         &       -0.10         &       -0.11         \\
\hline
\textbf{B: All Surveyed Firms}
                    &\multicolumn{1}{c}{(1)}         &\multicolumn{1}{c}{(2)}         &\multicolumn{1}{c}{(3)}         &\multicolumn{1}{c}{(4)}         &\multicolumn{1}{c}{(5)}         &\multicolumn{1}{c}{(6)}         \\
\hline
W                   &        0.27         &        0.01         &        0.07         &        0.11         &       -0.01         &       -0.03         \\
                    &      (0.40)         &      (0.05)         &      (0.04)         &      (0.34)         &      (0.04)         &      (0.03)         \\
[1em]
Algorithm           &       -0.54         &        0.19         &       -0.22         &        0.38         &       -0.38         &       -0.01         \\
                    &      (0.34)         &      (0.29)         &      (0.22)         &      (0.27)         &      (0.28)         &      (0.21)         \\
[1em]
Inspectors x Overlap&        0.25         &       -0.92\sym{**} &       -0.48         &        0.52         &       -0.16         &        0.27         \\
                    &      (0.37)         &      (0.40)         &      (0.47)         &      (0.36)         &      (0.43)         &      (0.28)         \\
[1em]
Algorithm x Random  &                     &       -0.16         &       -0.05         &                     &       -0.19         &       -0.11         \\
                    &                     &      (0.28)         &      (0.26)         &                     &      (0.34)         &      (0.25)         \\
\hline
N                   &          55         &          70         &         128         &          54         &          71         &         128         \\
R2                  &        0.14         &        0.28         &        0.14         &        0.26         &        0.25         &        0.12         \\
Mean outcome        &        0.00         &        0.01         &        0.02         &       -0.18         &       -0.06         &       -0.10         \\
\hline

\end{tabular}
\begin{tablenotes}[flushleft]
\footnotesize
\item Notes: OLS results for the Efficiency Index and Corruption Index. The specification is \(y_{ig}=\beta_0+\beta_1 W_i+\beta_2 \text{Algorithm}_i+\beta_3 \text{Overlap}_i+\beta_4 \text{Random}_i+\beta_5 \text{Safeties}_i+\beta_6 \text{Horsprogramme}_i+\gamma_g+\varepsilon_{ig}\). Here \(W_i=W_{\text{main},i}=\max\{(notificationvalue_i-confirmationvalue_i)/notificationvalue_i,0\}\). The control set follows Table D7 (algorithm-selection, overlap, random-selection, replacement, and ad hoc indicators), with the addition of \(W_i\). Panel A includes surveyed firms that report having been audited recently; Panel B includes all surveyed firms. Within each panel, columns report full audits, desk audits, and all audits. Fixed effects are at the selection-year \(\times\) tax-office level. Robust standard errors (Huber-White) are in parentheses. Stars denote significance at \(^{*}p<0.10\), \(^{**}p<0.05\), \(^{***}p<0.01\).
\end{tablenotes}
\end{threeparttable}
\end{adjustbox}
\end{table}

\begin{landscape}
\begin{table}[!htbp]
\centering
\begin{adjustbox}{max width=1.25\textwidth}
\begin{threeparttable}
\caption{Selection method, capped \(W\), and survey-question outcomes: two panels (Recent audit / All surveyed firms)}
\begin{tabular}{l*{12}{c}}
\textbf{A: Surveyed Firms (Self-Reporting a Recent Audit)}
                    &\multicolumn{3}{c}{Corruption in General}                        &\multicolumn{3}{c}{Corruption Experience}                        &\multicolumn{3}{c}{Grade on Honesty}                             &\multicolumn{3}{c}{Friend at tax Auth.}                          \\
                    &\multicolumn{1}{c}{(1)}&\multicolumn{1}{c}{(2)}&\multicolumn{1}{c}{(3)}&\multicolumn{1}{c}{(4)}&\multicolumn{1}{c}{(5)}&\multicolumn{1}{c}{(6)}&\multicolumn{1}{c}{(7)}&\multicolumn{1}{c}{(8)}&\multicolumn{1}{c}{(9)}&\multicolumn{1}{c}{(10)}&\multicolumn{1}{c}{(11)}&\multicolumn{1}{c}{(12)}\\
                    &\multicolumn{1}{c}{Full audits}&\multicolumn{1}{c}{Desk audits}&\multicolumn{1}{c}{All audits}&\multicolumn{1}{c}{Full audits}&\multicolumn{1}{c}{Desk audits}&\multicolumn{1}{c}{All audits}&\multicolumn{1}{c}{Full audits}&\multicolumn{1}{c}{Desk audits}&\multicolumn{1}{c}{All audits}&\multicolumn{1}{c}{Full audits}&\multicolumn{1}{c}{Desk audits}&\multicolumn{1}{c}{All audits}\\
\hline
W                   &      18.646         &      -0.120         &      -1.620         &      -0.131         &                     &      -0.007         &      -1.085         &       0.141         &       0.047         &      -0.080         &       0.162         &       0.142         \\
                    &    (11.367)         &     (3.229)         &     (1.690)         &     (0.172)         &                     &     (0.011)         &     (0.984)         &     (0.191)         &     (0.196)         &     (0.739)         &     (0.120)         &     (0.101)         \\
[1em]
Algorithm           &      -4.046         &     -44.528         &     -10.484         &       0.031         &                     &       0.025         &      -0.885         &      -2.254         &      -1.152         &       0.198         &      -0.139         &       0.147         \\
                    &     (6.547)         &    (38.381)         &     (7.186)         &     (0.090)         &                     &     (0.080)         &     (0.886)         &     (1.440)         &     (0.817)         &     (0.509)         &     (1.156)         &     (0.484)         \\
[1em]
Inspectors x Overlap&      33.691\sym{***}&                     &      36.936\sym{***}&       0.052         &                     &       0.007         &      -0.648         &                     &      -1.052         &       0.097         &                     &      -0.007         \\
                    &     (4.713)         &                     &     (4.160)         &     (0.064)         &                     &     (0.018)         &     (1.616)         &                     &     (1.719)         &     (0.600)         &                     &     (0.561)         \\
[1em]
Algorithm x Random  &                     &                     &                     &                     &                     &      -0.022         &                     &                     &                     &                     &      -0.918         &      -1.198         \\
                    &                     &                     &                     &                     &                     &     (0.079)         &                     &                     &                     &                     &     (1.423)         &     (0.806)         \\
\hline
N                   &          38         &          10         &          48         &          50         &          34         &          84         &          46         &          24         &          70         &          45         &          25         &          70         \\
R2                  &       0.294         &       0.638         &       0.310         &       0.203         &           .         &       0.207         &       0.195         &       0.478         &       0.248         &       0.156         &       0.593         &       0.318         \\
Mean outcome        &       11.81         &       20.50         &       13.62         &        0.08         &        0.00         &        0.04         &        6.78         &        6.45         &        6.67         &        3.40         &        3.48         &        3.42         \\
\hline
\textbf{B: All Surveyed Firms}
                    &\multicolumn{1}{c}{(1)}         &\multicolumn{1}{c}{(2)}         &\multicolumn{1}{c}{(3)}         &\multicolumn{1}{c}{(4)}         &\multicolumn{1}{c}{(5)}         &\multicolumn{1}{c}{(6)}         &\multicolumn{1}{c}{(7)}         &\multicolumn{1}{c}{(8)}         &\multicolumn{1}{c}{(9)}         &\multicolumn{1}{c}{(10)}         &\multicolumn{1}{c}{(11)}         &\multicolumn{1}{c}{(12)}         \\
\hline
W                   &      10.735         &      -2.130         &      -2.283\sym{*}  &      -0.213         &                     &      -0.011         &      -0.959         &       0.097         &       0.040         &      -0.215         &       0.130         &       0.121         \\
                    &    (11.482)         &     (1.904)         &     (1.314)         &     (0.174)         &                     &     (0.013)         &     (0.957)         &     (0.226)         &     (0.200)         &     (0.696)         &     (0.127)         &     (0.101)         \\
[1em]
Algorithm           &      -5.625         &     -22.938         &     -10.122         &       0.024         &                     &       0.020         &      -0.884         &      -1.065         &      -0.944         &       0.249         &       0.515         &       0.340         \\
                    &     (6.254)         &    (21.168)         &     (6.261)         &     (0.087)         &                     &     (0.068)         &     (0.863)         &     (1.034)         &     (0.731)         &     (0.499)         &     (0.895)         &     (0.444)         \\
[1em]
Inspectors x Overlap&      19.048         &                     &      23.436         &      -0.068         &                     &      -0.158         &      -0.176         &                     &      -0.606         &       0.174         &                     &       0.022         \\
                    &    (16.079)         &                     &    (15.200)         &     (0.149)         &                     &     (0.186)         &     (1.424)         &                     &     (1.412)         &     (0.541)         &                     &     (0.384)         \\
[1em]
Algorithm x Random  &                     &                     &                     &                     &                     &      -0.021         &                     &      -1.763\sym{*}  &      -1.901\sym{**} &                     &      -1.561         &      -1.383\sym{*}  \\
                    &                     &                     &                     &                     &                     &     (0.046)         &                     &     (0.956)         &     (0.918)         &                     &     (1.175)         &     (0.774)         \\
\hline
N                   &          41         &          16         &          57         &          55         &          46         &         101         &          49         &          31         &          80         &          47         &          32         &          79         \\
R2                  &       0.185         &       0.530         &       0.285         &       0.200         &           .         &       0.181         &       0.215         &       0.462         &       0.272         &       0.153         &       0.403         &       0.276         \\
Mean outcome        &       12.78         &       16.87         &       13.92         &        0.09         &        0.00         &        0.04         &        6.77         &        6.51         &        6.67         &        3.36         &        3.75         &        3.51         \\
\hline

\end{tabular}
\begin{tablenotes}[flushleft]
\footnotesize
\item Notes: OLS results for four taxpayer-survey outcomes: Corruption in general (Columns 1 to 3), Corruption experience (Columns 4 to 6), Grade on honesty (Columns 7 to 9), and Friend at tax authority (Columns 10 to 12). The specification is \(y_{ig}=\beta_0+\beta_1 W_i+\beta_2 \text{Algorithm}_i+\beta_3 \text{Overlap}_i+\beta_4 \text{Random}_i+\beta_5 \text{Safeties}_i+\beta_6 \text{Horsprogramme}_i+\gamma_g+\varepsilon_{ig}\). Here \(W_i=W_{\text{main},i}=\max\{(notificationvalue_i-confirmationvalue_i)/notificationvalue_i,0\}\). The control set follows Table D7 (algorithm-selection, overlap, random-selection, replacement, and ad hoc indicators), with the addition of \(W_i\). Panel A includes surveyed firms that report having been audited recently; Panel B includes all surveyed firms. Within each panel, columns report full audits, desk audits, and all audits. Corruption in general is the declared belief about the percentage of audits in Senegal affected by corruption; Corruption experience indicates that the respondent declared an instance of bribery to obtain tax favors; Grade on honesty is a 0-to-10 grade for the honesty of the inspector with whom the respondent last interacted during a full audit; and Friend at tax authority captures agreement with the statement that if the firm's owner has a friend at the tax authority, the firm will rarely be audited. Fixed effects are at the inspector-cluster-year level. Robust standard errors (Huber-White) are in parentheses. Stars denote significance at \(^{*}p<0.10\), \(^{**}p<0.05\), \(^{***}p<0.01\).
\end{tablenotes}
\end{threeparttable}
\end{adjustbox}
\end{table}
\end{landscape}

\subsection{Robustness regression tables}

\begin{table}[!htbp]
\centering
\resizebox{\textwidth}{!}{%
\begin{threeparttable}
\caption{Robustness: \(W_{\neq 0}\) and survey indexes (Recent audit / All surveyed firms)}
\begin{tabular}{l*{6}{c}}
\textbf{A: Surveyed Firms (Self-Reporting a Recent Audit)}
                    &\multicolumn{3}{c}{Efficiency Index}                             &\multicolumn{3}{c}{Corruption Index}                             \\
                    &\multicolumn{1}{c}{(1)}&\multicolumn{1}{c}{(2)}&\multicolumn{1}{c}{(3)}&\multicolumn{1}{c}{(4)}&\multicolumn{1}{c}{(5)}&\multicolumn{1}{c}{(6)}\\
                    &\multicolumn{1}{c}{Full audits}&\multicolumn{1}{c}{Desk audits}&\multicolumn{1}{c}{All audits}&\multicolumn{1}{c}{Full audits}&\multicolumn{1}{c}{Desk audits}&\multicolumn{1}{c}{All audits}\\
\hline
Downward revision (\% of notification) $\neq$ 0&        0.43         &       -0.07         &        0.21         &        0.29         &        0.21         &        0.15         \\
                    &      (0.41)         &      (0.27)         &      (0.22)         &      (0.33)         &      (0.22)         &      (0.17)         \\
[1em]
Algorithm           &       -0.52         &        0.43         &       -0.12         &        0.31         &       -0.42         &       -0.03         \\
                    &      (0.36)         &      (0.32)         &      (0.24)         &      (0.28)         &      (0.29)         &      (0.21)         \\
[1em]
Inspectors x Overlap&        0.29         &       -1.13\sym{*}  &       -0.68         &        0.68\sym{*}  &       -0.36         &        0.31         \\
                    &      (0.38)         &      (0.63)         &      (0.54)         &      (0.36)         &      (0.35)         &      (0.24)         \\
[1em]
Algorithm x Random  &                     &       -0.01         &        0.04         &                     &       -0.14         &       -0.10         \\
                    &                     &      (0.31)         &      (0.30)         &                     &      (0.32)         &      (0.25)         \\
\hline
N                   &          55         &          81         &         138         &          54         &          81         &         137         \\
R2                  &        0.15         &        0.25         &        0.09         &        0.28         &        0.33         &        0.19         \\
Mean outcome        &       -0.01         &       -0.01         &        0.00         &       -0.13         &       -0.10         &       -0.11         \\
\hline
\textbf{B: All Surveyed Firms}
                    &\multicolumn{1}{c}{(1)}         &\multicolumn{1}{c}{(2)}         &\multicolumn{1}{c}{(3)}         &\multicolumn{1}{c}{(4)}         &\multicolumn{1}{c}{(5)}         &\multicolumn{1}{c}{(6)}         \\
\hline
Downward revision (\% of notification) $\neq$ 0&        0.46         &       -0.08         &        0.18         &        0.08         &        0.21         &        0.09         \\
                    &      (0.35)         &      (0.23)         &      (0.19)         &      (0.31)         &      (0.21)         &      (0.16)         \\
[1em]
Algorithm           &       -0.52         &        0.17         &       -0.20         &        0.37         &       -0.38         &       -0.03         \\
                    &      (0.34)         &      (0.28)         &      (0.22)         &      (0.27)         &      (0.29)         &      (0.21)         \\
[1em]
Inspectors x Overlap&        0.30         &       -1.44\sym{**} &       -0.76         &        0.58         &       -0.17         &        0.41\sym{*}  \\
                    &      (0.33)         &      (0.60)         &      (0.54)         &      (0.35)         &      (0.37)         &      (0.22)         \\
[1em]
Algorithm x Random  &                     &        0.04         &        0.05         &                     &       -0.20         &       -0.16         \\
                    &                     &      (0.29)         &      (0.27)         &                     &      (0.32)         &      (0.25)         \\
\hline
N                   &          61         &          96         &         159         &          60         &          97         &         159         \\
R2                  &        0.18         &        0.20         &        0.09         &        0.30         &        0.22         &        0.12         \\
Mean outcome        &        0.01         &        0.03         &        0.03         &       -0.22         &       -0.05         &       -0.11         \\
\hline

\end{tabular}
\begin{tablenotes}[flushleft]
\footnotesize
\item Notes: OLS results for the Efficiency Index and Corruption Index. The specification is \(y_{ig}=\beta_0+\beta_1 W_i^{\neq 0}+\beta_2 \text{Algorithm}_i+\beta_3 \text{Overlap}_i+\beta_4 \text{Random}_i+\beta_5 \text{Safeties}_i+\beta_6 \text{Horsprogramme}_i+\gamma_g+\varepsilon_{ig}\). Here \(W_i^{\neq 0}=\mathbf{1}\{W_{\text{main},i}\neq 0\}\). The control set follows Table D7 (algorithm-selection, overlap, random-selection, replacement, and ad hoc indicators), with the addition of the discrepancy term \(W_i^{\neq 0}\). Panel A includes surveyed firms that report having been audited recently; Panel B includes all surveyed firms. Within each panel, columns report full audits, desk audits, and all audits. Fixed effects are at the selection-year \(\times\) tax-office level. Robust standard errors (Huber-White) are in parentheses. Stars denote significance at \(^{*}p<0.10\), \(^{**}p<0.05\), \(^{***}p<0.01\).
\end{tablenotes}
\end{threeparttable}
}
\end{table}

\begin{landscape}
\begin{table}[!htbp]
\centering
\begin{adjustbox}{max width=1.25\textwidth}
\begin{threeparttable}
\caption{Robustness: \(W_{\neq 0}\) and survey questions (Recent audit / All surveyed firms)}
\begin{tabular}{l*{12}{c}}
\textbf{A: Surveyed Firms (Self-Reporting a Recent Audit)}
                    &\multicolumn{3}{c}{Corruption in General}                        &\multicolumn{3}{c}{Corruption Experience}                        &\multicolumn{3}{c}{Grade on Honesty}                             &\multicolumn{3}{c}{Friend at tax Auth.}                          \\
                    &\multicolumn{1}{c}{(1)}&\multicolumn{1}{c}{(2)}&\multicolumn{1}{c}{(3)}&\multicolumn{1}{c}{(4)}&\multicolumn{1}{c}{(5)}&\multicolumn{1}{c}{(6)}&\multicolumn{1}{c}{(7)}&\multicolumn{1}{c}{(8)}&\multicolumn{1}{c}{(9)}&\multicolumn{1}{c}{(10)}&\multicolumn{1}{c}{(11)}&\multicolumn{1}{c}{(12)}\\
                    &\multicolumn{1}{c}{Full audits}&\multicolumn{1}{c}{Desk audits}&\multicolumn{1}{c}{All audits}&\multicolumn{1}{c}{Full audits}&\multicolumn{1}{c}{Desk audits}&\multicolumn{1}{c}{All audits}&\multicolumn{1}{c}{Full audits}&\multicolumn{1}{c}{Desk audits}&\multicolumn{1}{c}{All audits}&\multicolumn{1}{c}{Full audits}&\multicolumn{1}{c}{Desk audits}&\multicolumn{1}{c}{All audits}\\
\hline
Downward revision (\% of notification) $\neq$ 0&      17.226\sym{*}  &       3.684         &      10.975         &      -0.014         &       0.099         &       0.038         &       0.419         &       0.044         &       0.281         &       0.655         &       0.656         &       0.629\sym{**} \\
                    &     (9.243)         &    (10.664)         &     (7.172)         &     (0.075)         &     (0.069)         &     (0.050)         &     (0.768)         &     (0.631)         &     (0.532)         &     (0.458)         &     (0.422)         &     (0.315)         \\
[1em]
Algorithm           &      -4.663         &     -28.728         &      -8.520         &       0.025         &      -0.062         &       0.005         &      -0.907         &      -1.165         &      -0.942         &       0.200         &       0.509         &       0.245         \\
                    &     (6.720)         &    (21.678)         &     (7.599)         &     (0.089)         &     (0.071)         &     (0.076)         &     (0.873)         &     (1.234)         &     (0.774)         &     (0.481)         &     (0.782)         &     (0.456)         \\
[1em]
Inspectors x Overlap&      38.446\sym{***}&                     &      37.160\sym{***}&       0.005         &                     &       0.001         &      -0.977         &                     &      -0.986         &       0.067         &                     &       0.082         \\
                    &     (2.649)         &                     &     (3.725)         &     (0.018)         &                     &     (0.015)         &     (1.569)         &                     &     (1.673)         &     (0.499)         &                     &     (0.528)         \\
[1em]
Algorithm x Random  &                     &                     &                     &                     &      -0.287         &      -0.273         &                     &       0.111         &      -1.226         &                     &      -1.650\sym{*}  &      -0.835         \\
                    &                     &                     &                     &                     &     (0.223)         &     (0.195)         &                     &     (1.806)         &     (1.010)         &                     &     (0.950)         &     (0.691)         \\
\hline
N                   &          49         &          45         &          94         &          65         &          89         &         154         &          61         &          67         &         128         &          57         &          76         &         133         \\
R2                  &       0.319         &       0.429         &       0.319         &       0.136         &       0.395         &       0.240         &       0.141         &       0.537         &       0.325         &       0.108         &       0.418         &       0.259         \\
Mean outcome        &       12.63         &       17.33         &       14.88         &        0.07         &        0.04         &        0.05         &        6.68         &        6.44         &        6.56         &        3.45         &        3.35         &        3.39         \\
\hline
\textbf{B: All Surveyed Firms}
                    &\multicolumn{1}{c}{(1)}         &\multicolumn{1}{c}{(2)}         &\multicolumn{1}{c}{(3)}         &\multicolumn{1}{c}{(4)}         &\multicolumn{1}{c}{(5)}         &\multicolumn{1}{c}{(6)}         &\multicolumn{1}{c}{(7)}         &\multicolumn{1}{c}{(8)}         &\multicolumn{1}{c}{(9)}         &\multicolumn{1}{c}{(10)}         &\multicolumn{1}{c}{(11)}         &\multicolumn{1}{c}{(12)}         \\
\hline
Downward revision (\% of notification) $\neq$ 0&       9.151         &      -1.612         &       4.676         &      -0.116         &       0.016         &      -0.049         &       0.228         &       0.047         &       0.192         &       0.527         &       0.498         &       0.505\sym{*}  \\
                    &     (9.581)         &    (10.466)         &     (7.148)         &     (0.094)         &     (0.072)         &     (0.059)         &     (0.692)         &     (0.578)         &     (0.482)         &     (0.424)         &     (0.362)         &     (0.280)         \\
[1em]
Algorithm           &      -6.042         &     -20.476         &      -8.787         &       0.015         &       0.016         &       0.024         &      -0.874         &      -0.380         &      -0.747         &       0.250         &       0.402         &       0.301         \\
                    &     (6.344)         &    (14.568)         &     (6.301)         &     (0.088)         &     (0.069)         &     (0.063)         &     (0.848)         &     (0.866)         &     (0.681)         &     (0.475)         &     (0.686)         &     (0.410)         \\
[1em]
Inspectors x Overlap&      21.202         &                     &      21.634         &      -0.145         &                     &      -0.154         &      -0.620         &                     &      -0.588         &      -0.069         &                     &      -0.051         \\
                    &    (16.389)         &                     &    (16.345)         &     (0.144)         &                     &     (0.172)         &     (1.305)         &                     &     (1.366)         &     (0.425)         &                     &     (0.423)         \\
[1em]
Algorithm x Random  &                     &                     &                     &                     &      -0.303         &      -0.277\sym{*}  &                     &      -1.363         &      -1.570\sym{**} &                     &      -1.041         &      -0.754         \\
                    &                     &                     &                     &                     &     (0.185)         &     (0.165)         &                     &     (0.855)         &     (0.661)         &                     &     (0.947)         &     (0.628)         \\
\hline
N                   &          54         &          54         &         108         &          73         &         111         &         184         &          67         &          78         &         145         &          61         &          88         &         149         \\
R2                  &       0.193         &       0.407         &       0.285         &       0.125         &       0.407         &       0.261         &       0.145         &       0.527         &       0.336         &       0.102         &       0.343         &       0.238         \\
Mean outcome        &       12.85         &       15.64         &       14.25         &        0.09         &        0.06         &        0.07         &        6.74         &        6.50         &        6.61         &        3.39         &        3.40         &        3.40         \\
\hline

\end{tabular}
\begin{tablenotes}[flushleft]
\footnotesize
\item Notes: OLS results for four taxpayer-survey outcomes: Corruption in general (Columns 1 to 3), Corruption experience (Columns 4 to 6), Grade on honesty (Columns 7 to 9), and Friend at tax authority (Columns 10 to 12). The specification is \(y_{ig}=\beta_0+\beta_1 W_i^{\neq 0}+\beta_2 \text{Algorithm}_i+\beta_3 \text{Overlap}_i+\beta_4 \text{Random}_i+\beta_5 \text{Safeties}_i+\beta_6 \text{Horsprogramme}_i+\gamma_g+\varepsilon_{ig}\). Here \(W_i^{\neq 0}=\mathbf{1}\{W_{\text{main},i}\neq 0\}\). The control set follows Table D7 (algorithm-selection, overlap, random-selection, replacement, and ad hoc indicators), with the addition of the discrepancy term \(W_i^{\neq 0}\). Panel A includes surveyed firms that report having been audited recently; Panel B includes all surveyed firms. Within each panel, columns report full audits, desk audits, and all audits. Corruption in general is the declared belief about the percentage of audits in Senegal affected by corruption; Corruption experience indicates that the respondent declared an instance of bribery to obtain tax favors; Grade on honesty is a 0-to-10 grade for the honesty of the inspector with whom the respondent last interacted during a full audit; and Friend at tax authority captures agreement with the statement that if the firm's owner has a friend at the tax authority, the firm will rarely be audited. Fixed effects are at the inspector-cluster-year level. Robust standard errors (Huber-White) are in parentheses. Stars denote significance at \(^{*}p<0.10\), \(^{**}p<0.05\), \(^{***}p<0.01\).
\end{tablenotes}
\end{threeparttable}
\end{adjustbox}
\end{table}
\end{landscape}

\begin{table}[!htbp]
\centering
\begin{adjustbox}{max width=1\textwidth}
\begin{threeparttable}
\caption{Robustness: \(W_{>\mathrm{med}}\) and survey indexes (Recent audit / All surveyed firms)}
\begin{tabular}{l*{6}{c}}
\textbf{A: Surveyed Firms (Self-Reporting a Recent Audit)}
                    &\multicolumn{3}{c}{Efficiency Index}                             &\multicolumn{3}{c}{Corruption Index}                             \\
                    &\multicolumn{1}{c}{(1)}&\multicolumn{1}{c}{(2)}&\multicolumn{1}{c}{(3)}&\multicolumn{1}{c}{(4)}&\multicolumn{1}{c}{(5)}&\multicolumn{1}{c}{(6)}\\
                    &\multicolumn{1}{c}{Full audits}&\multicolumn{1}{c}{Desk audits}&\multicolumn{1}{c}{All audits}&\multicolumn{1}{c}{Full audits}&\multicolumn{1}{c}{Desk audits}&\multicolumn{1}{c}{All audits}\\
\hline
Downward revision (\% of notification) $>$ median&        0.44         &        0.06         &        0.29         &        0.07         &       -0.02         &       -0.07         \\
                    &      (0.36)         &      (0.24)         &      (0.19)         &      (0.32)         &      (0.23)         &      (0.17)         \\
[1em]
Algorithm           &       -0.51         &        0.43         &       -0.11         &        0.30         &       -0.41         &       -0.04         \\
                    &      (0.36)         &      (0.31)         &      (0.24)         &      (0.28)         &      (0.28)         &      (0.21)         \\
[1em]
Inspectors x Overlap&        0.22         &       -1.13\sym{*}  &       -0.72         &        0.66\sym{*}  &       -0.39         &        0.27         \\
                    &      (0.41)         &      (0.60)         &      (0.52)         &      (0.36)         &      (0.42)         &      (0.26)         \\
[1em]
Algorithm x Random  &                     &       -0.04         &       -0.00         &                     &       -0.10         &       -0.09         \\
                    &                     &      (0.33)         &      (0.31)         &                     &      (0.33)         &      (0.25)         \\
\hline
N                   &          55         &          81         &         138         &          54         &          81         &         137         \\
R2                  &        0.16         &        0.25         &        0.10         &        0.27         &        0.33         &        0.18         \\
Mean outcome        &       -0.01         &       -0.01         &        0.00         &       -0.13         &       -0.10         &       -0.11         \\
\hline
\textbf{B: All Surveyed Firms}
                    &\multicolumn{1}{c}{(1)}         &\multicolumn{1}{c}{(2)}         &\multicolumn{1}{c}{(3)}         &\multicolumn{1}{c}{(4)}         &\multicolumn{1}{c}{(5)}         &\multicolumn{1}{c}{(6)}         \\
\hline
Downward revision (\% of notification) $>$ median&        0.51\sym{*}  &        0.16         &        0.34\sym{**} &       -0.12         &        0.05         &       -0.09         \\
                    &      (0.30)         &      (0.21)         &      (0.17)         &      (0.28)         &      (0.21)         &      (0.16)         \\
[1em]
Algorithm           &       -0.50         &        0.18         &       -0.18         &        0.36         &       -0.38         &       -0.04         \\
                    &      (0.35)         &      (0.28)         &      (0.22)         &      (0.28)         &      (0.28)         &      (0.21)         \\
[1em]
Inspectors x Overlap&        0.25         &       -1.44\sym{***}&       -0.79         &        0.59\sym{*}  &       -0.22         &        0.39         \\
                    &      (0.36)         &      (0.55)         &      (0.51)         &      (0.34)         &      (0.41)         &      (0.24)         \\
[1em]
Algorithm x Random  &                     &       -0.01         &        0.02         &                     &       -0.18         &       -0.14         \\
                    &                     &      (0.31)         &      (0.28)         &                     &      (0.32)         &      (0.25)         \\
\hline
N                   &          61         &          96         &         159         &          60         &          97         &         159         \\
R2                  &        0.20         &        0.20         &        0.11         &        0.30         &        0.21         &        0.12         \\
Mean outcome        &        0.01         &        0.03         &        0.03         &       -0.22         &       -0.05         &       -0.11         \\
\hline

\end{tabular}
\begin{tablenotes}[flushleft]
\footnotesize
\item Notes: OLS results for the Efficiency Index and Corruption Index. The specification is \(y_{ig}=\beta_0+\beta_1 W_i^{>\mathrm{med}}+\beta_2 \text{Algorithm}_i+\beta_3 \text{Overlap}_i+\beta_4 \text{Random}_i+\beta_5 \text{Safeties}_i+\beta_6 \text{Horsprogramme}_i+\gamma_g+\varepsilon_{ig}\). Here \(W_i^{>\mathrm{med}}=\mathbf{1}\{W_{\text{main},i}>\widetilde{W}\}\), where \(\widetilde{W}\) is the pooled median of non-missing \(W_{\text{main}}\) in the regression-eligible sample. The control set follows Table D7 (algorithm-selection, overlap, random-selection, replacement, and ad hoc indicators), with the addition of the discrepancy term \(W_i^{>\mathrm{med}}\). Panel A includes surveyed firms that report having been audited recently; Panel B includes all surveyed firms. Within each panel, columns report full audits, desk audits, and all audits. Fixed effects are at the selection-year \(\times\) tax-office level. Robust standard errors (Huber-White) are in parentheses. Stars denote significance at \(^{*}p<0.10\), \(^{**}p<0.05\), \(^{***}p<0.01\).
\end{tablenotes}
\end{threeparttable}
\end{adjustbox}
\end{table}

\begin{landscape}
\begin{table}[!htbp]
\centering
\begin{adjustbox}{max width=1.25\textwidth}
\begin{threeparttable}
\caption{Robustness: \(W_{>\mathrm{med}}\) and survey questions (Recent audit / All surveyed firms)}
\begin{tabular}{l*{12}{c}}
\textbf{A: Surveyed Firms (Self-Reporting a Recent Audit)}
                    &\multicolumn{3}{c}{Corruption in General}                        &\multicolumn{3}{c}{Corruption Experience}                        &\multicolumn{3}{c}{Grade on Honesty}                             &\multicolumn{3}{c}{Friend at tax Auth.}                          \\
                    &\multicolumn{1}{c}{(1)}&\multicolumn{1}{c}{(2)}&\multicolumn{1}{c}{(3)}&\multicolumn{1}{c}{(4)}&\multicolumn{1}{c}{(5)}&\multicolumn{1}{c}{(6)}&\multicolumn{1}{c}{(7)}&\multicolumn{1}{c}{(8)}&\multicolumn{1}{c}{(9)}&\multicolumn{1}{c}{(10)}&\multicolumn{1}{c}{(11)}&\multicolumn{1}{c}{(12)}\\
                    &\multicolumn{1}{c}{Full audits}&\multicolumn{1}{c}{Desk audits}&\multicolumn{1}{c}{All audits}&\multicolumn{1}{c}{Full audits}&\multicolumn{1}{c}{Desk audits}&\multicolumn{1}{c}{All audits}&\multicolumn{1}{c}{Full audits}&\multicolumn{1}{c}{Desk audits}&\multicolumn{1}{c}{All audits}&\multicolumn{1}{c}{Full audits}&\multicolumn{1}{c}{Desk audits}&\multicolumn{1}{c}{All audits}\\
\hline
Downward revision (\% of notification) $>$ median&      12.917\sym{*}  &       5.271         &       8.785         &      -0.017         &       0.004         &      -0.017         &      -0.026         &      -0.276         &      -0.084         &       0.452         &       0.320         &       0.324         \\
                    &     (7.228)         &    (11.889)         &     (6.359)         &     (0.077)         &     (0.072)         &     (0.047)         &     (0.613)         &     (0.609)         &     (0.465)         &     (0.435)         &     (0.487)         &     (0.315)         \\
[1em]
Algorithm           &      -3.681         &     -28.964         &      -7.639         &       0.024         &      -0.042         &       0.008         &      -0.886         &      -1.062         &      -0.908         &       0.230         &       0.569         &       0.266         \\
                    &     (6.664)         &    (21.789)         &     (7.663)         &     (0.089)         &     (0.072)         &     (0.075)         &     (0.859)         &     (1.250)         &     (0.767)         &     (0.490)         &     (0.898)         &     (0.461)         \\
[1em]
Inspectors x Overlap&      38.773\sym{***}&                     &      37.454\sym{***}&       0.008         &                     &       0.005         &      -0.972         &                     &      -0.977         &       0.077         &                     &       0.089         \\
                    &     (2.486)         &                     &     (3.538)         &     (0.020)         &                     &     (0.017)         &     (1.567)         &                     &     (1.669)         &     (0.493)         &                     &     (0.523)         \\
[1em]
Algorithm x Random  &                     &                     &                     &                     &      -0.274         &      -0.262         &                     &       0.071         &      -1.288         &                     &      -1.587         &      -0.735         \\
                    &                     &                     &                     &                     &     (0.227)         &     (0.199)         &                     &     (1.845)         &     (1.022)         &                     &     (1.016)         &     (0.671)         \\
\hline
N                   &          49         &          45         &          94         &          65         &          89         &         154         &          61         &          67         &         128         &          57         &          76         &         133         \\
R2                  &       0.301         &       0.431         &       0.313         &       0.136         &       0.367         &       0.237         &       0.135         &       0.540         &       0.323         &       0.098         &       0.393         &       0.239         \\
Mean outcome        &       12.63         &       17.33         &       14.88         &        0.07         &        0.04         &        0.05         &        6.68         &        6.44         &        6.56         &        3.45         &        3.35         &        3.39         \\
\hline
\textbf{B: All Surveyed Firms}
                    &\multicolumn{1}{c}{(1)}         &\multicolumn{1}{c}{(2)}         &\multicolumn{1}{c}{(3)}         &\multicolumn{1}{c}{(4)}         &\multicolumn{1}{c}{(5)}         &\multicolumn{1}{c}{(6)}         &\multicolumn{1}{c}{(7)}         &\multicolumn{1}{c}{(8)}         &\multicolumn{1}{c}{(9)}         &\multicolumn{1}{c}{(10)}         &\multicolumn{1}{c}{(11)}         &\multicolumn{1}{c}{(12)}         \\
\hline
Downward revision (\% of notification) $>$ median&       8.636         &      -2.923         &       4.235         &      -0.063         &      -0.083         &      -0.076         &      -0.025         &      -0.145         &      -0.041         &       0.354         &       0.136         &       0.228         \\
                    &     (7.434)         &    (11.470)         &     (6.259)         &     (0.083)         &     (0.064)         &     (0.049)         &     (0.561)         &     (0.556)         &     (0.414)         &     (0.387)         &     (0.389)         &     (0.274)         \\
[1em]
Algorithm           &      -5.573         &     -20.132         &      -8.615         &       0.019         &       0.031         &       0.023         &      -0.864         &      -0.344         &      -0.731         &       0.272         &       0.462         &       0.323         \\
                    &     (6.444)         &    (14.723)         &     (6.352)         &     (0.088)         &     (0.071)         &     (0.063)         &     (0.837)         &     (0.854)         &     (0.674)         &     (0.480)         &     (0.713)         &     (0.414)         \\
[1em]
Inspectors x Overlap&      21.448         &                     &      21.788         &      -0.142         &                     &      -0.138         &      -0.568         &                     &      -0.538         &      -0.021         &                     &       0.024         \\
                    &    (16.514)         &                     &    (16.340)         &     (0.153)         &                     &     (0.169)         &     (1.308)         &                     &     (1.373)         &     (0.394)         &                     &     (0.393)         \\
[1em]
Algorithm x Random  &                     &                     &                     &                     &      -0.295\sym{*}  &      -0.278\sym{*}  &                     &      -1.420         &      -1.643\sym{**} &                     &      -0.940         &      -0.654         \\
                    &                     &                     &                     &                     &     (0.175)         &     (0.165)         &                     &     (0.908)         &     (0.691)         &                     &     (0.945)         &     (0.605)         \\
\hline
N                   &          54         &          54         &         108         &          73         &         111         &         184         &          67         &          78         &         145         &          61         &          88         &         149         \\
R2                  &       0.199         &       0.408         &       0.285         &       0.108         &       0.423         &       0.270         &       0.143         &       0.528         &       0.335         &       0.093         &       0.325         &       0.223         \\
Mean outcome        &       12.85         &       15.64         &       14.25         &        0.09         &        0.06         &        0.07         &        6.74         &        6.50         &        6.61         &        3.39         &        3.40         &        3.40         \\
\hline

\end{tabular}
\begin{tablenotes}[flushleft]
\footnotesize
\item Notes: OLS results for four taxpayer-survey outcomes: Corruption in general (Columns 1 to 3), Corruption experience (Columns 4 to 6), Grade on honesty (Columns 7 to 9), and Friend at tax authority (Columns 10 to 12). The specification is \(y_{ig}=\beta_0+\beta_1 W_i^{>\mathrm{med}}+\beta_2 \text{Algorithm}_i+\beta_3 \text{Overlap}_i+\beta_4 \text{Random}_i+\beta_5 \text{Safeties}_i+\beta_6 \text{Horsprogramme}_i+\gamma_g+\varepsilon_{ig}\). Here \(W_i^{>\mathrm{med}}=\mathbf{1}\{W_{\text{main},i}>\widetilde{W}\}\), where \(\widetilde{W}\) is the pooled median of non-missing \(W_{\text{main}}\) in the regression-eligible sample. The control set follows Table D7 (algorithm-selection, overlap, random-selection, replacement, and ad hoc indicators), with the addition of the discrepancy term \(W_i^{>\mathrm{med}}\). Panel A includes surveyed firms that report having been audited recently; Panel B includes all surveyed firms. Within each panel, columns report full audits, desk audits, and all audits. Corruption in general is the declared belief about the percentage of audits in Senegal affected by corruption; Corruption experience indicates that the respondent declared an instance of bribery to obtain tax favors; Grade on honesty is a 0-to-10 grade for the honesty of the inspector with whom the respondent last interacted during a full audit; and Friend at tax authority captures agreement with the statement that if the firm's owner has a friend at the tax authority, the firm will rarely be audited. Fixed effects are at the inspector-cluster-year level. Robust standard errors (Huber-White) are in parentheses. Stars denote significance at \(^{*}p<0.10\), \(^{**}p<0.05\), \(^{***}p<0.01\).
\end{tablenotes}
\end{threeparttable}
\end{adjustbox}
\end{table}
\end{landscape}

\begin{table}[!htbp]
\centering
\begin{adjustbox}{max width=1\textwidth}
\begin{threeparttable}
\caption{Robustness: \(W_{\ge q75}\) and survey indexes (Recent audit / All surveyed firms)}
\begin{tabular}{l*{6}{c}}
\textbf{A: Surveyed Firms (Self-Reporting a Recent Audit)}
                    &\multicolumn{3}{c}{Efficiency Index}                             &\multicolumn{3}{c}{Corruption Index}                             \\
                    &\multicolumn{1}{c}{(1)}&\multicolumn{1}{c}{(2)}&\multicolumn{1}{c}{(3)}&\multicolumn{1}{c}{(4)}&\multicolumn{1}{c}{(5)}&\multicolumn{1}{c}{(6)}\\
                    &\multicolumn{1}{c}{Full audits}&\multicolumn{1}{c}{Desk audits}&\multicolumn{1}{c}{All audits}&\multicolumn{1}{c}{Full audits}&\multicolumn{1}{c}{Desk audits}&\multicolumn{1}{c}{All audits}\\
\hline
Downward revision (\% of notification) in top quartile&        0.24         &       -0.32         &       -0.19         &        0.04         &        0.29         &        0.21         \\
                    &      (0.34)         &      (0.29)         &      (0.21)         &      (0.31)         &      (0.27)         &      (0.19)         \\
[1em]
Algorithm           &       -0.57         &        0.38         &       -0.13         &        0.29         &       -0.37         &       -0.04         \\
                    &      (0.37)         &      (0.32)         &      (0.24)         &      (0.30)         &      (0.29)         &      (0.21)         \\
[1em]
Inspectors x Overlap&        0.06         &       -1.02         &       -0.65         &        0.62         &       -0.49         &        0.19         \\
                    &      (0.43)         &      (0.67)         &      (0.61)         &      (0.41)         &      (0.38)         &      (0.25)         \\
[1em]
Algorithm x Random  &                     &        0.05         &        0.08         &                     &       -0.16         &       -0.12         \\
                    &                     &      (0.33)         &      (0.29)         &                     &      (0.33)         &      (0.25)         \\
\hline
N                   &          55         &          81         &         138         &          54         &          81         &         137         \\
R2                  &        0.13         &        0.27         &        0.09         &        0.27         &        0.34         &        0.19         \\
Mean outcome        &       -0.01         &       -0.01         &        0.00         &       -0.13         &       -0.10         &       -0.11         \\
\hline
\textbf{B: All Surveyed Firms}
                    &\multicolumn{1}{c}{(1)}         &\multicolumn{1}{c}{(2)}         &\multicolumn{1}{c}{(3)}         &\multicolumn{1}{c}{(4)}         &\multicolumn{1}{c}{(5)}         &\multicolumn{1}{c}{(6)}         \\
\hline
Downward revision (\% of notification) in top quartile&        0.30         &       -0.25         &       -0.10         &        0.10         &        0.47\sym{*}  &        0.31         \\
                    &      (0.31)         &      (0.24)         &      (0.18)         &      (0.29)         &      (0.25)         &      (0.19)         \\
[1em]
Algorithm           &       -0.60\sym{*}  &        0.16         &       -0.21         &        0.35         &       -0.37         &       -0.06         \\
                    &      (0.35)         &      (0.28)         &      (0.22)         &      (0.29)         &      (0.28)         &      (0.21)         \\
[1em]
Inspectors x Overlap&        0.08         &       -1.34\sym{**} &       -0.74         &        0.51         &       -0.39         &        0.26         \\
                    &      (0.39)         &      (0.63)         &      (0.59)         &      (0.39)         &      (0.37)         &      (0.24)         \\
[1em]
Algorithm x Random  &                     &        0.06         &        0.07         &                     &       -0.21         &       -0.15         \\
                    &                     &      (0.30)         &      (0.26)         &                     &      (0.33)         &      (0.25)         \\
\hline
N                   &          61         &          96         &         159         &          60         &          97         &         159         \\
R2                  &        0.16         &        0.21         &        0.08         &        0.30         &        0.24         &        0.13         \\
Mean outcome        &        0.01         &        0.03         &        0.03         &       -0.22         &       -0.05         &       -0.11         \\
\hline

\end{tabular}
\begin{tablenotes}[flushleft]
\footnotesize
\item Notes: OLS results for the Efficiency Index and Corruption Index. The specification is \(y_{ig}=\beta_0+\beta_1 W_i^{\ge q75}+\beta_2 \text{Algorithm}_i+\beta_3 \text{Overlap}_i+\beta_4 \text{Random}_i+\beta_5 \text{Safeties}_i+\beta_6 \text{Horsprogramme}_i+\gamma_g+\varepsilon_{ig}\). Here \(W_i^{\ge q75}=\mathbf{1}\{W_{\text{main},i}\ge\widetilde{W}_{75}\}\), where \(\widetilde{W}_{75}\) is the pooled 75th percentile of non-missing \(W_{\text{main}}\) in the regression-eligible sample. The control set follows Table D7 (algorithm-selection, overlap, random-selection, replacement, and ad hoc indicators), with the addition of the discrepancy term \(W_i^{\ge q75}\). Panel A includes surveyed firms that report having been audited recently; Panel B includes all surveyed firms. Within each panel, columns report full audits, desk audits, and all audits. Fixed effects are at the selection-year \(\times\) tax-office level. Robust standard errors (Huber-White) are in parentheses. Stars denote significance at \(^{*}p<0.10\), \(^{**}p<0.05\), \(^{***}p<0.01\).
\end{tablenotes}
\end{threeparttable}
\end{adjustbox}
\end{table}

\begin{landscape}
\begin{table}[!htbp]
\centering
\begin{adjustbox}{max width=1.25\textwidth}
\begin{threeparttable}
\caption{Robustness: \(W_{\ge q75}\) and survey questions (Recent audit / All surveyed firms)}
\begin{tabular}{l*{12}{c}}
\textbf{A: Surveyed Firms (Self-Reporting a Recent Audit)}
                    &\multicolumn{3}{c}{Corruption in General}                        &\multicolumn{3}{c}{Corruption Experience}                        &\multicolumn{3}{c}{Grade on Honesty}                             &\multicolumn{3}{c}{Friend at tax Auth.}                          \\
                    &\multicolumn{1}{c}{(1)}&\multicolumn{1}{c}{(2)}&\multicolumn{1}{c}{(3)}&\multicolumn{1}{c}{(4)}&\multicolumn{1}{c}{(5)}&\multicolumn{1}{c}{(6)}&\multicolumn{1}{c}{(7)}&\multicolumn{1}{c}{(8)}&\multicolumn{1}{c}{(9)}&\multicolumn{1}{c}{(10)}&\multicolumn{1}{c}{(11)}&\multicolumn{1}{c}{(12)}\\
                    &\multicolumn{1}{c}{Full audits}&\multicolumn{1}{c}{Desk audits}&\multicolumn{1}{c}{All audits}&\multicolumn{1}{c}{Full audits}&\multicolumn{1}{c}{Desk audits}&\multicolumn{1}{c}{All audits}&\multicolumn{1}{c}{Full audits}&\multicolumn{1}{c}{Desk audits}&\multicolumn{1}{c}{All audits}&\multicolumn{1}{c}{Full audits}&\multicolumn{1}{c}{Desk audits}&\multicolumn{1}{c}{All audits}\\
\hline
Downward revision (\% of notification) in top quartile&      11.780         &     -13.054         &       3.656         &      -0.095\sym{*}  &      -0.034         &      -0.075\sym{*}  &      -0.619         &      -0.620         &      -0.531         &      -0.205         &       0.535         &       0.080         \\
                    &     (8.289)         &    (15.636)         &     (7.835)         &     (0.052)         &     (0.059)         &     (0.038)         &     (0.631)         &     (0.612)         &     (0.503)         &     (0.458)         &     (0.408)         &     (0.345)         \\
[1em]
Algorithm           &      -6.734         &     -23.874         &      -8.631         &       0.040         &      -0.036         &       0.019         &      -0.779         &      -0.949         &      -0.801         &       0.203         &       0.544         &       0.248         \\
                    &     (6.332)         &    (21.692)         &     (7.504)         &     (0.089)         &     (0.070)         &     (0.076)         &     (0.875)         &     (1.071)         &     (0.765)         &     (0.501)         &     (0.944)         &     (0.479)         \\
[1em]
Inspectors x Overlap&      29.902\sym{***}&                     &      34.686\sym{***}&       0.084         &                     &       0.064         &      -0.481         &                     &      -0.552         &       0.273         &                     &       0.002         \\
                    &     (8.100)         &                     &     (6.991)         &     (0.057)         &                     &     (0.043)         &     (1.527)         &                     &     (1.627)         &     (0.695)         &                     &     (0.646)         \\
[1em]
Algorithm x Random  &                     &                     &                     &                     &      -0.264         &      -0.255         &                     &       0.461         &      -1.083         &                     &      -1.605         &      -0.670         \\
                    &                     &                     &                     &                     &     (0.209)         &     (0.181)         &                     &     (1.807)         &     (0.920)         &                     &     (1.054)         &     (0.631)         \\
\hline
N                   &          49         &          45         &          94         &          65         &          89         &         154         &          61         &          67         &         128         &          57         &          76         &         133         \\
R2                  &       0.283         &       0.444         &       0.296         &       0.154         &       0.369         &       0.248         &       0.147         &       0.546         &       0.331         &       0.080         &       0.399         &       0.229         \\
Mean outcome        &       12.63         &       17.33         &       14.88         &        0.07         &        0.04         &        0.05         &        6.68         &        6.44         &        6.56         &        3.45         &        3.35         &        3.39         \\
\hline
\textbf{B: All Surveyed Firms}
                    &\multicolumn{1}{c}{(1)}         &\multicolumn{1}{c}{(2)}         &\multicolumn{1}{c}{(3)}         &\multicolumn{1}{c}{(4)}         &\multicolumn{1}{c}{(5)}         &\multicolumn{1}{c}{(6)}         &\multicolumn{1}{c}{(7)}         &\multicolumn{1}{c}{(8)}         &\multicolumn{1}{c}{(9)}         &\multicolumn{1}{c}{(10)}         &\multicolumn{1}{c}{(11)}         &\multicolumn{1}{c}{(12)}         \\
\hline
Downward revision (\% of notification) in top quartile&      11.482         &     -16.220         &       2.365         &      -0.104\sym{**} &      -0.085\sym{*}  &      -0.096\sym{***}&      -0.671         &      -0.275         &      -0.445         &      -0.148         &       0.289         &       0.056         \\
                    &     (8.507)         &    (13.043)         &     (7.493)         &     (0.049)         &     (0.050)         &     (0.035)         &     (0.604)         &     (0.712)         &     (0.487)         &     (0.456)         &     (0.459)         &     (0.329)         \\
[1em]
Algorithm           &      -8.040         &     -16.979         &      -9.026         &       0.041         &       0.032         &       0.038         &      -0.730         &      -0.264         &      -0.619         &       0.252         &       0.415         &       0.301         \\
                    &     (6.046)         &    (14.546)         &     (6.419)         &     (0.090)         &     (0.068)         &     (0.065)         &     (0.860)         &     (0.782)         &     (0.685)         &     (0.493)         &     (0.720)         &     (0.421)         \\
[1em]
Inspectors x Overlap&      14.378         &                     &      20.969         &      -0.073         &                     &      -0.080         &      -0.009         &                     &      -0.168         &       0.211         &                     &       0.020         \\
                    &    (17.544)         &                     &    (16.704)         &     (0.163)         &                     &     (0.180)         &     (1.290)         &                     &     (1.359)         &     (0.596)         &                     &     (0.506)         \\
[1em]
Algorithm x Random  &                     &                     &                     &                     &      -0.303\sym{*}  &      -0.285\sym{*}  &                     &      -1.387         &      -1.650\sym{**} &                     &      -0.926         &      -0.600         \\
                    &                     &                     &                     &                     &     (0.172)         &     (0.158)         &                     &     (0.892)         &     (0.694)         &                     &     (0.930)         &     (0.570)         \\
\hline
N                   &          54         &          54         &         108         &          73         &         111         &         184         &          67         &          78         &         145         &          61         &          88         &         149         \\
R2                  &       0.210         &       0.438         &       0.281         &       0.115         &       0.418         &       0.270         &       0.156         &       0.529         &       0.341         &       0.080         &       0.327         &       0.217         \\
Mean outcome        &       12.85         &       15.64         &       14.25         &        0.09         &        0.06         &        0.07         &        6.74         &        6.50         &        6.61         &        3.39         &        3.40         &        3.40         \\
\hline

\end{tabular}
\begin{tablenotes}[flushleft]
\footnotesize
\item Notes: OLS results for four taxpayer-survey outcomes: Corruption in general (Columns 1 to 3), Corruption experience (Columns 4 to 6), Grade on honesty (Columns 7 to 9), and Friend at tax authority (Columns 10 to 12). The specification is \(y_{ig}=\beta_0+\beta_1 W_i^{\ge q75}+\beta_2 \text{Algorithm}_i+\beta_3 \text{Overlap}_i+\beta_4 \text{Random}_i+\beta_5 \text{Safeties}_i+\beta_6 \text{Horsprogramme}_i+\gamma_g+\varepsilon_{ig}\). Here \(W_i^{\ge q75}=\mathbf{1}\{W_{\text{main},i}\ge\widetilde{W}_{75}\}\), where \(\widetilde{W}_{75}\) is the pooled 75th percentile of non-missing \(W_{\text{main}}\) in the regression-eligible sample. The control set follows Table D7 (algorithm-selection, overlap, random-selection, replacement, and ad hoc indicators), with the addition of the discrepancy term \(W_i^{\ge q75}\). Panel A includes surveyed firms that report having been audited recently; Panel B includes all surveyed firms. Within each panel, columns report full audits, desk audits, and all audits. Corruption in general is the declared belief about the percentage of audits in Senegal affected by corruption; Corruption experience indicates that the respondent declared an instance of bribery to obtain tax favors; Grade on honesty is a 0-to-10 grade for the honesty of the inspector with whom the respondent last interacted during a full audit; and Friend at tax authority captures agreement with the statement that if the firm's owner has a friend at the tax authority, the firm will rarely be audited. Fixed effects are at the inspector-cluster-year level. Robust standard errors (Huber-White) are in parentheses. Stars denote significance at \(^{*}p<0.10\), \(^{**}p<0.05\), \(^{***}p<0.01\).
\end{tablenotes}
\end{threeparttable}
\end{adjustbox}
\end{table}
\end{landscape}




\end{document}










